% Leçon 13 — Vocabulaire et compréhension de texte : maladies en général et du SIDA en particulier — Chinois 1ère A — Cours signé OnBuch+
% Programme officiel MINESEC. Compiler avec : tectonic ZHA13-vocabulaire-maladies-en-general-et-du-sida-en-part.tex
\newcommand{\DOCMATIERE}{Chinois}
\newcommand{\DOCNIVEAU}{1ère A}
\newcommand{\DOCMODULE}{Module 2 : Environnement et santé}
\newcommand{\DOCLECON}{ZHA13}
\newcommand{\DOCTITRE}{Leçon 13 — Vocabulaire et compréhension de texte : maladies en général et du SIDA en particulier}
% OnBuch+ — préambule « cours signés » ENRICHI (compilé avec tectonic / XeLaTeX)
% Système visuel « Pop » d'OnBuch+ : fond crème (jamais blanc), bordures encre
% épaisses, ombres « dures » décalées, titres Archivo Black en capitales.
% Version MATHÉMATIQUES : graphes (pgfplots/tikz), boîtes pédagogiques
% (définition, propriété, méthode, attention, le savais-tu, exercice résolu,
% exercice, TP, bilan), page de garde et signature OnBuch+ ancrée en bas.
%
% Macros attendues dans main.tex AVANT \input{preamble} :
%   \DOCMATIERE  \DOCNIVEAU  \DOCMODULE  \DOCLECON (numéro)  \DOCTITRE
\documentclass[11pt]{article}

\usepackage{fontspec}
\usepackage[french]{babel} % français = langue principale ; le chinois via \zh{..} / textebox (xeCJK)
\usepackage{xeCJK}
\usepackage{pifont} % \ding{51} \ding{55} utilisés par les modèles
\usepackage[a4paper,margin=2cm,top=3.4cm,bottom=2.8cm,headheight=1.4cm,footskip=1.3cm]{geometry}
\usepackage{amsmath,amssymb,amsfonts}
\usepackage{mathtools} % \xmapsto, \coloneqq... souvent utilisés par les modèles
\usepackage{cancel} % simplifications barrées (bilans de forces)
% \abs{z} (module, valeur absolue) et \norm{u} (norme), utilisés par les modèles
\providecommand{\abs}{}\let\abs\relax\DeclarePairedDelimiter{\abs}{\lvert}{\rvert}
\providecommand{\norm}{}\let\norm\relax\DeclarePairedDelimiter{\norm}{\lVert}{\rVert}
\usepackage[table]{xcolor} % [table] : fournit aussi \rowcolors (alternance de lignes)
\usepackage{pagecolor}
\usepackage{tikz}
\usetikzlibrary{arrows.meta,positioning,calc,shapes.geometric,shapes.misc,shapes.arrows,decorations.pathmorphing,decorations.pathreplacing,decorations.markings,patterns,shadows,mindmap,trees,fit,backgrounds,matrix,intersections,angles,quotes}
\usepackage{pgfplots}
\usepackage[version=4]{mhchem} % équations nucléaires \ce{^{235}_{92}U}
\usepackage[european,siunitx]{circuitikz} % schémas électriques
\pgfplotsset{compat=1.18}
\usepgfplotslibrary{fillbetween,groupplots} % aires entre courbes (calcul intégral)
\usepackage{enumitem}
\usepackage{fancyhdr}
% [most] charge toutes les bibliothèques tcolorbox (listings, minted, raster,
% documentation...) et gonfle inutilement la mémoire TeX sur un document long
% (~300 boîtes) : on ne charge que ce qui est réellement utilisé.
\usepackage[skins,breakable]{tcolorbox}
\usepackage{graphicx}
\usepackage{colortbl}
\usepackage{booktabs}
\usepackage{tabularx}
\usepackage{diagbox} % case d'en-tête diagonale des tableaux à double entrée
% Colonnes extensibles centrée / gauche / droite (C, L, R), souvent utilisées par les modèles
\newcolumntype{C}{>{\centering\arraybackslash}X}
\newcolumntype{L}{>{\raggedright\arraybackslash}X}
\newcolumntype{R}{>{\raggedleft\arraybackslash}X}
\usepackage{array}
\usepackage{multicol}
\usepackage{multirow}
\usepackage{siunitx}
% parse-numbers=false : accepte aussi \SI{2,5 \times 10^{-3}}{...}
\sisetup{locale=FR,output-decimal-marker={,},group-minimum-digits=4,parse-numbers=false}
\DeclareSIUnit\ohm{\ensuremath{\symup{\Omega}}} % U+2126 absent de la police
\DeclareSIUnit\division{div} % réglages d'oscilloscope (ms/div, V/div)
\DeclareSIUnit\tr{tr} % tours (vitesse de rotation, ex. \SI{1500}{\tr\per\minute})
\DeclareSIUnit\molar{M} % concentration molaire (ex. \SI{3}{\molar}, \SI{1}{\micro\molar})
\DeclareSIUnit\an{an} % année (le modèle confond parfois avec \year, un compteur TeX) — ex. \SI{5}{\kilo\meter\per\an}
\DeclareSIUnit\FCFA{FCFA} % franc CFA (ex. \SI{500}{\FCFA\per\kilo\gram})
\DeclareSIUnit\degreeBrix{\textdegree Brix} % teneur en sucre d'un sirop/jus (réfractométrie)
\DeclareSIUnit\month{mois} % le modèle utilise parfois \month, un compteur TeX, comme unité de durée
\usepackage{qrcode}
% \degree : commande fréquemment utilisée par les modèles pour un angle ou une
% température en dehors de \SI{}{} (non fournie par les paquets chargés).
\providecommand{\degree}{^{\circ}}
% \qed (fin de démonstration), utilisé par les modèles sans amsthm
\providecommand{\qed}{\hfill$\square$}
% \barre : double barre des valeurs interdites dans un tableau de variations (tkz-tab)
\providecommand{\barre}{\ensuremath{\|}}
% environnement proof (amsthm non chargé) utilisé par les modèles
\ifdefined\proof\else\newenvironment{proof}[1][Démonstration]{\par\noindent\textit{#1.}\ }{\qed\par}\fi
\usepackage{lastpage}
\usepackage{hyperref}
\hypersetup{hidelinks}

% ============================================================
% Palette « Pop » OnBuch+ — mêmes hex que la classe P (plus_ui.dart)
% ============================================================
\definecolor{poporange}{HTML}{F59321}
\definecolor{popdark}{HTML}{E07A0C}
\definecolor{popcream}{HTML}{FFF6E8}
\definecolor{popink}{HTML}{1C1714}
\definecolor{poppurple}{HTML}{7A5AE0}
\definecolor{popgreen}{HTML}{2E9E6A}
\definecolor{poppink}{HTML}{E0457B}
\definecolor{popblue}{HTML}{2F7DE1}
\definecolor{popgold}{HTML}{C98A12}
\definecolor{popmuted}{HTML}{8A7F76}
\definecolor{popline}{HTML}{EADFD2}
\definecolor{popgray}{HTML}{8A7F76}
\colorlet{popgrey}{popgray}
% Alias tolérants (le modèle nomme parfois les couleurs différemment)
\colorlet{popwhite}{white}
\colorlet{popblack}{popink}
\colorlet{popred}{poppink}
\colorlet{popyellow}{popgold}
% Teintes claires (fonds de tableaux/encadrés) — le modèle nomme parfois
% les couleurs avec un suffixe L (ex. popblueL) sans les avoir définies.
\colorlet{poporangeL}{poporange!20}
\colorlet{popdarkL}{popdark!20}
\colorlet{poppurpleL}{poppurple!20}
\colorlet{popgreenL}{popgreen!20}
\colorlet{poppinkL}{poppink!20}
\colorlet{popblueL}{popblue!20}
\colorlet{popgoldL}{popgold!20}
\colorlet{popredL}{popred!20}
\colorlet{popgrayL}{popgray!20}
% Teintes claires pour fonds de tableaux / zones
\colorlet{popblueL}{popblue!10!white}
\colorlet{poporangeL}{poporange!14!white}
\colorlet{popgreenL}{popgreen!12!white}
\colorlet{poppurpleL}{poppurple!12!white}
\colorlet{poppinkL}{poppink!12!white}
\colorlet{popgoldL}{popgold!14!white}

\pagecolor{popcream}

% ============================================================
% Typographie
% ============================================================
\defaultfontfeatures{Ligatures=TeX}
\setmainfont{PlusJakartaSans-Medium.ttf}[Path=../../../fonts/,BoldFont=PlusJakartaSans-Bold.ttf]
% Symboles, API et emojis absents de la police principale : repli sur DejaVu Sans (caractères actifs)
\newfontface\symfont{DejaVuSans.ttf}[Path=../../../fonts/]
\makeatletter
\newcommand\symactive[1]{\catcode#1=13 \begingroup\lccode`\~=#1 \lowercase{\endgroup\def~}{{\symfont\char#1}}}
\newcommand\symblank[1]{\catcode#1=13 \begingroup\lccode`\~=#1 \lowercase{\endgroup\def~}{}}
\newcommand\symhyph[1]{\catcode#1=13 \begingroup\lccode`\~=#1 \lowercase{\endgroup\def~}{\mbox{-}}}
\newcount\symc
\newcommand\symrange[2]{\symc=#1 \@whilenum\symc<#2 \do{\expandafter\symactive\expandafter{\the\symc}\advance\symc by 1 }}
\newcommand\symblankrange[2]{\symc=#1 \@whilenum\symc<#2 \do{\expandafter\symblank\expandafter{\the\symc}\advance\symc by 1 }}
\makeatother
\symrange{"0250}{"02FF}\symactive{"2713}\symactive{"2714}\symactive{"2717}\symactive{"2718}\symactive{"2610}\symactive{"2611}\symactive{"2612}
\symrange{"2600}{"27BF}
\catcode"2705=13 \begingroup\lccode`\~="2705 \lowercase{\endgroup\def~}{{\symfont\char"2713}}\symblank{"26BD}\symblank{"26A1}
\symrange{"2460}{"257F}\symactive{"223C}\symactive{"2248}\symactive{"2260}
\symhyph{"2011}\symhyph{"2E1A}\symblankrange{"1F300}{"1FAFF}\symblank{"FE0F}
% Caractères chinois : WenQuanYi Zen Hei (sous-ensemble GB2312 + pinyin), gras/italique simulés
\setCJKmainfont{WenQuanYiZenHei-subset.ttf}[Path=../../../fonts/,AutoFakeBold=3,AutoFakeSlant=0.2]
\xeCJKsetup{CJKspace=false}
% Pinyin : voyelles à caron (ǎ ǐ ǒ ǔ ǖ ǘ ǚ ǜ) absentes de la police principale -> caractères actifs
\newfontface\pyfont{WenQuanYiZenHei-subset.ttf}[Path=../../../fonts/]
\newcommand\pyactive[1]{\catcode#1=13 \begingroup\lccode`\~=#1 \lowercase{\endgroup\def~}{{\pyfont\char#1}}}
\pyactive{"01CD}\pyactive{"01CE}\pyactive{"01CF}\pyactive{"01D0}\pyactive{"01D1}\pyactive{"01D2}\pyactive{"01D3}\pyactive{"01D4}\pyactive{"01D5}\pyactive{"01D6}\pyactive{"01D7}\pyactive{"01D8}\pyactive{"01D9}\pyactive{"01DA}\pyactive{"01DB}\pyactive{"01DC}
% Maths en sans-serif (Fira Math) avec chiffres et lettres droites de Plus
% Jakarta, pour que les formules se fondent dans le texte.
\usepackage[math-style=ISO,bold-style=ISO]{unicode-math}
\setmathfont{FiraMath-Regular.otf}[Path=../../../fonts/]
\setmathfont{PlusJakartaSans-Medium.ttf}[Path=../../../fonts/,range={up/num,up/{Latin,latin},\mathup}]
\usepackage{icomma} % virgule décimale sans espace en mode math
% Caractères mathématiques Unicode absents des polices de texte (Plus Jakarta,
% Archivo Black) : rendus par leur équivalent mathématique, en texte comme en
% formule — sinon ils s'affichent en carré vide (ex. « Arithmétique dans ℤ »).
\usepackage{newunicodechar}
\newunicodechar{ℝ}{\ensuremath{\mathbb{R}}}
\newunicodechar{ℤ}{\ensuremath{\mathbb{Z}}}
\newunicodechar{ℂ}{\ensuremath{\mathbb{C}}}
\newunicodechar{ℕ}{\ensuremath{\mathbb{N}}}
\newunicodechar{ℚ}{\ensuremath{\mathbb{Q}}}
\newunicodechar{𝔻}{\ensuremath{\mathbb{D}}}
\newunicodechar{α}{\ensuremath{\alpha}}
\newunicodechar{β}{\ensuremath{\beta}}
\newunicodechar{γ}{\ensuremath{\gamma}}
\newunicodechar{δ}{\ensuremath{\delta}}
\newunicodechar{ε}{\ensuremath{\varepsilon}}
\newunicodechar{θ}{\ensuremath{\theta}}
\newunicodechar{λ}{\ensuremath{\lambda}}
\newunicodechar{μ}{\ensuremath{\mu}}
\newunicodechar{σ}{\ensuremath{\sigma}}
\newunicodechar{τ}{\ensuremath{\tau}}
\newunicodechar{φ}{\ensuremath{\varphi}}
\newunicodechar{ψ}{\ensuremath{\psi}}
\newunicodechar{ω}{\ensuremath{\omega}}
\newunicodechar{Σ}{\ensuremath{\Sigma}}
% majuscules produites par \MakeUppercase dans les titres (α-aminés -> Α)
\newunicodechar{Α}{\ensuremath{\alpha}}
\newunicodechar{Β}{\ensuremath{\beta}}
\newunicodechar{Γ}{\ensuremath{\Gamma}}
\newunicodechar{Δ}{\ensuremath{\Delta}}
\newunicodechar{Ε}{\ensuremath{\varepsilon}}
\newunicodechar{Θ}{\ensuremath{\Theta}}
\newunicodechar{Λ}{\ensuremath{\Lambda}}
\newunicodechar{Μ}{\ensuremath{\mu}}
\newunicodechar{Τ}{\ensuremath{\tau}}
\newunicodechar{Φ}{\ensuremath{\Phi}}
\newunicodechar{Ψ}{\ensuremath{\Psi}}
\newunicodechar{Ω}{\ensuremath{\Omega}}
\newunicodechar{↦}{\ensuremath{\mapsto}}
\newunicodechar{∈}{\ensuremath{\in}}
\newunicodechar{∉}{\ensuremath{\notin}}
\newunicodechar{∧}{\ensuremath{\wedge}}
\newunicodechar{⇔}{\ensuremath{\Leftrightarrow}}
\newunicodechar{⟺}{\ensuremath{\Longleftrightarrow}}
\newunicodechar{⇒}{\ensuremath{\Rightarrow}}
\newunicodechar{⟹}{\ensuremath{\Longrightarrow}}
\newunicodechar{∘}{\ensuremath{\circ}}
\newunicodechar{∪}{\ensuremath{\cup}}
\newunicodechar{∩}{\ensuremath{\cap}}
\newunicodechar{≡}{\ensuremath{\equiv}}
\newunicodechar{⊂}{\ensuremath{\subset}}
\newunicodechar{⊥}{\ensuremath{\perp}}
\newunicodechar{∃}{\ensuremath{\exists}}
\newunicodechar{∀}{\ensuremath{\forall}}
\newunicodechar{∛}{\ensuremath{\sqrt[3]{\,}}}
\newunicodechar{∜}{\ensuremath{\sqrt[4]{\,}}}
\newunicodechar{⃗}{\ensuremath{{}^{\to}}}
\newunicodechar{ⁿ}{\ifmmode^{n}\else\textsuperscript{\ensuremath{n}}\fi}
\newunicodechar{ˣ}{\ifmmode^{x}\else\textsuperscript{\ensuremath{x}}\fi}
\newunicodechar{ᵇ}{\ifmmode^{b}\else\textsuperscript{\ensuremath{b}}\fi}
\newunicodechar{ʳ}{\ifmmode^{r}\else\textsuperscript{\ensuremath{r}}\fi}
\newunicodechar{ᵘ}{\ifmmode^{u}\else\textsuperscript{\ensuremath{u}}\fi}
\newunicodechar{ʸ}{\ifmmode^{y}\else\textsuperscript{\ensuremath{y}}\fi}
\newunicodechar{ᵃ}{\ifmmode^{a}\else\textsuperscript{\ensuremath{a}}\fi}
\newunicodechar{ᵉ}{\ifmmode^{e}\else\textsuperscript{\ensuremath{e}}\fi}
\newunicodechar{ᵏ}{\ifmmode^{k}\else\textsuperscript{\ensuremath{k}}\fi}
\newunicodechar{ᵖ}{\ifmmode^{p}\else\textsuperscript{\ensuremath{p}}\fi}
\newunicodechar{ᴺ}{\ifmmode^{N}\else\textsuperscript{\ensuremath{N}}\fi}
\newunicodechar{ᶜ}{\ifmmode^{c}\else\textsuperscript{\ensuremath{c}}\fi}
\newunicodechar{ʰ}{\ifmmode^{h}\else\textsuperscript{\ensuremath{h}}\fi}
\newunicodechar{ᵠ}{\ifmmode^{q}\else\textsuperscript{\ensuremath{q}}\fi}
\newunicodechar{ₙ}{\ifmmode_{n}\else\textsubscript{\ensuremath{n}}\fi}
\newunicodechar{ₐ}{\ifmmode_{a}\else\textsubscript{\ensuremath{a}}\fi}
\newunicodechar{ᵢ}{\ifmmode_{i}\else\textsubscript{\ensuremath{i}}\fi}
\newunicodechar{ᵣ}{\ifmmode_{r}\else\textsubscript{\ensuremath{r}}\fi}
\newunicodechar{ₖ}{\ifmmode_{k}\else\textsubscript{\ensuremath{k}}\fi}
\newunicodechar{ᵦ}{\ifmmode_{\beta}\else\textsubscript{\ensuremath{\beta}}\fi}
\newunicodechar{ₓ}{\ifmmode_{x}\else\textsubscript{\ensuremath{x}}\fi}
\newfontfamily\popdisplay{ArchivoBlack-Regular.ttf}[Path=../../../fonts/]
\newfontfamily\poplogo{SpaceGrotesk-Bold.ttf}[Path=../../../fonts/]
\newfontfamily\popsemi{PlusJakartaSans-SemiBold.ttf}[Path=../../../fonts/]
\newfontfamily\popxbold{PlusJakartaSans-ExtraBold.ttf}[Path=../../../fonts/]
\newfontfamily\popxboldls{PlusJakartaSans-ExtraBold.ttf}[Path=../../../fonts/,LetterSpace=3.0]

\setlist[enumerate]{leftmargin=1.7em,itemsep=5pt,topsep=4pt}
\setlist[itemize]{leftmargin=1.7em,itemsep=4pt,topsep=4pt,label={\color{poporange}\textbullet}}
\setlength{\parindent}{0pt}
\setlength{\parskip}{5pt}
\renewcommand{\arraystretch}{1.25}

% pgfplots : style Pop par défaut
\pgfplotsset{
  every axis/.append style={axis line style={popink,line width=1pt},tick style={popink},
    grid style={popline,line width=0.5pt},label style={font=\small},tick label style={font=\footnotesize}},
  popcurve/.style={poporange,line width=1.8pt,no markers,smooth},
}
% Même style disponible dans un \draw TikZ simple (hors pgfplots)
\tikzset{popcurve/.style={poporange,line width=1.8pt}}
\tikzset{popgrid/.style={popline,very thin}}
\colorlet{popcurve}{poporange} % le nom du style est parfois utilisé comme couleur

% ============================================================
% Pill / squiggle
% ============================================================
\newcommand{\poppill}[3]{%
  \tikz[baseline=-0.55ex]\node[draw=popink,line width=1.2pt,rounded corners=2.6pt,%
    fill=#2,inner xsep=6.5pt,inner ysep=3pt]{%
    \popxboldls\fontsize{7}{7}\selectfont\color{#3}\MakeUppercase{#1}};%
}
\newcommand{\popsquiggle}[1]{%
  \tikz[baseline=-0.4ex]{
    \draw[#1,line width=2.6pt,line cap=round]
      plot [smooth] coordinates {(0,0.09) (0.34,0.01) (0.68,0.21) (1.02,0.01) (1.36,0.10)};
  }%
}
% Mot-clé surligné (marqueur orange sous le texte)
\newcommand{\cle}[1]{{\popxbold\color{popdark}#1}}

% ============================================================
% En-tête / pied
% ============================================================
\pagestyle{fancy}
\fancyhf{}
\renewcommand{\headrulewidth}{0pt}
\renewcommand{\footrulewidth}{0pt}
\lhead{\raisebox{-0.15ex}{\poplogo\fontsize{13}{13}\selectfont\color{popink}OnBuch\color{poporange}+}}
\rhead{\poppill{\DOCMATIERE\ \textbullet\ \DOCNIVEAU\ \textbullet\ Leçon \DOCLECON}{white}{popink}}
\lfoot{\popxbold\fontsize{7.6}{7.6}\selectfont\color{popmuted}\MakeUppercase{Cours signé OnBuch+}\ \textbullet\ {\popsemi\DOCTITRE}}
\rfoot{\tikz[baseline=-0.6ex]\node[rounded corners=8.5pt,fill=popink,minimum height=17pt,minimum width=17pt,inner xsep=5pt,inner sep=0]{\popxbold\fontsize{8}{8}\selectfont\color{popcream}\thepage\,/\,\pageref*{LastPage}};}

% ============================================================
% Titres
% ============================================================
\newcommand{\chaptitre}[2]{%
  \begin{tcolorbox}[enhanced,colback=poporange,colframe=popink,boxrule=2.6pt,arc=11pt,
      drop shadow=popink,left=15pt,right=15pt,top=12pt,bottom=11pt,before skip=3pt,after skip=18pt]
    \poppill{#2}{popink}{popcream}\par\vspace{9pt}
    {\raggedright\hyphenpenalty=10000\exhyphenpenalty=10000\popdisplay\fontsize{22}{24}\selectfont\color{popcream}\MakeUppercase{#1}\par}
  \end{tcolorbox}
}
% Section numérotée : pastille orange avec le numéro + titre Archivo
\newcounter{popsec}
\newcommand{\coursec}[1]{%
  \stepcounter{popsec}\setcounter{popsub}{0}%
  \par\vspace{16pt}\noindent
  \begin{minipage}{\linewidth}
  \tikz[baseline=-0.7ex]\node[draw=popink,line width=1.6pt,fill=poporange,rounded corners=4pt,minimum size=22pt,inner sep=0]{\popdisplay\fontsize{12}{12}\selectfont\color{popcream}\thepopsec};%
  \hspace{8pt}{\popdisplay\fontsize{15}{17}\selectfont\color{popink}\MakeUppercase{#1}}\par
  \vspace{4pt}\noindent{\color{poporange}\rule{36pt}{3.6pt}}
  \end{minipage}\par\nopagebreak\vspace{8pt}
}
\newcounter{popsub}
\newcommand{\courssub}[1]{%
  \stepcounter{popsub}%
  \par\vspace{9pt}\noindent{\popxbold\fontsize{11.5}{13}\selectfont\color{popink}{\color{poporange}\thepopsec.\thepopsub}\hspace{6pt}#1}\par\nopagebreak\vspace{4pt}
}

% ============================================================
% Boîtes pédagogiques — même recette : fond blanc, bordure épaisse colorée,
% ombre dure encre, badge plein. Titre optionnel [#1].
% ============================================================
\tcbset{popbase/.style={breakable,enhanced,colback=white,boxrule=2.3pt,arc=9pt,
  shadow={3pt}{-3pt}{0pt}{popink},left=14pt,right=14pt,top=11pt,bottom=11pt,before skip=11pt,after skip=15pt,
  fontupper=\popsemi\fontsize{10.6}{14.5}\selectfont\color{popink},
  fontlower=\popsemi\fontsize{10.6}{14.5}\selectfont\color{popink}}}
% Coche dessinée (le glyphe ✓ n'existe pas dans les polices du document)
\newcommand{\popcheck}[1]{\tikz[baseline=-0.5ex]\draw[#1,line width=1.6pt,line cap=round,line join=round](-0.13,0.0)--(-0.04,-0.09)--(0.13,0.1);}
\providecommand{\coche}{\popcheck{popgreen}}
\providecommand{\resultat}[1]{\par\smallskip\noindent\fcolorbox{popgreen}{popgreenL}{\parbox{\dimexpr\linewidth-2\fboxsep-2\fboxrule\relax}{\textbf{Résultat.} #1}}\par\smallskip}
\newcommand{\popboxhead}[3]{\poppill{#1}{#2}{white}\ifx\relax#3\relax\else\hspace{6pt}{\popxbold\fontsize{10.6}{12}\selectfont\color{popink}#3}\fi\par\vspace{7pt}}

\newtcolorbox{aretenir}[1][]{popbase,colframe=popgreen,before upper={\popboxhead{À retenir}{popgreen}{#1}}}
% Texte en chinois (document de lecture, dialogue, extrait) : cadre
% d'étude. Un titre optionnel (source, auteur) s'affiche dans l'en-tête.
\newtcolorbox{textebox}[1][]{popbase,colframe=popblue,before upper={\popboxhead{文本 · Texte}{popblue}{#1}},after upper={}}
% Marques ✓ / ✗ (forme correcte / fautive) souvent utilisées par les modèles
\providecommand{\cross}{\textcolor{poppink}{\ensuremath{\times}}}
\providecommand{\xmark}{\cross}\providecommand{\crossmark}{\cross}\providecommand{\cmark}{\ensuremath{\checkmark}}
% Ligne de réponse / justification à compléter (inventée par les modèles)
\providecommand{\justification}[1]{\par\noindent\textit{Justification~:} \dotfill\par}
% \textipa (package tipa non chargé) : la police ne couvre pas l'API -> simple passage
\providecommand{\textipa}[1]{#1}
% Soulignement (ulem non chargé) : \uline, \ul
\providecommand{\uline}[1]{\underline{#1}}\providecommand{\ul}[1]{\underline{#1}}
% \glossaire{\item ...} inventée par les modèles : liste de glossaire
\providecommand{\glossaire}[1]{\begin{itemize}#1\end{itemize}}
% \alert (beamer) inventée par les modèles : texte mis en évidence
\providecommand{\alert}[1]{\textbf{\textcolor{poppink}{#1}}}
% variantes ASCII d'ordinaux écrites en macro
\providecommand{\iere}{\textsuperscript{re}}\providecommand{\ieme}{\textsuperscript{e}}\providecommand{\ere}{\textsuperscript{re}}\providecommand{\eme}{\textsuperscript{e}}\providecommand{\er}{\textsuperscript{er}}\providecommand{\ie}{\textsuperscript{e}}
% Ordinaux français écrits en macro par les modèles : 1\ière, 3\ième
\newcommand{\ière}{\textsuperscript{re}}\newcommand{\ième}{\textsuperscript{e}}
% \exemple (inventée par les modèles) : étiquette « Exemple : »
\providecommand{\exemple}{\textbf{Exemple~:}\ }
% \square utilisable aussi hors mode math (cases à cocher)
\AtBeginDocument{\let\squareMath\square\renewcommand{\square}{\ensuremath{\squareMath}}}
% Mot ou phrase en chinois à l'intérieur d'un paragraphe français
\newcommand{\zh}[1]{\ifmmode\text{#1}\else{#1}\fi}
\let\eng\zh \let\esp\zh \let\all\zh % alias tolérés
% flèches d'intonation inventées par les modèles
\providecommand{\textsearrow}{}\renewcommand{\textsearrow}{\ensuremath{\searrow}}\providecommand{\textnearrow}{}\renewcommand{\textnearrow}{\ensuremath{\nearrow}}
% \fr{..} : mot français (inventé par les modèles)
\providecommand{\fr}[1]{\textit{#1}}
% zhbox : encadré de texte chinois inventé par les modèles
\newenvironment{zhbox}{\begin{quote}}{\end{quote}}
% \vs : « versus » inventé par les modèles
\providecommand{\vs}{\textit{vs}\ }
% \blank : case à compléter inventée par les modèles
\providecommand{\blank}[1][]{\underline{\hspace{1.6em}}}
\newtcolorbox{exemplebox}[1][]{popbase,colframe=poporange,before upper={\popboxhead{Exemple}{poporange}{#1}}}
\newtcolorbox{definition}[1][]{popbase,colframe=popblue,before upper={\popboxhead{Définition}{popblue}{#1}}}
\newtcolorbox{propriete}[1][]{popbase,colframe=poppurple,before upper={\popboxhead{Propriété}{poppurple}{#1}}}
\newtcolorbox{methode}[1][]{popbase,colframe=popgold,before upper={\popboxhead{Méthode}{popgold}{#1}}}
\newtcolorbox{attention}[1][]{popbase,colframe=poppink,before upper={\popboxhead{Attention}{poppink}{#1}}}
\newtcolorbox{experience}[1][]{popbase,colframe=popblue,colback=popblueL,before upper={\popboxhead{Expérience}{popblue}{#1}}}
% « Le savais-tu ? » : carte violette pleine (contexte, culture, Cameroun)
\newtcolorbox{savaistu}[1][]{popbase,colback=poppurple,colframe=popink,
  fontupper=\popsemi\fontsize{10.4}{14.2}\selectfont\color{white},
  before upper={\poppill{Le savais-tu\,?}{white}{poppurple}\ifx\relax#1\relax\else\hspace{6pt}{\popxbold\color{white}#1}\fi\par\vspace{7pt}}}
% Exercice résolu : énoncé en haut, \tcblower puis la solution sur fond crème
\newcounter{popexo}\newcounter{popexr}
\newtcolorbox{exoresolu}[1][]{popbase,colframe=poporange,colbacklower=popcream,
  segmentation style={popink,line width=1.2pt,dashed},
  before upper={\stepcounter{popexr}\popboxhead{Exercice résolu \thepopexr}{poporange}{#1}},
  before lower={\poppill{Solution}{popink}{popcream}\par\vspace{6pt}}}
% Exercice (non résolu en place) — la correction est dans la partie Corrigés
\newtcolorbox{exercice}[1][]{popbase,colframe=popink,
  before upper={\stepcounter{popexo}\popboxhead{Exercice \thepopexo}{popink}{#1}}}
% Corrigé
\newtcolorbox{corrige}[1][]{popbase,colframe=popgreen,colback=popgreenL,
  before upper={\popboxhead{Corrigé}{popgreen}{#1}}}
% Objectifs / prérequis / bilan
\newtcolorbox{objectifs}{popbase,colframe=popink,colback=poporangeL,
  before upper={\popboxhead{Objectifs de la leçon}{popink}{}}}
\newtcolorbox{prerequis}{popbase,colframe=popink,colback=popblueL,
  before upper={\popboxhead{Prérequis}{popblue}{}}}
\newtcolorbox{bilan}{popbase,colframe=popink,colback=white,boxrule=2.8pt,
  before upper={\popboxhead{Fiche bilan}{poporange}{L'essentiel à connaître pour l'examen}}}
% Figure encadrée avec légende
\newtcolorbox{popfigure}[1][]{enhanced,breakable=false,colback=white,colframe=popline,boxrule=1.4pt,arc=8pt,
  left=10pt,right=10pt,top=10pt,bottom=8pt,before skip=10pt,after skip=12pt,halign=center,
  lower separated=false,halign lower=center,
  fontlower=\popsemi\fontsize{9}{11.5}\selectfont\color{popmuted},
  #1}
\newcommand{\legende}[1]{\par\vspace{6pt}{\popsemi\fontsize{9}{11.5}\selectfont\color{popmuted}#1}}

% ============================================================
% Signature OnBuch+
% ============================================================
\newcommand{\poplogoinline}[1]{{\poplogo\fontsize{#1}{#1}\selectfont\color{popink}OnBuch\color{poporange}+}}
\newcommand{\signatureonbuch}{%
  \begin{tcolorbox}[enhanced,colback=popink,colframe=popink,boxrule=2.2pt,arc=10pt,
      drop shadow=poporange,left=14pt,right=14pt,top=10pt,bottom=10pt,before skip=0pt,after skip=0pt]
    \tikz[baseline=-0.7ex]\node[circle,fill=popgreen,draw=popcream,line width=1.3pt,minimum size=22pt,inner sep=0]{\popcheck{white}};%
    \hspace{9pt}\begin{minipage}[c]{0.62\linewidth}
      {\popxbold\fontsize{12}{13}\selectfont\color{popcream}Signé\ \textbullet\ {\poplogo OnBuch\color{poporange}+}}\par\vspace{2pt}
      {\raggedright\popsemi\fontsize{8}{10.5}\selectfont\color{popline}\MakeUppercase{Équipe pédagogique OnBuch+}\\\MakeUppercase{Conforme au programme officiel MINESEC}\par}
    \end{minipage}\hfill
    {\poplogo\fontsize{22}{22}\selectfont\color{popcream}OnBuch\color{poporange}+}
  \end{tcolorbox}%
}

% ============================================================
% Page de garde
% #1 titre · #2 matière · #3 niveau · #4 module · #5 n° leçon · #6 durée ·
% #7 description / objectifs (texte ou itemize)
% ============================================================
\newcommand{\pagegarde}[7]{%
  \thispagestyle{empty}
  \newgeometry{a4paper,margin=1.7cm,top=1.6cm,bottom=1.4cm}%
  \begin{tikzpicture}[remember picture,overlay]
    \fill[poporange!18] ([xshift=-3.2cm,yshift=-3.6cm]current page.north east) circle (4.6cm);
    \fill[poppurple!14] ([xshift=2.4cm,yshift=7.6cm]current page.south west) circle (3.8cm);
  \end{tikzpicture}%
  {\poplogo\fontsize{30}{30}\selectfont\color{popink}OnBuch\color{poporange}+}\hfill
  \raisebox{0.6ex}{\poppill{Cours signé \textbullet\ Premium}{popink}{popcream}}\par
  \vspace{1pt}\popsquiggle{poppurple}\par
  \vspace{8pt}
  \begin{tcolorbox}[enhanced,colback=poporange,colframe=popink,boxrule=2.8pt,arc=13pt,
      drop shadow=popink,left=18pt,right=18pt,top=12pt,bottom=13pt,after skip=10pt]
    \poppill{#2 \textbullet\ #3}{popink}{popcream}\hspace{5pt}\poppill{#4}{white}{popink}\par\vspace{8pt}
    {\popdisplay\fontsize{14}{14}\selectfont\color{popink}LEÇON #5}\par\vspace{4pt}
    {\raggedright\hyphenpenalty=10000\exhyphenpenalty=10000\popdisplay\fontsize{25}{27}\selectfont\color{popcream}\MakeUppercase{#1}\par}
    \vspace{8pt}
    {\popxbold\fontsize{9.5}{9.5}\selectfont\color{popink}\MakeUppercase{Durée conseillée : #6}}
  \end{tcolorbox}
  \begin{tcolorbox}[enhanced,colback=white,colframe=popink,boxrule=2.2pt,arc=10pt,
      drop shadow=popink,left=16pt,right=16pt,top=10pt,bottom=10pt,after skip=8pt]
    \poppill{Dans cette leçon}{poppurple}{white}\par\vspace{5pt}
    \popsemi\fontsize{9.3}{12.6}\selectfont\color{popink}#7
  \end{tcolorbox}
  \noindent\begin{minipage}[t]{0.487\linewidth}
    \begin{tcolorbox}[enhanced,colback=popgreen,colframe=popink,boxrule=2.2pt,arc=10pt,
        drop shadow=popink,left=11pt,right=11pt,top=10pt,bottom=10pt,height=3.1cm,valign=center]
      \tcbox[colback=white,colframe=popink,boxrule=1.3pt,arc=5pt,boxsep=0pt,nobeforeafter,
        left=3pt,right=3pt,top=3pt,bottom=3pt,valign=center]{\qrcode[height=1.9cm]{https://play.google.com/store/apps/details?id=com.luvvix.onbuch&hl=fr}}%
      \hspace{8pt}\begin{minipage}[c]{0.5\linewidth}
        {\popdisplay\fontsize{11}{12.5}\selectfont\color{white}\MakeUppercase{Télécharge OnBuch}}\par\vspace{4pt}
        {\raggedright\popsemi\fontsize{8.4}{11}\selectfont\color{white}Gratuit \textbullet\ Google Play\\Tuteur IA, annales, concours, résultats\par}
      \end{minipage}
    \end{tcolorbox}
  \end{minipage}\hfill
  \begin{minipage}[t]{0.487\linewidth}
    \begin{tcolorbox}[enhanced,colback=poppurple,colframe=popink,boxrule=2.2pt,arc=10pt,
        drop shadow=popink,left=11pt,right=11pt,top=10pt,bottom=10pt,height=3.1cm,valign=center]
      \tikz[baseline=-0.7ex]\node[circle,fill=white,draw=popink,line width=1.3pt,minimum size=1.6cm,inner sep=0]{\popdisplay\fontsize{24}{24}\selectfont\color{poppurple}+};%
      \hspace{8pt}\begin{minipage}[c]{0.6\linewidth}
        {\popdisplay\fontsize{11}{12.5}\selectfont\color{white}\MakeUppercase{Passe à OnBuch+}}\par\vspace{4pt}
        {\raggedright\popsemi\fontsize{8.4}{11}\selectfont\color{white}Tous les cours signés, Tuteur IA illimité, fiches PDF — onglet OnBuch+ de l'app\par}
      \end{minipage}
    \end{tcolorbox}
  \end{minipage}
  \vspace{8pt}
  \popcontenu
  \vfill
  \signatureonbuch
  \restoregeometry
  \newpage
}

% Rangée « Ce cours contient » (page de garde)
\newcommand{\poptile}[3]{% #1 couleur #2 gros texte #3 légende
  \begin{tcolorbox}[enhanced,colback=#1,colframe=popink,boxrule=2pt,arc=9pt,shadow={2.5pt}{-2.5pt}{0pt}{popink},
      left=6pt,right=6pt,top=9pt,bottom=9pt,halign=center,valign=center,height=1.75cm,width=\linewidth,nobeforeafter]
    {\popdisplay\fontsize{12.5}{13}\selectfont\color{white}\MakeUppercase{#2}}\par\vspace{3pt}
    {\centering\popsemi\fontsize{7.8}{9.5}\selectfont\color{white}#3\par}
  \end{tcolorbox}}
\newcommand{\popcontenu}{%
  \par\noindent{\popxboldls\fontsize{7.5}{7.5}\selectfont\color{popmuted}CE COURS SIGNÉ CONTIENT}\par\vspace{7pt}
  \noindent\begin{minipage}[t]{0.235\linewidth}\poptile{popblue}{Cours}{complet, illustré, pas à pas}\end{minipage}\hfill
  \begin{minipage}[t]{0.235\linewidth}\poptile{poporange}{Méthodes}{et exercices résolus}\end{minipage}\hfill
  \begin{minipage}[t]{0.235\linewidth}\poptile{poppink}{Exercices}{type examen + corrigés}\end{minipage}\hfill
  \begin{minipage}[t]{0.235\linewidth}\poptile{popgreen}{Fiche bilan}{l'essentiel à retenir}\end{minipage}\par}

% Sommaire
\newtcolorbox{sommaire}{enhanced,colback=white,colframe=popink,boxrule=2.2pt,arc=10pt,
  drop shadow=popink,left=18pt,right=18pt,top=13pt,bottom=13pt,before skip=6pt,after skip=20pt,
  before upper={\poppill{Sommaire}{poppurple}{white}\par\vspace{9pt}\popsemi\fontsize{10.6}{14.5}\selectfont\color{popink}}}

% Signature de fin de cours — poussée tout en bas de la dernière page
\newcommand{\findecours}{%
  \par\vfill\nopagebreak
  \begin{center}\popsquiggle{poporange}\end{center}\vspace{4pt}
  \signatureonbuch
}
\AtBeginDocument{\def\llbracket{\mathopen{[\mkern-3.5mu[}}\def\rrbracket{\mathclose{]\mkern-3.5mu]}}} % intervalles d'entiers (glyphe ⟦ absent de la police)
% Glyphes absents de Fira Math : redéfinis à partir de glyphes disponibles
\AtBeginDocument{%
  \def\setminus{\mathbin{\backslash}}\let\smallsetminus\setminus
  \def\vdots{\mathord{\vbox{\baselineskip3.2pt\lineskiplimit0pt\kern4pt\hbox{.}\hbox{.}\hbox{.}}}}%
  \def\ddots{\mathinner{\mkern1mu\raise6.5pt\hbox{.}\mkern2mu\raise3.6pt\hbox{.}\mkern2mu\raise0.7pt\hbox{.}\mkern1mu}}%
  \def\triangle{\mathord{\scalebox{0.9}{$\symup{\Delta}$}}}%
  \def\nabla{\mathord{\raisebox{\depth}{\scalebox{1}[-1]{$\symup{\Delta}$}}}}%
  \def\checkmark{\popcheck{popdark}}%
  \def\star{\mathbin{\ast}}\def\bigstar{\mathord{\ast}}%
  \def\diamond{\mathbin{\scalebox{0.6}{$\Diamond$}}}%
  \def\bigcup{\mathop{\mathchoice{\raisebox{-0.3em}{\scalebox{1.7}{$\cup$}}}{\scalebox{1.25}{$\cup$}}{\cup}{\cup}}\displaylimits}%
  \def\bigcap{\mathop{\mathchoice{\raisebox{-0.3em}{\scalebox{1.7}{$\cap$}}}{\scalebox{1.25}{$\cap$}}{\cap}{\cap}}\displaylimits}%
  \def\lesseqqgtr{\mathrel{\lessgtr}}\def\bigoplus{\mathop{\oplus}\displaylimits}%
}
\providecommand{\ssi}{si et seulement si} % abréviation fréquente
\AtBeginDocument{\providecommand{\wideparen}{\overparen}} % arc de cercle

%% FIGFIX-BEGIN v4 — correctifs globaux des figures (docs/onbuch-generation/tools/figfix)
\usepackage{adjustbox}\usepackage{etoolbox}
% toute figure TikZ/pgfplots plus large que la ligne est réduite à la largeur disponible
\BeforeBeginEnvironment{tikzpicture}{\begin{adjustbox}{max width=\linewidth}}
\AfterEndEnvironment{tikzpicture}{\end{adjustbox}}
% légendes pgfplots lisibles par-dessus les courbes
\pgfplotsset{every axis/.append style={legend style={fill=white,fill opacity=0.92,text opacity=1,draw=black!15,inner sep=3pt}}}
% étiquettes d'arêtes (syntaxe edge["..."]) avec fond blanc
\tikzset{every edge quotes/.append style={fill=white,inner sep=1.2pt}}
%% FIGFIX-END
\begin{document}

\pagegarde{Leçon 13 — Vocabulaire et compréhension de texte : maladies en général et du SIDA en particulier}{Chinois}{1ère A}{Module 2 : Environnement et santé}{ZHA13}{2 h}{Cette leçon de 2 heures propose l'étude d'un texte authentique sur les maladies générales et le SIDA, l'acquisition d'un lexique médical ciblé (symptômes, prévention, terminologie SIDA), l'analyse de radicaux-clés (疒, 月/肉) et un entraînement intensif à la compréhension écrite et à la production guidée.
\vspace{4pt}
\begin{multicols}{2}\small
\begin{itemize}
\item Lire à voix haute le texte support avec une prononciation correcte (pinyin, tons) et une intonation naturelle.
\item Identifier et traduire le vocabulaire clé du thème : maladies courantes, symptômes, SIDA, prévention, transmission.
\item Reconnaître, écrire dans l'ordre des traits et classifier par radical 8 à 10 caractères nouveaux du thème.
\item Répondre par écrit en chinois (ou en français) à des questions de compréhension globale et détaillée sur le texte.
\item Utiliser les collocations verbales médicales (患病, 感染, 预防, 治疗) dans des phrases simples.
\item Produire un court paragraphe (5-6 phrases) sur la prévention d'une maladie en réinvestissant le lexique et les structures vues.
\end{itemize}\end{multicols}}

\begin{sommaire}
\begin{multicols}{2}
\begin{enumerate}
\item 1. Texte support : « 认识疾病，预防艾滋 »
\item 2. Glossaire thématique : Tableau de vocabulaire maître
\item 3. Clés (Radicaux) et Ordre des traits : Les briques du corps médical
\item 4. Compréhension de texte : Questions graduelles
\item 5. Collocations, Structures utiles et Production guidée
\item 6. Approfondissement socioculturel : Cameroun - Chine, regards croisés
\item Activité d'intégration
\item Méthodes et exercices résolus
\item Exercices
\item Corrigés des exercices
\item Fiche bilan
\end{enumerate}
\end{multicols}
\end{sommaire}

\begin{objectifs}
\begin{itemize}
\item Lire à voix haute le texte support avec une prononciation correcte (pinyin, tons) et une intonation naturelle.
\item Identifier et traduire le vocabulaire clé du thème : maladies courantes, symptômes, SIDA, prévention, transmission.
\item Reconnaître, écrire dans l'ordre des traits et classifier par radical 8 à 10 caractères nouveaux du thème.
\item Répondre par écrit en chinois (ou en français) à des questions de compréhension globale et détaillée sur le texte.
\item Utiliser les collocations verbales médicales (患病, 感染, 预防, 治疗) dans des phrases simples.
\item Produire un court paragraphe (5-6 phrases) sur la prévention d'une maladie en réinvestissant le lexique et les structures vues.
\item Différencier les termes proches : 病 vs 症, 传染 vs 感染, 预防 vs 控制.
\end{itemize}
\end{objectifs}

\begin{prerequis}
\begin{itemize}
\item Structure 基本句型 : Sujet + 时间 + 地点 + 动词 + 宾语 (ordre des mots de base).
\item Vocabulaire : Parties du corps (头, 手, 脚, 血, 身体) et adjectifs d'état (好, 坏, 严重, 轻).
\item Grammaire : Particule 了 (aspect accompli), 过 (expérience), 正在/在 (progressif), 得 (complément de degré).
\item Radicaux de base : 亻 (personne), 氵 (eau), 忄 (cœur), 口 (bouche).
\item Compétence : Lire un texte court (50-80 caractères) avec pinyin et en extraire l'idée principale.
\end{itemize}
\end{prerequis}

\newpage

\coursec{Texte support : « 认识疾病，预防艾滋 »}

\courssub{Situation d'accroche et problématique}

Au Cameroun, comme dans de nombreux pays africains, la sensibilisation aux maladies transmissibles, notamment le VIH/SIDA, reste un enjeu majeur de santé publique. À Yaoundé comme à Bamenda, les campagnes de dépistage, les affiches dans les centres de santé et les émissions de radio rappellent chaque jour l'importance de la prévention. Imaginez maintenant un élève camerounais qui, lors d'un échange universitaire à Pékin, doit lire une brochure médicale chinoise ou expliquer des symptômes à un médecin. La maîtrise du vocabulaire médical de base n'est pas seulement scolaire : elle est \cle{vitale}.

\textbf{Problématique :} comment lire, comprendre et restituer en chinois un court texte de santé publique portant sur les maladies en général et sur le SIDA en particulier ? Pour y répondre, nous allons : (1) découvrir un texte authentique, (2) en dégager les mots nouveaux (\zh{生词}), (3) analyser sa structure et ses connecteurs logiques, (4) nous entraîner à le traduire.

\courssub{Le texte support : lecture et découverte}

\begin{textebox}[认识疾病，预防艾滋 — Rènshi jíbìng, yùfáng Àizī — « Connaître les maladies, prévenir le SIDA »]
疾病是人体正常功能受到破坏的状态。\\
\emph{Jíbìng shì réntǐ zhèngcháng gōngnéng shòudào pòhuài de zhuàngtài.}\\
« La maladie est un état où les fonctions normales du corps humain sont détruites. »\\[4pt]
有些疾病会传染，比如感冒和艾滋病。\\
\emph{Yǒuxiē jíbìng huì chuánrǎn, bǐrú gǎnmào hé Àizībìng.}\\
« Certaines maladies sont contagieuses, comme le rhume et le SIDA. »\\[4pt]
艾滋病是由病毒引起的严重疾病。\\
\emph{Àizībìng shì yóu bìngdú yǐnqǐ de yánzhòng jíbìng.}\\
« Le SIDA est une maladie grave causée par un virus. »\\[4pt]
它主要通过血液、母婴和性接触传播。\\
\emph{Tā zhǔyào tōngguò xuèyè, mǔyīng hé xìng jiēchù chuánbō.}\\
« Il se transmet principalement par le sang, de la mère à l'enfant et par les rapports sexuels. »\\[4pt]
但是，握手、一起吃饭不会传染艾滋病。\\
\emph{Dànshì, wòshǒu, yìqǐ chīfàn bú huì chuánrǎn Àizībìng.}\\
« Mais se serrer la main ou manger ensemble ne transmet pas le SIDA. »\\[4pt]
因此，预防为主：要注意安全，尤其要保护年轻人。\\
\emph{Yīncǐ, yùfáng wéi zhǔ : yào zhùyì ānquán, yóuqí yào bǎohù niánqīngrén.}\\
« Par conséquent, la prévention est primordiale : il faut faire attention à la sécurité, et surtout protéger les jeunes. »\\[4pt]
总之，我们要关爱病人，不能歧视他们。\\
\emph{Zǒngzhī, wǒmen yào guān'ài bìngrén, bù néng qíshì tāmen.}\\
« En résumé, nous devons prendre soin des malades et ne pas les discriminer. »
\end{textebox}

\begin{definition}[生词 — les mots nouveaux]
Les \cle{生词} (\emph{shēngcí}, littéralement « mots vivants/nouveaux ») sont les mots de vocabulaire inconnu à maîtriser \textbf{prioritairement} pour la compréhension du texte. Dans une leçon de lecture, on les repère à la première lecture, on les surligne, puis on les mémorise avec leur pinyin, leurs tons et leur traduction. Ils feront l'objet du tableau de vocabulaire maître (section 2).
\end{definition}

\courssub{Lecture modèle : tons, liaisons et pauses}

La lecture à voix haute d'un texte chinois obéit à des règles précises. L'enseignant lit d'abord le texte entier à vitesse normale ; les élèves écoutent, livre fermé, puis livre ouvert en suivant du doigt.

\begin{propriete}[Règles de lecture expressive]
\begin{itemize}
\item \textbf{Pauses syntaxiques} : on marque une courte pause après la virgule chinoise \zh{，} et une pause longue après le point \zh{。} ou les deux-points \zh{：}. Exemple : \zh{因此，预防为主：要注意安全。} — pause après \zh{因此}, pause longue après \zh{为主}.
\item \textbf{Groupes de sens} : on ne coupe jamais un mot dissyllabique : \zh{疾病} (\emph{jíbìng}) se lit d'un bloc, jamais \emph{jí... bìng}.
\item \textbf{Sandhi du troisième ton} : deux 3\textsuperscript{e} tons qui se suivent : le premier devient un 2\textsuperscript{e} ton. Exemple : \zh{可以} se prononce \emph{kúyǐ}. Dans notre texte : \zh{主要} (\emph{zhǔyào}) ne pose pas de problème, mais \zh{所以} (\emph{suǒyǐ}) se prononce \emph{suóyǐ}.
\item \textbf{变调 de 不} : \zh{不} (4\textsuperscript{e} ton) devient 2\textsuperscript{e} ton devant un 4\textsuperscript{e} ton : \zh{不会} se prononce \emph{bú huì}, \zh{不能} reste \emph{bù néng}.
\item \textbf{变调 de 一} : \zh{一起} se prononce \emph{yìqǐ} (4\textsuperscript{e} ton devant un 3\textsuperscript{e} ton).
\end{itemize}
\end{propriete}

\begin{experience}[Lecture chorale puis individuelle]
\begin{enumerate}
\item \textbf{Lecture chorale} : la classe répète phrase par phrase après l'enseignant, en insistant sur les mots nouveaux : \zh{疾病}, \zh{病毒}, \zh{传染}, \zh{传播}, \zh{歧视}.
\item \textbf{Focus sur les groupes nominaux complexes} : faire répéter \zh{由病毒引起的严重疾病} (\emph{yóu bìngdú yǐnqǐ de yánzhòng jíbìng}, « maladie grave causée par un virus ») en le découpant : \zh{由病毒引起} + \zh{的} + \zh{严重疾病}.
\item \textbf{Lecture individuelle} : chaque élève lit un paragraphe ; les autres vérifient les tons et les pauses.
\end{enumerate}
\end{experience}

\courssub{Traduction littérale et communicative}

\begin{methode}[Lecture en 3 temps]
\begin{enumerate}
\item \textbf{Lecture rapide (silencieuse)} : parcourir le texte une fois pour saisir l'idée générale, sans s'arrêter sur les mots inconnus. Question-guide : « De quoi parle ce texte ? »
\item \textbf{Lecture détaillée} : surligner les \zh{生词}, cercler les radicaux connus (par exemple la clé de la maladie \zh{疒} dans \zh{疾病}), noter les connecteurs.
\item \textbf{Lecture à voix haute} : relire pour la fluidité, en respectant tons et pauses, jusqu'à pouvoir lire sans hésitation.
\end{enumerate}
\end{methode}

Pour traduire, on distingue deux niveaux :

\begin{itemize}
\item \textbf{Traduction littérale} : elle suit l'ordre chinois. \zh{它主要通过血液传播} → « Il principalement par le sang se transmet. » Elle sert à vérifier la compréhension de la structure.
\item \textbf{Traduction communicative} : elle reformule en français naturel → « Il se transmet principalement par le sang. »
\end{itemize}

\begin{attention}[L'ordre des compléments : piège majeur]
En chinois, le complément de moyen introduit par \zh{通过} (\emph{tōngguò}, « par, au moyen de ») se place \textbf{avant le verbe}, alors qu'en français il se place après :
\begin{itemize}
\item \zh{通过血液传播} — \emph{tōngguò xuèyè chuánbō} — « se transmettre \textbf{par le sang} » (et non « transmettre par le sang » dans cet ordre en chinois).
\item \zh{由病毒引起} — \emph{yóu bìngdú yǐnqǐ} — « causé \textbf{par un virus} » : \zh{由} marque l'agent et précède le verbe.
\end{itemize}
Ne dites jamais \zh{传播通过血液} : c'est une erreur typique des francophones qui calquent l'ordre français.
\end{attention}

\begin{exemplebox}[Phrases clés traduites]
\begin{itemize}
\item \zh{因此，预防为主。} — \emph{Yīncǐ, yùfáng wéi zhǔ.} — « Par conséquent, la prévention est primordiale. » (litt. « la prévention fait office de principal »)
\item \zh{我们要关爱病人。} — \emph{Wǒmen yào guān'ài bìngrén.} — « Nous devons prendre soin des malades / leur témoigner de la compassion. »
\item \zh{不能歧视他们。} — \emph{Bù néng qíshì tāmen.} — « On ne doit pas les discriminer. »
\item \zh{尤其要保护年轻人。} — \emph{Yóuqí yào bǎohù niánqīngrén.} — « Il faut surtout protéger les jeunes. »
\end{itemize}
\end{exemplebox}

\courssub{Analyse de la structure textuelle}

Le texte suit une progression classique du texte explicatif-argumentatif chinois : du \textbf{général} (les maladies) vers le \textbf{spécifique} (le SIDA), puis vers un \textbf{appel à l'action}.

\begin{popfigure}
\begin{tikzpicture}[
node distance=0.35cm,
box/.style={draw=popblue, fill=popblueL, rounded corners=2pt, align=center, text width=10.5cm, inner sep=4pt, font=\small},
spec/.style={draw=poporange, fill=poporangeL, rounded corners=2pt, align=center, text width=10.5cm, inner sep=4pt, font=\small},
conc/.style={draw=popgreen, fill=popgreenL, rounded corners=2pt, align=center, text width=10.5cm, inner sep=4pt, font=\small},
fleche/.style={-{Stealth[length=2mm]}, thick, popdark}]
\node[box] (titre) {\textbf{Titre} : 认识疾病，预防艾滋 (Connaître les maladies, prévenir le SIDA)};
\node[box, below=of titre] (p1) {\textbf{§1} Maladies en général : définition + exemples de maladies contagieuses};
\node[spec, below=of p1] (p2) {\textbf{§2} Qu'est-ce que le SIDA ? Maladie grave causée par un virus};
\node[spec, below=of p2] (p3) {\textbf{§3} Transmission : sang, mère-enfant, rapports sexuels (mais pas la poignée de main !)};
\node[spec, below=of p3] (p4) {\textbf{§4} Prévention : sécurité, protection des jeunes};
\node[conc, below=of p4] (p5) {\textbf{Conclusion} : appel à l'action — aimer les malades, refuser la discrimination};
\draw[fleche] (titre) -- (p1);
\draw[fleche] (p1) -- (p2);
\draw[fleche] (p2) -- (p3);
\draw[fleche] (p3) -- (p4);
\draw[fleche] (p4) -- (p5);
\end{tikzpicture}
\legende{Organigramme de la structure du texte : du général au particulier, puis à l'appel à l'action.}
\end{popfigure}

\courssub{Les connecteurs logiques du texte}

Les connecteurs sont les « charnières » du texte : ils indiquent la relation logique entre les phrases. Les reconnaître, c'est comprendre deux fois plus vite.

\begin{center}
\begin{tabularx}{\linewidth}{|X|X|X|}
\hline
\rowcolor{popblueL} Connecteur & Emploi et exemple & Traduction \\
\hline
因此 yīncǐ & conséquence : \zh{因此，预防为主。} & par conséquent, donc \\
\hline
但是 dànshì & opposition : \zh{但是，握手不会传染。} & mais, cependant \\
\hline
尤其 yóuqí & mise en relief : \zh{尤其要保护年轻人。} & surtout, en particulier \\
\hline
总之 zǒngzhī & conclusion : \zh{总之，我们要关爱病人。} & en résumé, bref \\
\hline
比如 bǐrú & exemple : \zh{比如感冒和艾滋病。} & par exemple, comme \\
\hline
主要 zhǔyào & restriction : \zh{主要通过血液传播。} & principalement \\
\hline
\end{tabularx}
\end{center}

\begin{attention}[艾滋病 ou VIH ?]
Ne pas confondre \zh{艾滋病} (\emph{Àizībìng}, le SIDA — le \textbf{syndrome}, la maladie déclarée) et le VIH (le \textbf{virus}, en chinois \zh{艾滋病毒} ou \emph{HIV}). En chinois courant, on utilise souvent \zh{艾滋病} pour désigner l'infection dans son ensemble, mais le texte médical distingue soigneusement le \zh{病毒} (virus) de la \zh{病} (maladie). Dans notre texte : \zh{艾滋病是由病毒引起的} — « le SIDA est causé par un virus » : les deux notions sont bien séparées.
\end{attention}

\begin{exoresolu}[Repérage des connecteurs]
Énoncé : dans la phrase \zh{艾滋病很严重，但是一起吃饭不会传染，因此不用害怕病人。}, relevez les connecteurs, donnez leur fonction logique, puis traduisez la phrase en français.
\tcblower
Solution détaillée :
\begin{enumerate}
\item \zh{但是} (\emph{dànshì}) : connecteur d'\textbf{opposition} — il oppose la gravité de la maladie à l'absence de risque dans la vie quotidienne.
\item \zh{因此} (\emph{yīncǐ}) : connecteur de \textbf{conséquence} — il tire la conclusion logique des deux constats précédents.
\item Traduction : « Le SIDA est très grave, \textbf{mais} manger ensemble ne le transmet pas ; \textbf{par conséquent}, il n'y a pas lieu d'avoir peur des malades. »
\end{enumerate}
Remarque : en français, on ponctue avec un point-virgule avant « par conséquent » ; en chinois, \zh{因此} s'introduit simplement après une virgule \zh{，}.
\end{exoresolu}

\begin{savaistu}[Le ruban rouge 红丝带]
Le ruban rouge (\zh{红丝带}, \emph{hóng sīdài}) est le symbole international de la lutte contre le SIDA, adopté en Chine depuis les années 1990. Chaque 1\textsuperscript{er} décembre, Journée mondiale de lutte contre le SIDA (\zh{世界艾滋病日}, \emph{shìjiè Àizībìng rì}), on le voit épinglé sur les vêtements, à Pékin comme à Yaoundé. En chinois, le caractère \zh{丝} (\emph{sī}, « soie ») rappelle la fameuse Route de la soie : un beau pont symbolique entre la Chine et le monde.
\end{savaistu}

\begin{exercice}[À vous de jouer]
\begin{enumerate}
\item Lisez le texte à voix haute en respectant les pauses et le 变调 de \zh{不} dans \zh{不会} et \zh{不能}.
\item Soulignez dans le texte les quatre connecteurs \zh{因此}, \zh{但是}, \zh{尤其}, \zh{总之} et indiquez leur fonction.
\item Traduisez en français communicatif : \zh{它主要通过血液、母婴和性接触传播。}
\item Répondez en chinois : \zh{握手会传染艾滋病吗？} (Se serrer la main transmet-il le SIDA ?)
\end{enumerate}
\end{exercice}

\begin{corrige}[Exercice « À vous de jouer »]
\begin{enumerate}
\item \zh{不会} se prononce \emph{bú huì} (2\textsuperscript{e} ton) ; \zh{不能} se prononce \emph{bù néng} (4\textsuperscript{e} ton conservé devant un 2\textsuperscript{e} ton).
\item \zh{但是} : opposition ; \zh{因此} : conséquence ; \zh{尤其} : mise en relief ; \zh{总之} : conclusion.
\item « Il se transmet principalement par le sang, de la mère à l'enfant et par les rapports sexuels. »
\item \zh{不会。握手不会传染艾滋病。} — \emph{Bú huì. Wòshǒu bú huì chuánrǎn Àizībìng.} — « Non. Se serrer la main ne transmet pas le SIDA. »
\end{enumerate}
\end{corrige}

\begin{aretenir}[L'essentiel de la section]
\begin{itemize}
\item Le texte « 认识疾病，预防艾滋 » va du général (les maladies) au particulier (le SIDA : définition, transmission, prévention) et se termine par un appel à la non-discrimination.
\item Quatre connecteurs à maîtriser : \zh{因此} (conséquence), \zh{但是} (opposition), \zh{尤其} (mise en relief), \zh{总之} (conclusion).
\item Les compléments introduits par \zh{通过} et \zh{由} se placent \textbf{avant} le verbe en chinois.
\item Méthode de lecture en 3 temps : rapide, détaillée, à voix haute.
\item \zh{艾滋病} = le SIDA (syndrome) ; \zh{病毒} = le virus : ne pas les confondre.
\end{itemize}
\end{aretenir}

\coursec{Leçon 13 — Vocabulaire et compréhension de texte : maladies en général et du SIDA en particulier}

\coursec{Texte support : « 认识疾病，预防艾滋 »}

\begin{textebox}[Texte original : 认识疾病，预防艾滋]
\zh{认识疾病，预防艾滋} \\[4pt]
\zh{疾病是人类健康的敌人。常见的疾病有感冒、发烧、咳嗽、头疼等，症状通常比较轻微，经过治疗和休息就能好转。但是，也有一些严重的传染病，比如艾滋病。} \\[4pt]
\zh{艾滋病是由艾滋病毒（HIV）引起的。病毒进入人体后，攻击免疫系统，使人体失去抵抗疾病的能力。艾滋病主要通过三条途径传播：血液传播、性传播和母婴传播。日常接触，如握手、拥抱、共用餐具、蚊虫叮咬，不会传染艾滋病。} \\[4pt]
\zh{预防艾滋病，关键在于避免高危行为。要做到：不吸毒，不共用针头；性行为时正确使用安全套；接受输血或注射时，确保血液和器械安全；孕妇做检查，阻断母婴传播。定期做艾滋病检查，早发现、早治疗，患者也能像正常人一样生活、工作。} \\[4pt]
\zh{面对艾滋病患者，我们不应歧视，而应给予关爱和支持。预防疾病，保护健康，是我们每一个人的责任。} \\[8pt]
\emph{Rènshi jíbìng, yùfáng Àizī.} \\[2pt]
\emph{Jíbìng shì rénlèi jiànkāng de dírén. Chángjiàn de jíbìng yǒu gǎnmào, fāshāo, késou, tóuténg děng, zhèngzhuàng tōngcháng bǐjiào qīngwēi, jīngguò zhìliáo xiūxi jiù néng hǎozhuǎn. Dànshì, yě yǒu yīxiē yánzhòng de chuánrǎnbìng, bǐrú Àizībìng.} \\[4pt]
\emph{Àizībìng shì yóu Àizī bìngdú (HIV) yǐnqǐ de. Bìngdú jìnrù réntǐ hòu, gōngjī miǎnyì xìtǒng, shǐ réntǐ shīqù dǐkàng jíbìng de nénglì. Àizībìng zhǔyào tōngguò sān tiáo tújìng chuánbō: xuèyè chuánbō, xìng chuánbō hé mǔyīng chuánbō. Rìcháng jiēchù, rú wòshǒu, yǒngbào, gòngyòng cānjù, wénchóng dīngyǎo, bù huì chuánrǎn Àizībìng.} \\[4pt]
\emph{Yùfáng Àizībìng, guānjiàn zài yú bìmiǎn gāowēi xíngwéi. Yào dào: bù xīdú, bù gòngyòng zhēntóu; xìng xíngwéi shí zhèngquè shǐyòng ānquántào; jiēshòu shūxiě huò zhùshè shí, quèbǎo xuèyè hé qìxiè ānquán; yùnfù zuò jiǎnchá, zǔduàn mǔyīng chuánbō. Dìngqī zuò Àizībìng jiǎnchá, zǎo fāxiàn, zǎo zhìliáo, huànzhě yě néng xiàng zhèngcháng rén yīyàng shēnghuó, gōngzuò.} \\[4pt]
\emph{Miànduì Àizībìng huànzhě, wǒmen bù yīng qīshì, ér yīng gěiyǔ guānaì hé zhīchí. Yùfáng jíbìng, bǎohù jiànkāng, shì wǒmen měi yī gè rén de zérèn.} \\[8pt]
\textbf{Traduction :} \emph{Connaître les maladies, prévenir le sida.} \\
Les maladies sont les ennemies de la santé humaine. Les maladies courantes sont le rhume, la fièvre, la toux, le mal de tête, etc. Les symptômes sont généralement légers ; après traitement et repos, on guérit. Mais il existe de graves maladies infectieuses, comme le sida. \\
Le sida est causé par le virus du sida (VIH). Une fois dans le corps, le virus attaque le système immunitaire, faisant perdre au corps la capacité de résister aux maladies. Le se transmet principalement par trois voies : transmission par le sang, transmission sexuelle et transmission mère-enfant. Les contacts quotidiens — serrer la main, étreindre, partager des ustensiles, piqûres de moustiques — ne transmettent pas le sida. \\
Prévenir le sida, la clé est d'éviter les comportements à haut risque. Il faut : ne pas prendre de drogue, ne pas partager d'aiguilles ; utiliser correctement le préservatif lors des rapports sexuels ; lors de transfusions ou d'injections, garantir la sécurité du sang et du matériel ; les femmes enceintes doivent se faire dépister pour bloquer la transmission mère-enfant. Faire régulièrement le test du sida, dépister tôt, traiter tôt, les patients peuvent aussi vivre et travailler comme des personnes normales. \\
Face aux malades du sida, nous ne devons pas discriminer, mais apporter soin et soutien. Prévenir les maladies, protéger la santé, est le devoir de chacun de nous.
\end{textebox}

\coursec{Glossaire thématique : Tableau de vocabulaire maître}

Dans la section précédente, nous avons lu le texte \zh{认识疾病，预防艾滋} (\emph{Rènshi jíbìng, yùfáng Àizī} « Connaître les maladies, prévenir le sida »). Vous avez rencontré de nombreux mots nouveaux : \zh{疾病}, \zh{病毒}, \zh{传染}, \zh{预防}… Cette deuxième section les organise méthodiquement pour que vous puissiez les \cle{mémoriser}, les \cle{prononcer} et les \cle{employer} correctement. Un bon vocabulaire ne s'apprend pas mot par mot, mais par \cle{familles}, par \cle{collocations} et par \cle{radicaux}.

\courssub{Les quatre familles de mots du thème}

Avant de présenter le tableau, observons comment les mots se regroupent naturellement. Quand on parle de santé, on suit toujours le même chemin logique : la \textbf{maladie} apparaît (symptômes), on cherche sa \textbf{cause} (virus, transmission), on agit pour \textbf{prévenir}, puis pour \textbf{soigner}. Notre glossaire suit exactement cette logique en quatre catégories.

\begin{definition}[Caractère (字) et mot (词)]
En chinois, un \cle{caractère} (\zh{字} \emph{zì}) est une unité graphique et sonore ; un \cle{mot} (\zh{词} \emph{cí}) est souvent composé de \textbf{deux caractères}. Ainsi \zh{病} (\emph{bìng} « maladie ») est un caractère, mais le mot courant est \zh{病毒} (\emph{bìngdú} « virus »), composé de \zh{病} « maladie » + \zh{毒} « poison ». De même : \zh{免疫} (\emph{miǎnyì} « immunité » : \zh{免} « éviter » + \zh{疫} « épidémie »), \zh{系统} (\emph{xìtǒng} « système »), \zh{接触} (\emph{jiēchù} « contact » : \zh{接} « recevoir » + \zh{触} « toucher »). Apprenez toujours le \textbf{mot complet}, pas seulement le caractère isolé.
\end{definition}

\textbf{Catégorie 1 : Maladies et symptômes}\par

\begin{center}
\begin{tabularx}{\linewidth}{|X|X|X|X|X|}
\hline
\rowcolor{popblueL} Caractères & Pinyin & Nature & Traduction & Collocation \\
\hline
疾病 & jíbìng & n. & maladie (terme général) & 预防疾病 \\
\hline
病 & bìng & n./v. & maladie ; être malade & 看病 \\
\hline
症状 & zhèngzhuàng & n. & symptôme & 出现症状 \\
\hline
发烧 & fāshāo & v./n. & avoir de la fièvre & 发高烧 \\
\hline
咳嗽 & késou & v./n. & tousser ; toux & 咳嗽得很厉害 \\
\hline
感冒 & gǎnmào & v./n. & attraper un rhume & 得了感冒 \\
\hline
头疼 & tóuténg & v./adj. & avoir mal à la tête & 突然头疼 \\
\hline
严重 & yánzhòng & adj. & grave, sérieux & 病情严重 \\
\hline
轻微 & qīngwēi & adj. & léger, bénin & 症状轻微 \\
\hline
\end{tabularx}
\end{center}

\textbf{Catégorie 2 : SIDA et VIH}\par

\begin{center}
\begin{tabularx}{\linewidth}{|X|X|X|X|X|}
\hline
\rowcolor{popblueL} Caractères & Pinyin & Nature & Traduction & Collocation \\
\hline
艾滋病 & àizībìng & n. & sida & 患有艾滋病 \\
\hline
艾滋病毒 & àizī bìngdú & n. & VIH (virus du sida) & 感染艾滋病毒 \\
\hline
病毒 & bìngdú & n. & virus & 抵抗病毒 \\
\hline
免疫 & miǎnyì & n./v. & immunité ; immuniser & 免疫系统 \\
\hline
系统 & xìtǒng & n. & système & 免疫系统受损 \\
\hline
患者 & huànzhě & n. & malade, patient & 艾滋病患者 \\
\hline
检查 & jiǎnchá & v./n. & examiner ; dépistage & 做艾滋病检查 \\
\hline
\end{tabularx}
\end{center}

\textbf{Catégorie 3 : Transmission et prévention}\par

\begin{center}
\begin{tabularx}{\linewidth}{|X|X|X|X|X|}
\hline
\rowcolor{popblueL} Caractères & Pinyin & Nature & Traduction & Collocation \\
\hline
传染 & chuánrǎn & v./n. & transmettre ; contagion & 传染疾病 \\
\hline
遗传 & yíchuán & v./n. & hériter ; hérédité & 遗传疾病 \\
\hline
传播 & chuánbō & v. & propager, diffuser & 传播途径 \\
\hline
途径 & tújìng & n. & voie, canal & 传播途径 \\
\hline
血液 & xuèyè & n. & sang & 接触血液 \\
\hline
接触 & jiēchù & v./n. & toucher ; contact & 避免接触 \\
\hline
预防 & yùfáng & v. & prévenir & 预防艾滋 \\
\hline
安全 & ānquán & adj./n. & sûr ; sécurité & 安全措施 \\
\hline
高危 & gāowēi & adj. & à haut risque & 高危行为 \\
\hline
\end{tabularx}
\end{center}

\textbf{Catégorie 4 : Verbes d'action médicale}\par

\begin{center}
\begin{tabularx}{\linewidth}{|X|X|X|X|X|}
\hline
\rowcolor{popblueL} Caractères & Pinyin & Nature & Traduction & Collocation \\
\hline
治疗 & zhìliáo & v. & traiter, soigner & 治疗疾病 \\
\hline
看病 & kànbìng & v. & consulter un médecin & 去医院看病 \\
\hline
避免 & bìmiǎn & v. & éviter & 避免接触血液 \\
\hline
保护 & bǎohù & v. & protéger & 保护自己 \\
\hline
感染 & gǎnrǎn & v. & être infecté & 感染病毒 \\
\hline
患 & huàn & v. & souffrir de (maladie) & 患有艾滋病 \\
\hline
体检 & tǐjiǎn & v./n. & bilan de santé & 定期体检 \\
\hline
\end{tabularx}
\end{center}

\courssub{Paires de synonymes et d'antonymes}

Comprendre les oppositions aide à mémoriser deux mots à la fois. Voici les quatre paires essentielles de cette leçon.

\begin{exemplebox}[Paires contraires et proches]
\zh{健康} (\emph{jiànkāng} « santé ») contre \zh{疾病} (\emph{jíbìng} « maladie ») : \zh{预防疾病，保护健康。} \emph{Yùfáng jíbìng, bǎohù jiànkāng.} « Prévenir les maladies, protéger la santé. »\\[2pt]
\zh{传染} (\emph{chuánrǎn} « contagieux ») contre \zh{遗传} (\emph{yíchuán} « héréditaire ») : \zh{艾滋病是传染病，不是遗传病。} \emph{Àizībìng shì chuánrǎnbìng, bú shì yíchuánbìng.} « Le sida est une maladie contagieuse, pas héréditaire. »\\[2pt]
\zh{预防} (\emph{yùfáng} « prévenir ») contre \zh{治疗} (\emph{zhìliáo} « traiter ») : \zh{预防比治疗更重要。} \emph{Yùfáng bǐ zhìliáo gèng zhòngyào.} « La prévention est plus importante que le traitement. »\\[2pt]
\zh{严重} (\emph{yánzhòng} « grave ») contre \zh{轻微} (\emph{qīngwēi} « léger ») : \zh{他的病很严重，我的病很轻微。} \emph{Tā de bìng hěn yánzhòng, wǒ de bìng hěn qīngwēi.} « Sa maladie est grave, la mienne est légère. »
\end{exemplebox}

\begin{attention}[传染 ou 遗传 ?]
Les deux mots contiennent \zh{传} (\emph{chuán} « transmettre »), mais : \zh{传染} = transmission \textbf{entre personnes} (contagion : grippe, sida) ; \zh{遗传} = transmission \textbf{des parents aux enfants} (hérédité : couleur des yeux, certaines maladies génétiques). Ne dites jamais *\zh{艾滋病遗传} pour dire que le sida se transmet : dites \zh{艾滋病传染} ou précisez \zh{母婴传播} (transmission mère-enfant).
\end{attention}

\courssub{Nature grammaticale des mots : un mot, plusieurs fonctions}

\begin{propriete}[Un mot chinois peut changer de nature]
Contrairement au français, beaucoup de mots chinois fonctionnent comme \textbf{verbe} ou comme \textbf{nom} sans changer de forme :
\begin{itemize}
\item \zh{传染} : \textbf{verbe} dans \zh{这种病传染人。} \emph{Zhè zhǒng bìng chuánrǎn rén.} « Cette maladie infecte les gens. » ; \textbf{nom} dans \zh{防止传染} \emph{fángzhǐ chuánrǎn} « empêcher la contagion ».
\item \zh{免疫} : \textbf{nom} dans \zh{免疫系统} \emph{miǎnyì xìtǒng} « le système immunitaire » ; \textbf{verbe} dans le sens d'« immuniser ».
\item \zh{检查} : \textbf{verbe} \zh{检查身体} « examiner le corps » ; \textbf{nom} \zh{做检查} « passer un examen (médical) ».
\end{itemize}
C'est le \textbf{contexte} et la \textbf{place dans la phrase} qui décident de la nature du mot.
\end{propriete}

\begin{attention}[Piège francophone : contaminer, polluer, infecter]
En français, « contaminer » peut évoquer une pollution chimique. En chinois, \zh{污染} (\emph{wūrǎn} « polluer ») s'emploie pour l'\textbf{environnement} (\zh{污染空气} « polluer l'air »), \textbf{jamais pour une personne malade}. Pour une maladie qui passe d'une personne à une autre, utilisez \zh{传染} ; pour dire « être infecté », utilisez \zh{感染} : \zh{他感染了病毒。} \emph{Tā gǎnrǎn le bìngdú.} « Il a été infecté par le virus. »
\end{attention}

\courssub{Les collocations (搭配) : les mots qui vont ensemble}

\begin{definition}[Collocation (搭配 \emph{dāpèi})]
Une \cle{collocation} est une association naturelle et habituelle de mots. En chinois comme en français, on ne peut pas combiner n'importe quel verbe avec n'importe quel nom. On dit \zh{治疗疾病} (\emph{zhìliáo jíbìng} « traiter une maladie ») ou \zh{治病} / \zh{看病}, mais jamais *\zh{医疾病}. Le verbe \zh{医} (\emph{yī}) existe, mais il est littéraire et ne s'emploie pas seul dans la langue courante.
\end{definition}

\begin{exemplebox}[Collocations modèles de la leçon]
\zh{患有艾滋病} \emph{huànyǒu àizībìng} « être atteint du sida »\\[2pt]
\zh{免疫系统受损} \emph{miǎnyì xìtǒng shòu sǔn} « le système immunitaire est affaibli »\\[2pt]
\zh{避免接触血液} \emph{bìmiǎn jiēchù xuèyè} « éviter le contact avec le sang »\\[2pt]
\zh{定期体检} \emph{dìngqī tǐjiǎn} « faire régulièrement un bilan de santé »
\end{exemplebox}

\begin{methode}[Apprendre le vocabulaire efficacement]
\begin{enumerate}
\item \textbf{Par familles de radicaux} : tous les caractères contenant la clé \zh{疒} (\emph{nè}, la « clé de la maladie ») sont liés à la santé : \zh{病} (maladie), \zh{症} (symptôme), \zh{疼} (douleur), \zh{疫} (épidémie), \zh{疗} (soigner). Un seul radical = toute une famille de sens (détail dans la section 3).
\item \textbf{Par « chunks » (blocs) } : mémorisez des blocs entiers prêts à l'emploi : \zh{高危行为} (comportement à haut risque), \zh{传播途径} (voie de transmission), \zh{定期体检} (bilan régulier).
\item \textbf{Par paires contraires} : \zh{预防}/\zh{治疗}, \zh{严重}/\zh{轻微} s'apprennent ensemble.
\item \textbf{Par la prononciation} : lisez chaque mot à voix haute en respectant les tons, trois fois de suite.
\end{enumerate}
\end{methode}

\courssub{Carte mentale du vocabulaire}

\begin{popfigure}
\begin{center}
\begin{tikzpicture}[mindmap, grow cyclic,
  every node/.style={concept, align=center, font=\small},
  root concept/.append style={concept color=popblue, font=\large},
  level 1/.append style={level distance=3.4cm, sibling angle=90},
  level 2/.append style={level distance=2.4cm, sibling angle=45, font=\footnotesize}]
\node[root concept, align=center]{疾病\\ jíbìng}
  child[concept color=poporange]{ node[align=center]{症状\\ zhèngzhuàng}
    child { node[align=center]{发烧\\ fāshāo} }
    child { node[align=center]{咳嗽\\ késou} }
    child { node[align=center]{头疼\\ tóuténg} }
  }
  child[concept color=popgreen]{ node[align=center]{传播\\ chuánbō}
    child { node[align=center]{传染\\ chuánrǎn} }
    child { node[align=center]{血液\\ xuèyè} }
    child { node[align=center]{接触\\ jiēchù} }
  }
  child[concept color=poppurple]{ node[align=center]{预防\\ yùfáng}
    child { node[align=center]{避免\\ bìmiǎn} }
    child { node[align=center]{安全\\ ānquán} }
    child { node[align=center]{体检\\ tǐjiǎn} }
  }
  child[concept color=poppink]{ node[align=center]{治疗\\ zhìliáo}
    child { node[align=center]{看病\\ kànbìng} }
    child { node[align=center]{保护\\ bǎohù} }
    child { node[align=center]{免疫\\ miǎnyì} }
  };
\end{tikzpicture}
\end{center}
\legende{Carte mentale du vocabulaire de la leçon : quatre branches autour de 疾病.}
\end{popfigure}

\courssub{Classification par radical : préparer la mémorisation}

\begin{center}
\begin{tabularx}{\linewidth}{|X|X|X|}
\hline
\rowcolor{popblueL} Radical & Sens du radical & Caractères de la leçon \\
\hline
疒 (nè) & maladie & 病, 症, 疼, 疫, 疗 \\
\hline
氵(shuǐ, eau) & liquide & 液 (liquide), 治 (soigner) \\
\hline
血 (xuè) & sang & 血 (sang) \\
\hline
宀 / 心 & maison / cœur & 安 (paix, sécurité), 患 (souffrir de) \\
\hline
\end{tabularx}
\end{center}

Retenez l'idée maîtresse : \textbf{le radical donne le sens général, l'autre composant donne souvent le son}. Ainsi \zh{病} = \zh{疒} (maladie) + \zh{丙} (\emph{bǐng}, indice phonétique). Nous étudierons l'ordre des traits de ces caractères dans la section 3.

\courssub{Exercices de reconnaissance rapide}

\begin{exoresolu}[Caractère vers pinyin et sens]
\textbf{Énoncé.} Donnez le pinyin (avec tons) et la traduction de : 病毒, 预防, 严重, 接触.
\tcblower
\textbf{Solution.}
\begin{itemize}
\item \zh{病毒} : \emph{bìngdú} — « virus » (4\up{e} ton + 2\up{e} ton).
\item \zh{预防} : \emph{yùfáng} — « prévenir » (4\up{e} ton + 2\up{e} ton).
\item \zh{严重} : \emph{yánzhòng} — « grave » (2\up{e} ton + 4\up{e} ton).
\item \zh{接触} : \emph{jiēchù} — « contact ; toucher » (1\up{er} ton + 4\up{e} ton).
\end{itemize}
Piège : ne confondez pas \zh{预防} (\emph{yùfáng}, prévenir) et \zh{治疗} (\emph{zhìliáo}, traiter) : on prévient \textbf{avant}, on traite \textbf{après}.
\end{exoresolu}

\begin{exercice}[Français vers caractères]
Traduisez en chinois (caractères + pinyin) : 1. le système immunitaire ; 2. éviter le contact avec le sang ; 3. faire un bilan de santé régulier ; 4. une maladie contagieuse ; 5. être atteint du sida.
\end{exercice}

\begin{corrige}[Exercice « Français vers caractères »]
1. \zh{免疫系统} \emph{miǎnyì xìtǒng}. 2. \zh{避免接触血液} \emph{bìmiǎn jiēchù xuèyè}. 3. \zh{定期体检} \emph{dìngqī tǐjiǎn}. 4. \zh{传染病} \emph{chuánrǎnbìng}. 5. \zh{患有艾滋病} \emph{huànyǒu àizībìng}.
\end{corrige}

\begin{exercice}[Classement par radical]
Classez ces caractères selon leur radical : 病, 液, 症, 治, 疼, 血, 疫. Que vous apprend ce classement sur leur sens ?
\end{exercice}

\begin{corrige}[Exercice « Classement par radical »]
Radical \zh{疒} (maladie) : \zh{病, 症, 疼, 疫}. Radical \zh{氵} (eau/liquide) : \zh{液, 治}. Radical autonome : \zh{血} (sang). On constate que la clé \zh{疒} regroupe presque tous les mots liés à la maladie : c'est un excellent repère de lecture et de mémorisation.
\end{corrige}

\begin{savaistu}[爱滋 ou 艾滋 ?]
Le mot « sida » (AIDS) a été transcrit phonétiquement en chinois. À Taïwan et à Hong Kong, on écrivait \zh{爱滋} (\emph{àizī}), avec le caractère \zh{爱} « amour ». En Chine continentale, on a choisi \zh{艾滋} : le caractère \zh{艾} désigne une herbe médicinale (l'armoise, utilisée en moxibustion), ce qui donne au mot une couleur plus médicale et moins ambiguë. Les deux formes se prononcent pareil, mais dans vos examens, écrivez toujours \zh{艾滋病}.
\end{savaistu}

\begin{aretenir}[L'essentiel du glossaire]
\begin{itemize}
\item Quatre familles : symptômes (\zh{发烧, 咳嗽}), sida/VIH (\zh{艾滋病, 病毒, 免疫}), transmission/prévention (\zh{传染, 预防, 避免}), action médicale (\zh{治疗, 看病, 体检}).
\item Un mot = souvent deux caractères : apprenez \zh{病毒, 免疫, 系统, 接触} en bloc.
\item Les paires contraires se mémorisent ensemble : \zh{健康}/\zh{疾病}, \zh{传染}/\zh{遗传}, \zh{预防}/\zh{治疗}, \zh{严重}/\zh{轻微}.
\item La clé \zh{疒} signale la maladie : c'est votre meilleure amie pour lire et retenir.
\end{itemize}
\end{aretenir}

\coursec{Clés (Radicaux) et Ordre des traits : Les briques du corps médical}

\courssub{Règles générales d'ordre des traits (笔顺 \emph{bǐshùn})}

\begin{propriete}[Principes fondamentaux de l'écriture]
\begin{enumerate}
\item \textbf{Haut $\to$ Bas} (\zh{上到下}) : \zh{三} (sān), \zh{病} (partie haute \zh{疒}).
\item \textbf{Gauche $\to$ Droite} (\zh{左到右}) : \zh{好} (hǎo), \zh{症} (\zh{疒} + \zh{正}).
\item \textbf{Horizontal avant Vertical} (\zh{横竖}) : \zh{十} (shí), \zh{血} (premier trait horizontal).
\item \textbf{Extérieur avant Intérieur} (\zh{外到内}) : \zh{同} (tóng), \zh{疗} (cadre \zh{疒} avant \zh{了}).
\item \textbf{Fermer en dernier} (\zh{最后封口}) : \zh{国} (guó), \zh{囗} se ferme après le contenu.
\item \textbf{Trait central avant les ailes} (\zh{中间两边}) : \zh{小} (xiǎo), \zh{水} (shuǐ).
\end{enumerate}
\end{propriete}

\courssub{Analyse détaillée de 8 caractères clés}

\begin{center}
\begin{tabularx}{\linewidth}{|X|X|X|X|X|}
\hline
\rowcolor{popblueL} Caractère & Pinyin & Radical (Clé) & Nombre de traits & Décomposition / Ordre des traits principal \\
\hline
\zh{病} & bìng & \zh{疒} (nè, maladie) & 10 & \zh{疒} (5 traits : 丶, ㇆, 丿, 丶, 一) + \zh{丙} (5 traits : 一, 丨, 一, 丿, ㇏) \\
\hline
\zh{症} & zhèng & \zh{疒} (nè) & 13 & \zh{疒} (5) + \zh{正} (5 : 一, 丨, 一, 丨, 一) = 10 + 3 traits restants ? Non. \zh{正} 5 traits. Total 10. Attendez : \zh{症} = \zh{疒} (5) + \zh{正} (5) = 10 traits. (Correction : \zh{症} a 10 traits). Ordre : clé \zh{疒} puis \zh{正} (haut \zh{一}, verticale \zh{丨}, horizontales). \\
\hline
\zh{疫} & yì & \zh{疒} (nè) & 9 & \zh{疒} (5) + \zh{殳} (shū, 4 traits : ㇆, 丿, 丶, ㇏). \\
\hline
\zh{疗} & liáo & \zh{疒} (nè) & 8 & \zh{疒} (5) + \zh{了} (liǎo, 2 traits : ㇆, ㇏) + 1 trait ? Non, \zh{疒} (5) + \zh{了} (2) = 7. Vérification : \zh{疗} 8 traits. \zh{疒} (5) + \zh{了} (2) + point final ? Standard : \zh{疒} puis \zh{了}. \\
\hline
\zh{血} & xuè & \zh{血} (xuè, sang) & 6 & Caractère unique (clé). Ordre : \zh{丶} (point haut), \zh{一} (horiz.), \zh{丿} (piě), \zh{㇏} (nà), \zh{一} (horiz. bas), \zh{丨} (verticale centrale traversant). \\
\hline
\zh{患} & huàn & \zh{忄} (xīn, cœur vertical) & 11 & \zh{忄} (3 traits : 丶, 丶, 丿) + \zh{串} (chuàn, 8 traits : \zh{丨}, \zh{(trait brisé)}, \zh{一}, \zh{一}, \zh{丨}, \zh{(trait brisé)}, \zh{一}, \zh{一}). \\
\hline
\zh{检} & jiǎn & \zh{木} (mù, bois) & 12 & \zh{木} (4) + \zh{柬} (jiǎn, 8 traits : \zh{竹} simplifié \zh{⺮} 6 + \zh{㐱} ?). Simplifié : \zh{木} + \zh{柬} (partie droite). Ordre : \zh{木} puis partie droite (haut \zh{⺮}, bas \zh{㐱}). \\
\hline
\zh{治} & zhì & \zh{氵} (shuǐ, eau) & 8 & \zh{氵} (3 traits : 丶, 丶, 丿) + \zh{台} (tái, 5 traits : \zh{一}, \zh{丨}, \zh{一}, \zh{丨}, \zh{一}). \\
\hline
\end{tabularx}
\end{center}

\begin{exercice}[Écriture guidée]
Sur votre cahier d'exercices, écrivez chaque caractère ci-dessus \textbf{5 fois} en respectant scrupuleusement l'ordre des traits. Pour chaque caractère, soulignez le radical au crayon rouge et notez sa signification à côté.
\end{exercice}

\coursec{Compréhension de texte : Questions graduelles}

Relisez le texte \zh{认识疾病，预防艾滋} (Section 1) pour répondre aux questions suivantes. Répondez en français, sauf indication contraire.

\begin{exercice}[Questions de compréhension]
\begin{enumerate}
\item \textbf{Global :} Quel est le thème principal du texte ? Citez deux maladies « courantes » et une maladie « grave » mentionnées.
\item \textbf{Détail (Vocabulaire) :} Retrouvez dans le texte les mots chinois correspondant à : « système immunitaire », « transmission mère-enfant », « comportement à haut risque », « discriminer », « dépistage / examen ».
\item \textbf{Détail (Grammaire/Structure) :} La phrase \zh{病毒进入人体后，攻击免疫系统} contient le mot \zh{后}. Quelle structure grammaticale cela forme-t-il ? Traduisez la phrase.
\item \textbf{Inférence :} D'après le texte, pourquoi les piqûres de moustiques ne transmettent-elles pas le VIH ? (Réponse courte attendue).
\item \textbf{Expression personnelle (Chinois) :} Répondez en chinois (caractères + pinyin) : \zh{我们应该怎么对待艾滋病患者？} (Comment devons-nous traiter les malades du sida ?)
\end{enumerate}
\end{exercice}

\begin{corrige}[Exercice « Compréhension de texte »]
\begin{enumerate}
\item Thème : Connaissance des maladies et prévention du sida. Courantes : 感冒 (rhume), 发烧 (fièvre), 咳嗽 (toux), 头疼 (mal de tête). Grave : 艾滋病 (sida).
\item \zh{免疫系统} (miǎnyì xìtǒng), \zh{母婴传播} (mǔyīng chuánbō), \zh{高危行为} (gāowēi xíngwéi), \zh{歧视} (qīshì), \zh{检查} (jiǎnchá).
\item Structure : \textbf{Verbe/Action + 后} (hòu) = « Après avoir... / Une fois que... ». Traduction : « Après être entré dans le corps humain, le virus attaque le système immunitaire. »
\item Le virus ne survit pas / ne se réplique pas dans le moustique ; ce n'est pas un vecteur biologique. Le texte dit explicitement : « 日常接触... 蚊虫叮咬，不会传染 ».
\item \zh{我们不应该歧视他们，应该给予关爱和支持。} \emph{Wǒmen bù yīnggāi qīshì tāmen, yīnggāi gěiyǔ guānaì hé zhīchí.}
\end{enumerate}
\end{corrige}

\coursec{Collocations, Structures utiles et Production guidée}

\courssub{Structures grammaticales du thème (Niveau HSK 3)}

\begin{propriete}[La comparaison avec 比 (\emph{bǐ})]
\textbf{Structure :} Sujet A + \zh{比} + Sujet B + Adjectif (SV) \textbf{(sans \zh{很})}.
\textbf{Emploi :} Comparer deux entités sur une qualité. La négation : A \zh{不比} B Adj (A n'est pas plus... que B) ou A \zh{没比} B Adj (rare).
\end{propriete}

\begin{exemplebox}[Exemples du texte et nouveaux]
\zh{预防比治疗更重要。} \emph{Yùfáng bǐ zhìliáo gèng zhòngyào.} (Prévenir est plus important que traiter.)\\[2pt]
\zh{艾滋病比感冒严重得多。} \emph{Àizībìng bǐ gǎnmào yánzhòng de duō.} (Le sida est bien plus grave que le rhume.)\\[2pt]
\zh{我的病没比他的病严重。} \emph{Wǒ de bìng méi bǐ tā de bìng yánzhòng.} (Ma maladie n'est pas plus grave que la sienne.)
\end{exemplebox}


\begin{propriete}[Le changement progressif : 越来越 (\emph{yuè lái yuè})]
\textbf{Structure :} Sujet + \zh{越来越} + Adjectif / Verbe d'état mental.
\textbf{Emploi :} Indique une intensification progressive dans le temps.
\end{propriete}

\begin{exemplebox}[Exemples]
\zh{他的病情越来越严重。} \emph{Tā de bìngqíng yuè lái yuè yánzhòng.} (Son état s'aggrave de plus en plus.)\\[2pt]
\zh{大家对艾滋病的了解越来越多。} \emph{Dàjiā duì Àizībìng de liǎojiě yuè lái yuè duō.} (Tout le monde en sait de plus en plus sur le sida.)
\end{exemplebox}


\begin{propriete}[L'inclusion / exception : 除了... 以外 / 还 (\emph{chúle... yǐwài / hái})]
\textbf{Structure :} \zh{除了} A \zh{以外} (, 还/也) B. « En plus de A, (aussi) B. »
\textbf{Emploi :} Énumérer des voies de transmission, des mesures de prévention.
\end{propriete}

\begin{exemplebox}[Exemples]
\zh{除了血液传播，还有性传播和母婴传播。} \emph{Chúle xuèyè chuánbō, hái yǒu xìng chuánbō hé mǔyīng chuánbō.} (En plus de la transmission sanguine, il y a la transmission sexuelle et mère-enfant.)\\[2pt]
\zh{除了避免接触血液，还要正确使用安全套。} \emph{Chúle bìmiǎn jiēchù xuèyè, hái yào zhèngquè shǐyòng ānquántào.}
\end{exemplebox}


\begin{propriete}[Aspect accompli 了 (\emph{le}) vs Expérience 过 (\emph{guò})]
\textbf{了 (le) :} Action terminée, changement d'état, événement nouveau. \zh{他得了感冒。} (Il a attrapé un rhume — nouvel état). \zh{我做了检查。} (J'ai fait l'examen — action achevée).
\textbf{过 (guò) :} Expérience vécue au moins une fois dans le passé. \zh{我做过艾滋病检查。} (J'ai déjà fait un test du sida — expérience de vie).
\end{propriete}

\begin{attention}[Piège] Ne dites pas *\zh{我做了过检查}. Choisissez \textbf{soit} \zh{了} (focus sur la fin de l'action récente) \textbf{soit} \zh{过} (focus sur l'expérience passée).
\end{attention}


\courssub{Production guidée : Écriture et Oral}

\begin{methode}[Rédiger un message de prévention (Affiche / SMS)]
\textbf{Contexte :} La Journée mondiale de lutte contre le sida (1er décembre). Vous êtes délégué santé au lycée de Yaoundé / Douala / Buea.
\textbf{Étapes :}
\begin{enumerate}
\item \textbf{Accroche} : Titre court, percutant. Ex: \zh{预防艾滋，从我做起} (Prévenir le sida, commence par moi).
\item \textbf{Faits clés (3 points)} : Voies de transmission (\zh{三条途径}), Moyens de prévention (\zh{安全套, 不共用针头}), Dépistage (\zh{早检查, 早治疗}).
\item \textbf{Appel à l'action / Lutte contre la stigmatisation} : \zh{不歧视, 给关爱} (Pas de discrimination, donnons de l'amour).
\item \textbf{Contact} : \zh{咨询电话: 12345} / \zh{校医室}.
\end{enumerate}
\textbf{Modèle :}
\zh{【预防艾滋，守护健康】} \\
\zh{同学们：艾滋病通过血液、性、母婴传播。日常接触不传染。} \\
\zh{请记住：不吸毒、不共用针头、性行为用安全套。} \\
\zh{定期体检，早发现早治疗。不歧视患者，给予关爱。} \\
\zh{校医室咨询：XXXX-XXXX} \\
\emph{Pinyin et traduction fournis en annexe.}
\end{methode}

\begin{experience}[Jeu de rôle : « Chez le médecin » / « Au centre de dépistage »]
\textbf{Rôles :} A = Patient (élève) / B = Médecin / Infirmier(ère).
\textbf{Scénario 1 (Maladie banale) :} L'élève a \zh{发烧, 咳嗽, 头疼} depuis 2 jours. Le médecin demande les symptômes (\zh{什么症状？}), prescrit (\zh{开药}), conseille repos (\zh{多休息, 多喝水}).
\textbf{Scénario 2 (Dépistage VIH) :} L'élève veut faire \zh{艾滋病检查}. L'infirmier explique la procédure (\zh{抽血, 等结果}), la confidentialité (\zh{保密}), le sens du résultat (\zh{阴性/阳性}), et la fenêtre sérologique (\zh{窗口期} — vocabulaire bonus).
\textbf{Contraintes linguistiques :} Utiliser au moins 5 mots du glossaire + structure \zh{得了...} / \zh{要...} / \zh{应该...}.
\end{experience}

\begin{exercice}[Production écrite : Lettre ouverte / Petit essai]
Rédigez un court texte en chinois (80--100 caractères) sur le thème : \zh{《健康生活，远离艾滋》} (\emph{« Vie saine, loin du sida »}).
\textbf{Points à inclure :}
\begin{itemize}
\item L'importance de la santé (\zh{健康很重要}).
\item Connaissance des voies de transmission (\zh{知道传播途径}).
\item Mesures de prévention concrètes (\zh{采取预防措施} : \zh{安全套, 不共用针头}).
\item Attitude envers les malades (\zh{不歧视, 给关爱}).
\end{itemize}
\textbf{Aide :} Utilisez les connecteurs \zh{首先} (d'abord), \zh{其次} (ensuite), \zh{最后} (enfin), \zh{因为}... \zh{所以}... (parce que... donc...).
\end{exercice}

\begin{corrige}[Modèle de rédaction — Niveau visé]
\zh{健康生活，远离艾滋} \\[2pt]
\zh{健康是幸福的基础。首先，我们要认识艾滋病，知道它通过血液、性行为和母婴传播，日常接触不传染。其次，要采取预防措施：不吸毒，不共用针头，性行为正确使用安全套。定期体检，早发现早治疗。最后，我们不应该歧视艾滋病患者，而要给予他们关爱和支持。预防艾滋，人人有责，让我们一起守护健康！} \\[4pt]
\emph{Jiànkāng shēnghuó, yuǎnlí Àizī.} \\
\emph{Jiànkāng shì xìngfú de jīchǔ. Shǒuxiān, wǒmen yào rènshi Àizībìng, zhīdào tā tōngguò xuèyè, xìng xíngwéi hé mǔyīng chuánbō, rìcháng jiēchù bù chuánrǎn. Qícì, yào cǎiqǔ yùfáng cuòshī: bù xīdú, bù gòngyòng zhēntóu, xìng xíngwéi zhèngquè shǐyòng ānquántào. Dìngqī tǐjiǎn, zǎo fāxiàn zǎo zhìliáo. Zuìhòu, wǒmen bù yīnggāi qīshì Àizībìng huànzhě, ér yào gěiyǔ tāmen guānaì hé zhīchí. Yùfáng Àizī, rénrén yǒu zé, ràng wǒmen yīqǐ shǒuhù jiànkāng!}
\end{corrige}

\coursec{Approfondissement socioculturel : Cameroun - Chine, regards croisés}

\courssub{Le SIDA au Cameroun : données et actions}

\begin{savaistu}[Situation épidémiologique et riposte au Cameroun]
Selon les rapports d'ONUSIDA et du Ministère de la Santé Publique (MINSANTÉ), le Cameroun a fait des progrès notables dans la riposte au VIH/Sida. La prévalence chez les adultes (15-49 ans) est passée d'environ 5,4\% (2004) à environ 2,7\% (estimations récentes), grâce à la gratuité des ARV (antirétroviraux) décrétée en 2007 et au déploiement des services de \zh{防止母婴传播} (PMTCT / \zh{阻断母婴传播}). Cependant, des défis persistent : taux de dépistage encore insuffisant chez les jeunes, stigmatisation des \zh{重点人群} (populations clés : travailleurs du sexe, HSH, usagers de drogues), et disparités régionales (Extrême-Nord, Adamaoua plus touchés). Le \zh{Comité National de Lutte contre le Sida} (CNLS) coordonne la stratégie nationale alignée sur les objectifs « 95-95-95 » de l'ONU pour 2030.
\end{savaistu}

\courssub{La Chine : politique « Quatre gratuités et une prise en charge » (四免一关怀)}

\begin{savaistu}[Politique chinoise de lutte contre le SIDA]
Depuis 2003, la Chine met en œuvre la politique des \zh{四免一关怀} (\emph{Sì miǎn yī guānhuái}) :
\begin{enumerate}
\item \zh{免费抗病毒治疗} (Traitement antirétroviral gratuit pour les patients pauvres/ruraux).
\item \zh{免费自愿咨询检测} (Conseil et dépistage volontaires gratuits).
\item \zh{免费阻断母婴传播} (Prévention mère-enfant gratuite : dépistage, ARV, substituts de lait maternel, accouchement sécurisé).
\item \zh{免费救助艾滋病孤儿} (Prise en charge gratuite des orphelins du sida : scolarité, vie quotidienne).
\item \zh{关怀救助生活困难的艾滋病病人} (Aide sociale aux patients en difficulté : allocation minimale, assurance maladie).
\end{enumerate}
Résultat : La transmission par transfusion sanguine est quasi éliminée ; la transmission mère-enfant chutée drastiquement. Le défi actuel : la transmission sexuelle (hétéro et homo) devient la voie majoritaire (> 90\%), et le dépistage tardif reste fréquent.
\end{savaistu}

\courssub{Comparatif : Regards croisés Cameroun -- Chine}

\begin{center}
\begin{tabularx}{\linewidth}{|X|X|X|}
\hline
\rowcolor{popblueL} Dimension & Cameroun & Chine \\
\hline
\textbf{Prévalence (adultes)} & \~ 2,7\% (généralisée) & \textless 0,1\% (faible, concentrée) \\
\hline
\textbf{Voie principale} & Hétérosexuelle (> 90\%) & Sexuelle (hétéro + homo), ex-hétéro dominant \\
\hline
\textbf{Accès ARV} & Gratuit depuis 2007, couverture \~ 80\% & Gratuit (politique 4 gratuités), couverture élevée \\
\hline
\textbf{PMTCT (Mère-enfant)} & Programme national, défis logistiques (zones rurales) & Intégré aux soins maternels, taux transmission \textless 4\% \\
\hline
\textbf{Stigmatisation} & Forte, frein au dépistage & Forte historiquement, en baisse grâce aux campagnes médiatiques \\
\hline
\textbf{Médecine traditionnelle} & \zh{中医} (Médecine chinoise) vs \zh{非洲传统医学} : rôle complémentaire reconnu, mais \textbf{ne remplacent pas les ARV}. & Intégration \zh{中西医结合} (Médecine chinoise + occidentale) officielle pour effets secondaires ARV / immunité. \\
\hline
\end{tabularx}
\end{center}

\begin{experience}[Débat guidé : « Stigmatisation et droits des malades »]
\textbf{Sujet :} \zh{艾滋病患者应该享有和普通人一样的就业权和受教育权吗？} (Les malades du sida doivent-ils avoir les mêmes droits à l'emploi et à l'éducation que les autres ?)
\textbf{Modalités :} Deux groupes (Pour / Contre -- jouer le rôle). Arguments attendus : Droits humains (Charte ONU, Constitution), Non-contagiosité au travail/école (pas de risque par contact quotidien), Coût social de l'exclusion, Peur irrationnelle vs faits scientifiques.
\textbf{Vocabulaire débat :} \zh{平等} (égalité), \zh{歧视} (discrimination), \zh{隐私} (vie privée), \zh{就业} (emploi), \zh{入学} (scolarisation), \zh{法律保护} (protection légale).
\end{experience}

\begin{aretenir}[Synthèse socioculturelle]
\begin{itemize}
\item \textbf{Point commun :} Les deux pays ont rendu les ARV gratuits et ciblent l'élimination de la transmission mère-enfant.
\item \textbf{Différence d'épidémie :} Cameroun = épidémie généralisée (population générale) ; Chine = épidémie concentrée (populations clés), mais en expansion sexuelle.
\item \textbf{Leçon commune :} La \cle{connaissance scientifique} (connaître les \zh{传播途径}) et la \cle{lutte contre la stigmatisation} (\zh{反歧视}) sont les deux piliers pour atteindre l'objectif « Zéro nouvelle infection, Zéro discrimination, Zéro décès lié au sida ».
\end{itemize}
\end{aretenir}

\coursec{Clés (Radicaux) et Ordre des traits : Les briques du corps médical}

Dans la section précédente, nous avons découvert le vocabulaire maître du thème « maladies et SIDA » : \zh{疾病}, \zh{艾滋病}, \zh{免疫}, \zh{检查}, \zh{歧视}… Vous avez peut-être remarqué que plusieurs de ces caractères se ressemblent : \zh{病}, \zh{疾}, \zh{症}, \zh{痛} partagent tous un même élément en haut à gauche ; \zh{肝}, \zh{肺}, \zh{脏}, \zh{胃} partagent un élément à gauche ou en bas. Ce n'est pas un hasard : les caractères chinois sont construits comme des briques Lego, et ces briques récurrentes s'appellent les \cle{clés} ou \cle{radicaux} (\zh{部首} \emph{bùshǒu}). Comprendre les radicaux, c'est pouvoir deviner le sens d'un caractère inconnu, le chercher dans un dictionnaire et le mémoriser durablement. C'est l'objet de cette section.

\courssub{Qu'est-ce qu'une clé de dictionnaire ?}

\begin{definition}[La clé de dictionnaire \zh{部首}]
Une \cle{clé} (\zh{部首} \emph{bùshǒu}, « tête de section ») est l'élément graphique sous lequel un caractère est classé dans un dictionnaire chinois. Elle donne généralement une indication de \textbf{sens} (le champ sémantique), tandis que l'autre partie du caractère donne souvent une indication de \textbf{prononciation} (la phonétique).\\
Exemples : pour \zh{病}, \zh{疾}, \zh{症}, la clé est \zh{疒} (\emph{bìng zì tóu}, « tête du caractère maladie »). Pour \zh{肝}, \zh{肺}, la clé est \zh{月} (\emph{ròu zì páng}, « côté du caractère chair »).
\end{definition}

\begin{methode}[Chercher un caractère inconnu dans un dictionnaire]
\begin{enumerate}
\item Identifier le radical (souvent à gauche, en haut, ou en enveloppe).
\item Compter les traits du radical, puis trouver le radical dans l'index des clés.
\item Compter les traits \emph{restants} du caractère (hors radical).
\item Retrouver le caractère dans la liste des caractères ayant ce nombre de traits sous cette clé.
\end{enumerate}
Exemple : \zh{瘦} → radical \zh{疒} (5 traits) + 9 traits restants → on le trouve dans la section \zh{疒} à « +9 traits ».
\end{methode}

\courssub{Le radical \zh{疒} : la clé de la maladie}

Observez d'abord ces caractères du texte support : \zh{病} (maladie), \zh{疾} (maladie), \zh{症} (symptôme), \zh{痛} (douleur), \zh{疲} (fatigué), \zh{瘦} (maigre). Tous contiennent \zh{疒}, et tous évoquent le corps souffrant. Le radical agit comme une étiquette de sens.

\begin{propriete}[Le radical \zh{疒} (\emph{nè})]
\begin{itemize}
\item \textbf{Sens} : « maladie, état maladif ». Il n'existe presque jamais seul ; il sert de composant.
\item \textbf{Position} : toujours en \textbf{enveloppe haut-gauche} : il recouvre le haut et le côté gauche du caractère, comme un toit de lit d'hôpital. On l'appelle familièrement \zh{病字头} (\emph{bìng zì tóu}).
\item \textbf{Nombre de traits} : 5 (点 → 横 → 竖 → 提 → 点).
\item \textbf{Variante} : aucune variante graphique ; il ne se confond pas avec \zh{广} (« vaste », radical de \zh{店}), qui n'a que 3 traits et n'a pas les deux petits traits à gauche.
\end{itemize}
\end{propriete}

\begin{savaistu}[D'où vient \zh{疒} ?]
Le radical \zh{疒} vient d'une ancienne image : un homme \textbf{alité} sur un lit (\zh{爿}, le lit vu de profil), avec de petits points représentant la \textbf{sueur} ou la fièvre du malade. Quand vous écrivez \zh{疒}, imaginez un malade couché qui transpire : vous n'oublierez jamais ni sa forme ni son sens ! Cette étymologie pictographique est typique de l'écriture chinoise, née de dessins il y a plus de trois mille ans.
\end{savaistu}

\begin{center}
\begin{tabularx}{\linewidth}{|X|X|X|X|}
\hline
\rowcolor{popblueL} Caractère & Pinyin & Décomposition & Traduction \\
\hline
病 & bìng & 疒 + 丙 (phonétique bǐng) & maladie ; tomber malade \\
\hline
疾 & jí & 疒 + 矢 (flèche) & maladie (rapide comme une flèche) \\
\hline
症 & zhèng & 疒 + 正 (phonétique zhèng) & symptôme \\
\hline
痛 & tòng & 疒 + 甬 (phonétique yǒng) & douleur ; douloureux \\
\hline
疲 & pí & 疒 + 皮 (phonétique pí) & fatigué \\
\hline
瘦 & shòu & 疒 + 叟 (phonétique sǒu) & maigre \\
\hline
\end{tabularx}
\end{center}

Remarquez la logique : \zh{疲} (fatigué) et \zh{瘦} (maigre) sont classés sous « maladie » car, dans la conception ancienne, fatigue extrême et maigreur étaient des signes de mauvaise santé. La phonétique, elle, guide la prononciation : \zh{病} se lit \emph{bìng} comme sa phonétique \zh{丙} (\emph{bǐng}) — seul le ton change.

\courssub{Le radical \zh{月}/\zh{肉} : la clé de la chair et des organes}

Deuxième famille : \zh{肝} (foie), \zh{肺} (poumons), \zh{肾} (reins), \zh{脏} (organe), \zh{腑} (entrailles), \zh{胃} (estomac), \zh{肌} (muscle). Tous contiennent \zh{月}… mais attention, ce n'est \emph{pas} la lune !

\begin{propriete}[Le radical \zh{月} = \zh{肉} (\emph{ròu}, « chair »)]
\begin{itemize}
\item \textbf{Sens} : dans les caractères du corps humain, \zh{月} est la forme comprimée de \zh{肉} (\emph{ròu}, « viande, chair »). On l'appelle \zh{肉字旁} (\emph{ròu zì páng}, « côté chair »). Il indique un \textbf{organe, un muscle, une partie du corps}.
\item \textbf{Position} : le plus souvent à \textbf{gauche} (\zh{肝}, \zh{肺}, \zh{脏}, \zh{肌}), parfois en \textbf{bas} (\zh{胃}, \zh{肾}).
\item \textbf{Piège} : le vrai \zh{月} « lune/mois » existe aussi comme radical (\zh{明}, \zh{期}). Astuce : si le caractère désigne une partie du corps, c'est « chair » ; s'il désigne le temps ou la lumière, c'est « lune ».
\end{itemize}
\end{propriete}

\begin{center}
\begin{tabularx}{\linewidth}{|X|X|X|X|}
\hline
\rowcolor{popblueL} Caractère & Pinyin & Position de 月 & Traduction \\
\hline
肝 & gān & gauche & foie \\
\hline
肺 & fèi & gauche & poumons \\
\hline
肾 & shèn & bas & reins \\
\hline
脏 & zàng & gauche & organe interne \\
\hline
腑 & fǔ & gauche & entrailles, viscères \\
\hline
胃 & wèi & bas & estomac \\
\hline
肌 & jī & gauche & muscle \\
\hline
\end{tabularx}
\end{center}

\begin{attention}[Ne confondez pas \zh{月}-chair et \zh{月}-lune]
Les francophones écrivent souvent \zh{明} (clair : soleil + lune) et \zh{肝} (foie : chair + phonétique) avec le même réflexe mental. Or \zh{肝脏} (\emph{gānzàng}, le foie) n'a rien à voir avec la lune ! De même, ne confondez pas \zh{脏} (\emph{zàng}, organe) et \zh{脏} lu \emph{zāng} (sale) : même caractère, deux lectures — le contexte décide.
\end{attention}

\courssub{Le radical \zh{血} : la clé du sang}

Troisième brique, essentielle pour notre thème : \zh{血} (\emph{xuè} / \emph{xiě}, « sang »). C'est à la fois un \textbf{caractère autonome} et un \textbf{radical} (clé n°143). Sa forme ancienne dessine un récipient (\zh{皿}) contenant une goutte de sang offerte en sacrifice.

\begin{itemize}
\item \zh{血} seul : \zh{血液} \emph{xuèyè} « le sang (terme médical) » ; \zh{血型} \emph{xuèxíng} « groupe sanguin » ; \zh{输血} \emph{shūxuè} « transfusion sanguine ».
\item Comme composant, il est rare dans les caractères courants : \zh{衅} (\emph{xìn}, « querelle », caractère rare) ; on le retrouve surtout dans les mots composés médicaux.
\end{itemize}

\begin{exemplebox}[Le sang dans le vocabulaire médical]
\zh{输血可能传播艾滋病毒。}\\
\emph{Shūxuè kěnéng chuánbō àizībìngdú.}\\
« La transfusion sanguine peut transmettre le virus du sida. »\\[4pt]
\zh{你知道自己的血型吗？}\\
\emph{Nǐ zhīdào zìjǐ de xuèxíng ma?}\\
« Connais-tu ton groupe sanguin ? »
\end{exemplebox}

\courssub{L'ordre des traits : les six règles d'or}

Écrire un caractère chinois n'est pas dessiner librement : chaque trait s'écrit dans un \textbf{ordre fixe} (norme GB13000.1). Cet ordre garantit une écriture rapide, lisible et équilibrée — et il est exigé à l'examen.

\begin{propriete}[Les six règles d'ordre des traits]
\begin{enumerate}
\item \textbf{Horizontal avant vertical} : \zh{十} (\emph{shí}, dix) : d'abord 一, ensuite 丨.
\item \textbf{Gauche avant droite} : \zh{人} (\emph{rén}, personne) : le trait tombant à gauche d'abord.
\item \textbf{Haut avant bas} : \zh{三} (\emph{sān}, trois) : du trait du haut vers celui du bas.
\item \textbf{Extérieur avant intérieur} : \zh{同} (\emph{tóng}, ensemble) : l'enveloppe d'abord, le contenu ensuite, puis on « ferme la porte ».
\item \textbf{Milieu avant côtés} : \zh{小} (\emph{xiǎo}, petit) : le crochet central, puis les deux points.
\item \textbf{Trait traversant en dernier} : \zh{中} (\emph{zhōng}, milieu) : la verticale qui traverse s'écrit après la bouche \zh{口}.
\end{enumerate}
\end{propriete}

\begin{attention}[Erreur fréquente : le radical \zh{疒}]
Beaucoup d'élèves écrivent \zh{疒} en traçant le point final (le 5\textsuperscript{e} trait) \emph{avant} le long trait horizontal. C'est faux. L'ordre correct est : \textbf{点 (1) → 横 (2) → 竖 (3) → 提 (4) → 点 (5)}. Le long horizontal est le \textbf{2\textsuperscript{e} trait}, tout de suite après le point du sommet ; les deux petits traits de gauche (竖 puis 提, puis le point) viennent ensuite. Retenez : « le toit d'abord, les gouttes de sueur ensuite ».
\end{attention}

\courssub{Décomposition guidée : \zh{病} et \zh{免}}

\begin{popfigure}
\begin{tikzpicture}[scale=1.0, every node/.style={font=\small}]
% Grille 田字格 pour 病
\draw[popline, thick] (0,0) rectangle (3,3);
\draw[popline, dashed] (0,1.5)--(3,1.5);
\draw[popline, dashed] (1.5,0)--(1.5,3);
\draw[popline, dashed] (0,0)--(3,3);
\draw[popline, dashed] (0,3)--(3,0);
\node[popblue] at (1.5,1.5) {\Huge 病};
\node[below, popmuted] at (1.5,-0.1) {1. 点 en haut};
% Flèches ordre
\draw[-{Stealth}, poporange, very thick] (3.4,2.6) -- (4.6,2.6) node[midway, above, popdark]{疒};
\draw[popline, thick] (4.8,0) rectangle (7.8,3);
\draw[popline, dashed] (4.8,1.5)--(7.8,1.5);
\draw[popline, dashed] (6.3,0)--(6.3,3);
\node[popgreen] at (6.3,1.5) {\Huge 疒};
\node[below, popmuted] at (6.3,-0.1) {2. 疒 : 5 traits (点横竖提点)};
\draw[-{Stealth}, poporange, very thick] (8.2,2.6) -- (9.4,2.6) node[midway, above, popdark]{+ 丙};
\draw[popline, thick] (9.6,0) rectangle (12.6,3);
\draw[popline, dashed] (9.6,1.5)--(12.6,1.5);
\draw[popline, dashed] (11.1,0)--(11.1,3);
\node[poppurple] at (11.1,1.5) {\Huge 丙};
\node[below, popmuted] at (11.1,-0.1) {3. 丙 : 5 traits (phonétique bǐng)};
\end{tikzpicture}
\legende{Décomposition de \zh{病} (10 traits) : l'enveloppe \zh{疒} (radical, sens « maladie ») s'écrit d'abord, puis le contenu \zh{丙} (phonétique). Règle appliquée : extérieur avant intérieur.}
\end{popfigure}

\begin{popfigure}
\begin{tikzpicture}[scale=1.0, every node/.style={font=\small}]
\draw[popline, thick] (0,0) rectangle (3,3);
\draw[popline, dashed] (0,1.5)--(3,1.5);
\draw[popline, dashed] (1.5,0)--(1.5,3);
\node[popblue] at (1.5,1.5) {\Huge 免};
\node[below, popmuted, align=center] at (1.5,-0.1) {7 traits : 撇 → 横撇 → 竖 →\\ 横折 → 横 → 撇 → 竖弯钩};
\draw[-{Stealth}, poporange, very thick] (3.4,1.5) -- (4.6,1.5);
\node[align=left, popdark] at (7.6,1.5) {Le 5\textsuperscript{e} trait (横 long) ne touche\\ \textbf{pas} le bord droit : il s'arrête\\ avant la fermeture du « cadre ».};
\end{tikzpicture}
\legende{Écriture de \zh{免} (\emph{miǎn}, « exempter ») : attention au long horizontal qui reste ouvert à droite.}
\end{popfigure}

\begin{exemplebox}[Comptage des traits des caractères cibles]
\begin{itemize}
\item \zh{病} : \textbf{10 traits} (疒 = 5 + 丙 = 5).
\item \zh{疾} : \textbf{10 traits} (疒 = 5 + 矢 = 5).
\item \zh{症} : \textbf{10 traits} (疒 = 5 + 正 = 5 ; 正 s'écrit en 5 traits : 横、竖、横、竖、横).
\item \zh{艾} : \textbf{5 traits} (艹 = 3 + 乂 = 2).
\item \zh{滋} : \textbf{12 traits} (氵 = 3 + 兹 = 9).
\item \zh{免} : \textbf{7 traits} (撇、横撇、竖、横折、横、撇、竖弯钩) — attention au 5\textsuperscript{e} trait : le long horizontal \textbf{ne touche pas le bord droit}.
\item \zh{疫} : \textbf{9 traits} (疒 = 5 + 殳 = 4).
\item \zh{检} : \textbf{11 traits} (木 = 4 + 佥 = 7).
\item \zh{歧} : \textbf{8 traits} (止 = 4 + 支 = 4).
\item \zh{视} : \textbf{8 traits} (礻 = 4 + 见 = 4).
\end{itemize}
\end{exemplebox}

\begin{methode}[Mémorisation par l'étymologie et l'imagerie mnémotechnique]
\begin{enumerate}
\item \textbf{\zh{病}} = \zh{疒} (le lit du malade) + \zh{丙} (le son \emph{bǐng}) → « le malade alité qui gémit \emph{bing} ». Sens : maladie.
\item \textbf{\zh{免}} = \zh{儿} (enfant) + trait supérieur évoquant « sans/non » → image mnémotechnique : « un enfant sans souci » est \textbf{dispensé} de tout → sens : exempter, éviter. D'où \zh{免疫} (\emph{miǎnyì}) : « dispensé de l'épidémie » = immunité !
\item \textbf{\zh{疫}} = \zh{疒} (maladie) + \zh{殳} (arme, son \emph{shū}) → la maladie qui frappe comme une arme : l'épidémie.
\item \textbf{\zh{检}} = \zh{木} (bois) + \zh{佥} (son \emph{qiān}) → autrefois : marque de bois vérifiée → vérifier, examiner (\zh{检查} : examen médical).
\end{enumerate}
\end{methode}

\courssub{Écriture guidée sur quadrillage \zh{田字格}}

Le \zh{田字格} (\emph{tiánzìgé}, « grille-champ ») partage le carré en quatre cases par un axe horizontal et un axe vertical. Il oblige à respecter les \textbf{proportions} : dans \zh{病}, l'enveloppe \zh{疒} occupe le haut et la gauche, tandis que \zh{丙} se loge dans le quart inférieur droit, sans dépasser. Dans \zh{肝}, le radical \zh{月} à gauche est \emph{étroit} (un tiers de la largeur), la phonétique \zh{干} occupe les deux tiers.

\begin{exoresolu}[Écrire \zh{疫} dans le respect des règles]
Énoncé : décrivez l'ordre complet des traits de \zh{疫} (\emph{yì}, « épidémie ») et justifiez chaque étape par une règle d'or.
\tcblower
Solution détaillée : \zh{疫} = \zh{疒} + \zh{殳}, 9 traits.
\begin{enumerate}
\item Traits 1 à 5 : \zh{疒} — 点 (haut), 横 (long horizontal, 2\textsuperscript{e} trait !), 竖, 提, 点. Règle : \emph{extérieur avant intérieur} (l'enveloppe d'abord) et \emph{haut avant bas}.
\item Traits 6 à 9 : \zh{殳} — 撇, 横折弯 (la partie \zh{几} en haut), puis 横撇, 捺 (la partie \zh{又} en bas). Règle : \emph{haut avant bas} à l'intérieur du contenu, et \emph{gauche avant droite} pour le couple 撇/捺 final.
\end{enumerate}
Vérification : 5 + 4 = 9 traits. Le caractère reste dans le quadrillage : \zh{疒} couvre le haut-gauche, \zh{殳} s'inscrit en bas à droite.
\end{exoresolu}

\begin{exercice}[Jeu : « Quel radical manque-t-il ? »]
Complétez chaque caractère avec le bon radical (\zh{疒}, \zh{月}, \zh{氵}, \zh{礻}) :
\begin{enumerate}
\item \zh{丙} + ? = maladie
\item \zh{干} + ? = foie
\item \zh{兹} + ? = (dans \zh{艾滋})
\item \zh{见} + ? = regarder (\zh{视})
\end{enumerate}
Puis classez ces caractères sous le bon radical : \zh{痛}, \zh{肺}, \zh{瘦}, \zh{胃}, \zh{症}, \zh{肌}.
\end{exercice}

\begin{corrige}[Exercice « Quel radical manque-t-il ? »]
1. \zh{疒} + \zh{丙} = \zh{病}. \quad 2. \zh{月} + \zh{干} = \zh{肝}. \quad 3. \zh{氵} + \zh{兹} = \zh{滋}. \quad 4. \zh{礻} + \zh{见} = \zh{视}.\\
Classement : sous \zh{疒} : \zh{痛}, \zh{瘦}, \zh{症} ; sous \zh{月}(肉) : \zh{肺}, \zh{胃}, \zh{肌}.
\end{corrige}

\begin{aretenir}[L'essentiel de la section]
\begin{itemize}
\item \zh{疒} (5 traits, enveloppe haut-gauche) = famille de la \textbf{maladie} : \zh{病}, \zh{疾}, \zh{症}, \zh{痛}, \zh{疲}, \zh{瘦}, \zh{疫}.
\item \zh{月} = \zh{肉} (chair) = famille des \textbf{organes} : \zh{肝}, \zh{肺}, \zh{肾}, \zh{脏}, \zh{胃}, \zh{肌}. Ne pas confondre avec \zh{月} « lune ».
\item \zh{血} = le \textbf{sang} : \zh{血液}, \zh{血型}, \zh{输血}.
\item Six règles d'ordre : horizontal avant vertical ; gauche avant droite ; haut avant bas ; extérieur avant intérieur ; milieu avant côtés ; traversant en dernier.
\item Dans \zh{疒}, le long horizontal est le \textbf{2\textsuperscript{e}} trait, jamais le dernier.
\end{itemize}
\end{aretenir}

Vous savez désormais lire l'architecture interne des caractères du corps médical. Cette compétence vous servira immédiatement dans la section suivante, où nous vérifierons votre compréhension globale du texte \zh{认识疾病，预防艾滋} à l'aide de questions graduelles.

\coursec{Compréhension de texte : Questions graduelles}

Dans la section précédente, nous avons appris à décomposer les caractères du vocabulaire médical grâce aux clés (radicaux) et à l'ordre des traits. Nous savons donc désormais \emph{lire} les mots. Il reste à franchir l'étape décisive : \emph{comprendre} un texte entier et \emph{répondre} à des questions dessus, comme à l'examen. C'est l'objet de cette section : une méthode de lecture rigoureuse, puis des questions graduelles, du plus facile (Vrai/Faux) au plus exigeant (expression personnelle).

\courssub{Rappel du texte support}

Avant de répondre aux questions, relisons le texte de la leçon. Pour vous entraîner, voici la version intégrale du texte support, à relire attentivement :

\begin{textebox}[认识疾病，预防艾滋]
生病是每个人都可能遇到的事情。感冒、发烧、头痛是小病，很容易治好；但是有些病很严重，比如艾滋病。\\
Shēng bìng shì měi ge rén dōu kě néng yù dào de shì qing. Gǎn mào, fā shāo, tóu téng shì xiǎo bìng, hěn róng yì zhì hǎo; dàn shì yǒu xiē bìng hěn yán zhòng, bǐ rú ài zī bìng.\\
« Tomber malade est une chose que tout le monde peut rencontrer. Le rhume, la fièvre, le mal de tête sont des petites maladies, faciles à guérir ; mais certaines maladies sont graves, comme le SIDA. »\\[4pt]
艾滋病是一种严重的传染病。艾滋病毒（HIV）主要攻击人体的免疫系统，破坏免疫细胞，使人体慢慢丧失抵抗疾病的能力。\\
Ài zī bìng shì yì zhǒng yán zhòng de chuán rǎn bìng. Ài zī bìng dú (HIV) zhǔ yào gōng jī rén tǐ de miǎn yì xì tǒng, pò huài miǎn yì xì bāo, shǐ rén tǐ màn màn sàng shī dǐ kàng jí bìng de néng lì.\\
« Le SIDA est une grave maladie infectieuse. Le virus du SIDA (VIH) attaque principalement le système immunitaire du corps humain, détruit les cellules immunitaires et fait perdre peu à peu au corps sa capacité à résister aux maladies. »\\[4pt]
艾滋病的传播途径主要有三种：血液传播、性传播和母婴传播。日常生活接触，比如握手、一起吃饭、拥抱，不会传播艾滋病。\\
Ài zī bìng de chuán bō tú jìng zhǔ yào yǒu sān zhǒng: xuè yè chuán bō, xìng chuán bō hé mǔ yīng chuán bō. Rì cháng shēng huó jiē chù, bǐ rú wò shǒu, yī qǐ chī fàn, yōng bào, bù huì chuán bō ài zī bìng.\\
« Les voies de transmission du SIDA sont principalement au nombre de trois : la transmission par le sang, la transmission sexuelle et la transmission de la mère à l'enfant. Les contacts de la vie quotidienne, comme se serrer la main, manger ensemble, s'embrasser (accolade), ne transmettent pas le SIDA. »\\[4pt]
现在还没有艾滋病疫苗，也不能完全治好，治疗费用很高，所以预防为主。我们应该了解艾滋病知识，保护自己，也不应该歧视艾滋病人，要关心他们、帮助他们。\\
Xiàn zài hái méi yǒu ài zī bìng yì miáo, yě bù néng wán quán zhì hǎo, zhì liáo fèi yòng hěn gāo, suǒ yǐ yù fáng wéi zhǔ. Wǒ men yīng gāi liǎo jiě ài zī bìng zhī shi, bǎo hù zì jǐ, yě bù yīng gāi qí shì ài zī bìng bìng rén, yào guān xīn tā men, bāng zhù tā men.\\
« Il n'existe pas encore de vaccin contre le SIDA, on ne peut pas le guérir complètement non plus, et le coût du traitement est élevé : c'est pourquoi la prévention est primordiale. Nous devons connaître les informations sur le SIDA, nous protéger, et nous ne devons pas discriminer les malades du SIDA : il faut les soutenir et les aider. »
\end{textebox}

\courssub{La méthode de lecture : « Localiser — Extraire — Reformuler »}

Face à un texte chinois à l'examen, beaucoup d'élèves paniquent et traduisent tout, mot à mot. C'est une erreur de stratégie. Le bon réflexe est de partir de la \emph{question}, pas du texte.

\begin{definition}[定位 (dìngwèi) : la localisation]
\cle{定位} (dìng wèi), littéralement « fixer la position », désigne la technique de repérage rapide d'une information dans un texte grâce aux \cle{mots-clés} : noms propres, chiffres (三种), termes techniques (免疫系统), mots répétés. Au lieu de relire tout le texte, on scanne les lignes à la recherche du mot-clé, puis on lit seulement la phrase qui le contient.
\end{definition}

\begin{methode}[Stratégie de lecture « Localiser — Extraire — Reformuler »]
\begin{enumerate}
\item \textbf{Lire la question d'abord.} Soulignez le mot-clé de la question. Exemple : dans \zh{艾滋病毒主要攻击什么系统？}, le mot-clé est \zh{攻击} (gōng jī, attaquer).
\item \textbf{Chercher le mot-clé dans le texte} (定位). On balaie les lignes des yeux jusqu'à trouver \zh{攻击} : il apparaît au paragraphe 2.
\item \textbf{Lire la phrase cible entière.} \zh{艾滋病毒主要攻击人体的免疫系统。} On extrait l'information : \zh{免疫系统}.
\item \textbf{Répondre en ses propres mots, en chinois simple.} \zh{它主要攻击免疫系统。}(Tā zhǔ yào gōng jī miǎn yì xì tǒng.) — « Il attaque principalement le système immunitaire. » Une réponse courte et correcte vaut mieux qu'une phrase longue et fautive.
\end{enumerate}
\end{methode}

\begin{propriete}[Les quatre types de questions à l'examen]
Les questions de compréhension de texte en chinois se rangent en quatre catégories. Les reconnaître permet d'adapter sa réponse :
\begin{enumerate}
\item \zh{细节理解题} (xì jié lǐ jiě tí) — \textbf{question de détail} : la réponse est écrite mot pour mot dans le texte. Technique : 定位.
\item \zh{主旨大意题} (zhǔ zhǐ dà yì tí) — \textbf{question d'idée principale} : il faut résumer le but global du texte (souvent dans la première ou la dernière phrase).
\item \zh{词义猜测题} (cí yì cāi cè tí) — \textbf{deviner le sens d'un mot} : on déduit le sens d'un mot inconnu grâce au contexte, sans dictionnaire.
\item \zh{推理判断题} (tuī lǐ pàn duàn tí) — \textbf{question d'inférence} : la réponse n'est pas écrite directement, il faut la déduire logiquement du texte.
\end{enumerate}
\end{propriete}

\begin{attention}[Piège d'examen : « D'après le texte... »]
Quand la question commence par \zh{根据课文} (gēn jù kè wén, « d'après le texte ») ou \zh{文中说} (wén zhōng shuō, « le texte dit »), votre réponse doit être \textbf{uniquement basée sur le texte}, pas sur vos connaissances personnelles. Même si vous savez, par vos cours de SVT, qu'il existe d'autres modes de transmission ou des traitements récents, \textbf{ne les mentionnez pas} si le texte ne les cite pas. Le correcteur évalue votre compréhension du texte, pas votre culture générale. Citer la phrase du texte (定位) est la meilleure justification.
\end{attention}

\courssub{Compréhension globale : idée principale, but, public}

Commençons par les questions les plus larges (主旨大意题). Elles portent sur le texte dans son ensemble.

\begin{exoresolu}[Vrai ou Faux ? Justifiez par une phrase du texte]
\textbf{Énoncé.} Dites si les affirmations suivantes sont vraies (对 duì) ou fausses (不对 bù duì), puis citez la phrase du texte qui le prouve.
\begin{enumerate}
\item 只有身体弱的人才会生病。(Zhǐ yǒu shēn tǐ ruò de rén cái huì shēng bìng.) — « Seules les personnes faibles tombent malades. »
\item 握手会传播艾滋病。(Wò shǒu huì chuán bō ài zī bìng.) — « Se serrer la main transmet le SIDA. »
\item 现在可以完全治好艾滋病。(Xiàn zài kě yǐ wán quán zhì hǎo ài zī bìng.) — « Aujourd'hui on peut guérir complètement le SIDA. »
\end{enumerate}
\tcblower
\textbf{Solution détaillée.}
\begin{enumerate}
\item \textbf{不对 (Faux).} Le texte dit : \zh{生病是每个人都可能遇到的事情。}(« Tomber malade est une chose que tout le monde peut rencontrer. ») Donc tout le monde peut tomber malade, pas seulement les personnes faibles.
\item \textbf{不对 (Faux).} Le texte dit : \zh{日常生活接触，比如握手、一起吃饭、拥抱，不会传播艾滋病。}(« Les contacts quotidiens, comme se serrer la main... ne transmettent pas le SIDA. »)
\item \textbf{不对 (Faux).} Le texte dit : \zh{现在还没有艾滋病疫苗，也不能完全治好。}(« Il n'y a pas encore de vaccin, et on ne peut pas guérir complètement. »)
\end{enumerate}
Remarquez la méthode : pour chaque réponse, on cite la phrase exacte du texte (定位). C'est ce qui donne tous les points.
\end{exoresolu}

\begin{exemplebox}[Choix multiple : l'idée principale]
\textbf{Q :} 这篇课文的主要意思是什么？(Zhè piān kè wén de zhǔ yào yì si shì shén me ?) — « Quelle est l'idée principale de ce texte ? »\\
A. 介绍感冒怎么治 (présenter comment soigner le rhume)\\
B. 介绍艾滋病和预防方法 (présenter le SIDA et les moyens de prévention)\\
C. 讲医院的历史 (raconter l'histoire des hôpitaux)\\[4pt]
\textbf{Réponse : B.} Le titre \zh{认识疾病，预防艾滋} (« Connaître les maladies, prévenir le SIDA ») et le dernier paragraphe (\zh{预防为主}) le confirment. Le but du texte est \textbf{d'informer et de sensibiliser} ; le public cible est \textbf{les jeunes / les élèves} (\zh{我们应该了解...} « nous devons connaître... »).
\end{exemplebox}

\courssub{Compréhension détaillée : réponses courtes en chinois}

Passons aux questions de détail (细节理解题). La règle d'or : répondre par une \textbf{phrase courte et complète}, en reprenant les mots du texte.

\begin{exemplebox}[Questions — Réponses modèles]
\textbf{Q1 :} 艾滋病是一种什么病？(Ài zī bìng shì yì zhǒng shén me bìng ?) — « Quel type de maladie est le SIDA ? »\\
\textbf{R :} 艾滋病是一种严重的传染病。(Ài zī bìng shì yì zhǒng yán zhòng de chuán rǎn bìng.) — « Le SIDA est une grave maladie infectieuse. »\\[4pt]
\textbf{Q2 :} 艾滋病毒主要攻击人体的什么系统？(Ài zī bìng dú zhǔ yào gōng jī rén tǐ de shén me xì tǒng ?) — « Quel système du corps humain le VIH attaque-t-il principalement ? »\\
\textbf{R :} 它主要攻击免疫系统。(Tā zhǔ yào gōng jī miǎn yì xì tǒng.) — « Il attaque principalement le système immunitaire. »\\[4pt]
\textbf{Q3 :} 艾滋病的传播途径有哪三种？(Ài zī bìng de chuán bō tú jìng yǒu nǎ sān zhǒng ?) — « Quelles sont les trois voies de transmission du SIDA ? »\\
\textbf{R :} 血液传播、性传播和母婴传播。(Xuè yè chuán bō, xìng chuán bō hé mǔ yīng chuán bō.) — « La transmission par le sang, la transmission sexuelle et la transmission de la mère à l'enfant. »\\[4pt]
\textbf{Q4 :} 日常生活接触会传播艾滋病吗？(Rì cháng shēng huó jiē chù huì chuán bō ài zī bìng ma ?) — « Les contacts quotidiens transmettent-ils le SIDA ? »\\
\textbf{R :} 不会。(Bù huì.) — « Non. » (Réponse minimale acceptée ; mieux : \zh{不会，比如握手、一起吃饭都不会传播。})\\[4pt]
\textbf{Q5 :} 我们应该怎么对待艾滋病人？(Wǒ men yīng gāi zěn me duì dài ài zī bìng bìng rén ?) — « Comment devons-nous nous comporter envers les malades du SIDA ? »\\
\textbf{R :} 我们不应该歧视他们，要关心他们、帮助他们。(Wǒ men bù yīng gāi qí shì tā men, yào guān xīn tā men, bāng zhù tā men.) — « Nous ne devons pas les discriminer, il faut les soutenir et les aider. »
\end{exemplebox}

\begin{attention}[Erreur fréquente des francophones : la réponse « oui/non »]
En français, on répond « oui » ou « non ». En chinois, il n'existe pas de mot unique pour « oui » : on \textbf{répète le verbe de la question}. Question avec \zh{会} : réponse \zh{会} / \zh{不会}. Question avec \zh{是} : réponse \zh{是} / \zh{不是}. Question avec \zh{有} : réponse \zh{有} / \zh{没有}. Ne traduisez jamais « oui » par un mot isolé inventé : reprenez le verbe !
\end{attention}

\courssub{Deviner le sens des mots en contexte (词义猜测题)}

À l'examen, un mot inconnu n'est pas une catastrophe : le contexte donne presque toujours la réponse. Entraînons-nous avec quatre mots du texte.

\begin{center}
\begin{tabularx}{\linewidth}{|X|X|X|X|}
\hline
\rowcolor{popblueL} Mot & Contexte dans le texte & Indice du contexte & Sens deviné \\
\hline
破坏 (pò huài) & 破坏免疫细胞 & le virus « fait quelque chose de mal » aux cellules & détruire, endommager \\
\hline
丧失 (sàng shī) & 丧失抵抗疾病的能力 & conséquence de la destruction : on « n'a plus » la capacité & perdre (définitivement) \\
\hline
途径 (tú jìng) & 传播途径主要有三种 & suivi d'une liste de trois « façons » de transmettre & voie, moyen, chemin \\
\hline
歧视 (qí shì) & 不应该歧视艾滋病人，要关心他们 & opposé à 关心 (soutenir) et 帮助 (aider) & discriminer, mépriser \\
\hline
\end{tabularx}
\end{center}

\begin{methode}[Deviner un mot inconnu en 4 étapes]
\begin{enumerate}
\item Identifier sa \textbf{nature grammaticale} (verbe ? nom ?) grâce à sa place dans la phrase.
\item Regarder les \textbf{mots voisins} : \zh{破坏免疫细胞} — \zh{细胞} (cellules) est un nom, donc \zh{破坏} est un verbe d'action.
\item Chercher les \textbf{oppositions} : \zh{不应该歧视... 要关心} — \zh{歧视} est le contraire de \zh{关心}.
\item Vérifier que le sens deviné \textbf{fonctionne dans la phrase} en la retraduisant mentalement.
\end{enumerate}
\end{methode}

\courssub{Restructuration : remettre les étapes dans l'ordre}

Les questions de restructuration demandent de réorganiser des informations dispersées dans le texte. Ici, il faut reconstituer la chaîne logique de l'infection décrite au paragraphe 2.

\begin{exoresolu}[Remettez les étapes dans l'ordre logique]
\textbf{Énoncé.} Voici quatre étapes mélangées. Remettez-les dans l'ordre décrit par le texte :\\
A. 人体丧失抵抗疾病的能力 (le corps perd sa capacité de résistance)\\
B. 艾滋病毒进入人体 (le VIH entre dans le corps)\\
C. 人得了艾滋病 (la personne développe le SIDA)\\
D. 病毒破坏免疫细胞，别的病来了 (le virus détruit les cellules immunitaires, d'autres maladies arrivent)
\tcblower
\textbf{Solution.} Ordre correct : \textbf{B → D → A → C}.\\
Justification par le texte (paragraphe 2) : le virus \zh{攻击免疫系统} (attaque le système immunitaire) → \zh{破坏免疫细胞} (détruit les cellules immunitaires) → \zh{使人体慢慢丧失抵抗疾病的能力} (fait perdre peu à peu la capacité de résistance) → la personne devient malade du SIDA. Notez le mot \zh{慢慢} (màn màn, « peu à peu ») : il indique que le processus est \emph{progressif}, ce qui confirme l'ordre chronologique.
\end{exoresolu}

\begin{popfigure}
\begin{tikzpicture}[node distance=0.45cm, every node/.style={font=\small}]
\node[draw=popblue, fill=popblueL, rounded corners, align=center, text width=4.6cm] (q) {\textbf{Question}\\ 攻击什么系统？};
\node[draw=poporange, fill=poporangeL, rounded corners, align=center, text width=4.6cm, right=1.1cm of q] (l) {\textbf{1. 定位}\\ mot-clé : 攻击\\ paragraphe 2};
\node[draw=popgreen, fill=popgreenL, rounded corners, align=center, text width=4.6cm, right=1.1cm of l] (e) {\textbf{2. Extraire}\\ 攻击人体的\\ 免疫系统};
\node[draw=poppurple, fill=poppurpleL, rounded corners, align=center, text width=4.6cm, below=0.9cm of e] (r) {\textbf{3. Reformuler}\\ 它主要攻击\\ 免疫系统。};
\draw[-{Stealth}, thick, popdark] (q) -- (l);
\draw[-{Stealth}, thick, popdark] (l) -- (e);
\draw[-{Stealth}, thick, popdark] (e) -- (r);
\end{tikzpicture}
\legende{La chaîne « Localiser — Extraire — Reformuler » appliquée à une question de détail.}
\end{popfigure}

\courssub{Expression personnelle guidée}

Dernière marche : la question ouverte, où il faut \emph{argumenter} en s'appuyant sur le texte.

\textbf{Question :} 根据课文，为什么预防为主？(Gēn jù kè wén, wèi shén me yù fáng wéi zhǔ ?) — « D'après le texte, pourquoi la prévention est-elle primordiale ? »

Le texte donne trois arguments, reliés par 所以 (suǒ yǐ, « donc »). Modèle de réponse en trois points :

\begin{exemplebox}[Réponse modèle argumentée]
因为第一，现在还没有艾滋病疫苗；第二，艾滋病不能完全治好；第三，治疗费用很高。所以预防为主。\\
Yīn wèi dì yī, xiàn zài hái méi yǒu ài zī bìng yì miáo; dì èr, ài zī bìng bù néng wán quán zhì hǎo; dì sān, zhì liáo fèi yòng hěn gāo. Suǒ yǐ yù fáng wéi zhǔ.\\
« Parce que premièrement, il n'y a pas encore de vaccin contre le SIDA ; deuxièmement, le SIDA ne peut pas être complètement guéri ; troisièmement, le coût du traitement est élevé. Donc la prévention est primordiale. »
\end{exemplebox}

Remarquez la structure argumentative très utile à réutiliser dans toutes vos productions : 因为第一... 第二... 第三... 所以... (« Parce que premièrement... deuxièmement... troisièmement... donc... »). C'est l'équivalent chinois du plan dialectique français, en plus direct.

\begin{savaistu}[Le Cameroun et la lutte contre le VIH]
Au Cameroun, le taux de prévalence du VIH a fortement diminué en vingt ans, passant d'environ 5,4\,\% en 2004 à environ 2,7\,\% en 2023, grâce notamment au dépistage gratuit et à la distribution gratuite des antirétroviraux (ARV). Cette politique de santé publique illustre concrètement le principe 预防为主 (yù fáng wéi zhǔ, « la prévention d'abord ») défendu par notre texte : quand guérir est difficile et coûteux, empêcher la transmission reste l'arme la plus efficace. La Chine mène elle aussi de vastes campagnes de dépistage et d'éducation, notamment dans les écoles et les universités.
\end{savaistu}

\courssub{Correction collective : justifier par la citation du texte}

En classe, la correction se fait au tableau : pour chaque question, un élève donne sa réponse, un autre doit \textbf{citer la phrase du texte} qui la justifie (定位). Voici le tableau de correspondance complet, à recopier dans votre cahier :

\begin{center}
\begin{tabularx}{\linewidth}{|c|X|X|X|}
\hline
\rowcolor{popblueL} N° & Paragraphe & Mots-clés du texte & Réponse attendue \\
\hline
1 (définition) & §2 & 严重的传染病 & une grave maladie infectieuse \\
\hline
2 (cible du virus) & §2 & 攻击免疫系统 & le système immunitaire \\
\hline
3 (transmission) & §3 & 三种：血液、性、母婴 & sang, sexuel, mère-enfant \\
\hline
4 (contacts quotidiens) & §3 & 不会传播 & non (不会) \\
\hline
5 (prévention) & §4 & 没有疫苗、不能治好、费用高 & 3 arguments + 所以预防为主 \\
\hline
6 (attitude) & §4 & 不应该歧视、要关心帮助 & ne pas discriminer, soutenir \\
\hline
\end{tabularx}
\end{center}

\begin{exercice}[À vous de jouer !]
Répondez en chinois, par des phrases complètes, en citant le texte :
\begin{enumerate}
\item 感冒和艾滋病一样吗？为什么？(Gǎn mào hé ài zī bìng yí yàng ma ? Wèi shén me ?)
\item 艾滋病毒使人体丧失什么能力？(Ài zī bìng dú shǐ rén tǐ sàng shī shén me néng lì ?)
\item 根据课文，我们应该怎么保护自己？(Gēn jù kè wén, wǒ men yīng gāi zěn me bǎo hù zì jǐ ?)
\item Devinez le sens de 疫苗 (yì miáo) dans la phrase 现在还没有艾滋病疫苗 grâce au contexte.
\end{enumerate}
\end{exercice}

\begin{corrige}[Exercice « À vous de jouer ! »]
\begin{enumerate}
\item 不一样。感冒是小病，很容易治好；艾滋病是严重的传染病。(Bù yí yàng. Gǎn mào shì xiǎo bìng, hěn róng yì zhì hǎo; ài zī bìng shì yán zhòng de chuán rǎn bìng.) — « Non. Le rhume est une petite maladie, facile à guérir ; le SIDA est une grave maladie infectieuse. » (Justification : §1.)
\item 它使人体丧失抵抗疾病的能力。(Tā shǐ rén tǐ sàng shī dǐ kàng jí bìng de néng lì.) — « Il fait perdre au corps la capacité de résister aux maladies. » (§2.)
\item 我们应该了解艾滋病知识，保护自己。(Wǒ men yīng gāi liǎo jiě ài zī bìng zhī shi, bǎo hù zì jǐ.) — « Nous devons connaître les informations sur le SIDA et nous protéger. » (§4.)
\item 疫苗 = « vaccin ». Indice du contexte : la phrase dit qu'il n'en existe « pas encore » (还没有) et que la maladie ne peut pas être guérie ; or un vaccin est précisément ce qui \emph{prévient} une maladie — d'où la conclusion logique 预防为主.
\end{enumerate}
\end{corrige}

\begin{aretenir}[L'essentiel de la section]
\begin{itemize}
\item Partez toujours de la \textbf{question}, pas du texte : méthode \cle{定位} (localiser) → extraire → reformuler en chinois simple.
\item Quatre types de questions : 细节理解题 (détail), 主旨大意题 (idée principale), 词义猜测题 (sens deviné), 推理判断题 (inférence).
\item « D'après le texte » = réponse \textbf{uniquement} tirée du texte, justifiée par une citation.
\item Pour répondre « oui/non », répétez le \textbf{verbe} de la question : 会 / 不会, 是 / 不是, 有 / 没有.
\item Structure argumentative à retenir : 因为第一... 第二... 第三... 所以...
\end{itemize}
\end{aretenir}

\coursec{Collocations, Structures utiles et Production guidée}

Dans la section précédente, vous avez répondu à des questions de compréhension sur le texte \zh{认识疾病，预防艾滋}. Vous savez donc \emph{comprendre} ce texte. Il reste maintenant à passer de la compréhension à la \cle{production} : réutiliser les structures et les collocations du texte pour parler et écrire à votre tour. C'est l'objectif de cette section : d'abord observer trois structures modèles, puis ancrer les collocations verbe + complément, puis s'entraîner aux transformations de phrases, et enfin produire une affiche de prévention pour la Journée mondiale du sida.

\courssub{Trois structures modèles extraites du texte}

\textbf{Observation.}\par
Relisons trois phrases du texte support :

\begin{exemplebox}[Observer avant de comprendre]
\zh{艾滋病不仅会破坏免疫系统，而且会让人丧失抵抗疾病的能力。}\\
\emph{Àizībìng bùjǐn huì pòhuài miǎnyì xìtǒng, érqiě huì ràng rén sàngshī dǐkàng jíbìng de nénglì.}\\
« Le sida non seulement détruit le système immunitaire, mais en plus fait perdre à l'homme la capacité de résister aux maladies. »

\zh{只有通过正确的途径，才能预防艾滋病。}\\
\emph{Zhǐyǒu tōngguò zhèngquè de tújìng, cái néng yùfáng Àizībìng.}\\
« Ce n'est qu'en passant par les bonnes voies que l'on peut prévenir le sida. »

\zh{艾滋病通过血液、母婴和性接触传播。}\\
\emph{Àizībìng tōngguò xuèyè, mǔyīng hé xìng jiēchù chuánbō.}\\
« Le sida se transmet par le sang, de la mère à l'enfant et par les contacts sexuels. »
\end{exemplebox}

\begin{propriete}[Structure 1 : 不仅……而且…… « non seulement... mais aussi »]
\textbf{Forme :} Sujet + \zh{不仅} + proposition 1, \zh{而且} + proposition 2.\\
\textbf{Emploi :} ajoute une information \emph{plus forte} à la première. Si les deux propositions ont le même sujet, \zh{不仅} se place après le sujet ; si les sujets sont différents, \zh{不仅} se place avant le premier sujet et \zh{而且} introduit le nouveau sujet.\\
\textbf{Comparaison avec le français :} le français dit « non seulement... mais aussi... » ; le chinois garde les deux morceaux dans le même ordre, sans inversion du sujet.
\end{propriete}

\begin{exemplebox}[不仅……而且…… en action]
\zh{运动不仅对身体好，而且让人心情愉快。} \emph{Yùndòng bùjǐn duì shēntǐ hǎo, érqiě ràng rén xīnqíng yúkuài.} « Le sport est non seulement bon pour le corps, mais il rend aussi de bonne humeur. »\\
\zh{不仅医生要关爱患者，而且全社会都要消除歧视。} \emph{Bùjǐn yīshēng yào guān'ài huànzhě, érqiě quán shèhuì dōu yào xiāochú qíshì.} « Non seulement les médecins doivent prendre soin des malades, mais toute la société doit éliminer la discrimination. » (sujets différents : \zh{不仅} avant le premier sujet)
\end{exemplebox}

\begin{propriete}[Structure 2 : 只有……才…… « seulement si... alors » — condition nécessaire]
\textbf{Forme :} \zh{只有} + condition, \zh{才} + résultat.\\
\textbf{Emploi :} exprime une condition \cle{nécessaire} : sans la condition, le résultat est impossible. Le sujet est souvent omis dans la deuxième proposition s'il est identique.\\
Exemple : \zh{只有坚持用药，才能控制病情。} \emph{Zhǐyǒu jiānchí yòngyào, cái néng kòngzhì bìngqíng.} « Ce n'est qu'en prenant régulièrement ses médicaments que l'on peut contrôler la maladie. »\\
\textbf{Attention :} ne pas confondre avec \zh{只要……就……} (condition \emph{suffisante} : « il suffit de... pour que... »). \zh{只有} = « seulement si » ; \zh{只要} = « du moment que ».
\end{propriete}

\begin{attention}[Erreur typique des francophones]
Le français « Si tu ne te fais pas dépister, tu ne sais pas » se traduit souvent par la phrase maladroite *\zh{如果你不检查，你不知道}. En chinois naturel, on préfère la condition nécessaire : \zh{只有检查了，才知道。} \emph{Zhǐyǒu jiǎnchá le, cái zhīdào.} « Ce n'est qu'en se faisant dépister que l'on sait. » Retenez : « seulement si testé, alors savoir ».
\end{attention}

\begin{propriete}[Structure 3 : 通过…… « par la voie de, au moyen de »]
\textbf{Forme :} \zh{通过} + nom / groupe verbal, (Sujet) + verbe. Variante du texte : \zh{通过……途径} « par la voie de... ».\\
\textbf{Emploi :} introduit le \cle{moyen} ou le \cle{canal} d'une action. Très fréquent dans les textes de santé et les textes officiels.\\
Exemples : \zh{通过检查，我们知道了自己的健康状况。} « Grâce au dépistage, nous connaissons notre état de santé. » / \zh{通过宣传教育，人们消除了恐惧。} « Par la sensibilisation, les gens ont fait disparaître la peur. »
\end{propriete}

\courssub{Collocations verbe + complément à ancrer}

En chinois, un verbe ne s'emploie pas avec n'importe quel complément : on dit \zh{患病} « contracter une maladie » et non *\zh{得病} dans un registre soutenu (ce dernier est correct mais familier). Voici les dix collocations du texte à mémoriser comme des blocs.

\begin{center}
\begin{tabularx}{\linewidth}{|X|X|X|X|}
\hline
\rowcolor{popblueL} Collocation & Pinyin & Nature & Traduction \\
\hline
\zh{患病} & huànbìng & v. + n. & contracter une maladie \\
\hline
\zh{感染病毒} & gǎnrǎn bìngdú & v. + n. & être infecté par un virus \\
\hline
\zh{破坏系统} & pòhuài xìtǒng & v. + n. & détruire un système \\
\hline
\zh{丧失能力} & sàngshī nénglì & v. + n. & perdre une capacité \\
\hline
\zh{传播疾病} & chuánbō jíbìng & v. + n. & transmettre une maladie \\
\hline
\zh{采取措施} & cǎiqǔ cuòshī & v. + n. & prendre des mesures \\
\hline
\zh{进行检查} & jìnxíng jiǎnchá & v. + n. & effectuer un dépistage \\
\hline
\zh{避免接触} & bìmiǎn jiēchù & v. + n. & éviter le contact \\
\hline
\zh{关爱患者} & guān'ài huànzhě & v. + n. & prendre soin des malades \\
\hline
\zh{消除歧视} & xiāochú qíshì & v. + n. & éliminer la discrimination \\
\hline
\end{tabularx}
\end{center}

\begin{definition}[高危行为 — comportements à risque]
\zh{高危行为} \emph{gāowēi xíngwéi} : terme technique incontournable désignant les comportements qui exposent fortement au virus : rapports non protégés, partage de seringues, transfusion de sang non dépisté. Exemple : \zh{我们要避免高危行为。} « Nous devons éviter les comportements à risque. »
\end{definition}

\begin{popfigure}
\begin{tikzpicture}[mindmap, grow cyclic, every node/.style={concept}, root concept/.append style={concept color=popblue, font=\large\bfseries}, level 1/.append style={level distance=3.4cm, sibling angle=72}, concept color=popblueL]
\node[align=center]{\zh{健康}\\ Santé}
  child { node[align=center]{\zh{患病}\\ tomber malade} }
  child { node[align=center]{\zh{感染病毒}\\ être infecté} }
  child { node[align=center]{\zh{进行检查}\\ se faire dépister} }
  child { node[align=center]{\zh{采取措施}\\ prendre des mesures} }
  child { node[align=center]{\zh{关爱患者}\\ soigner les malades} };
\end{tikzpicture}
\legende{Carte mentale lexicale : les collocations autour du thème de la santé.}
\end{popfigure}

\courssub{Transformations de phrases : négation, passif, question}

\begin{propriete}[Rappel des trois transformations]
\textbf{Affirmative $\rightarrow$ Négative :} \zh{不} pour le présent/futur et l'habitude ; \zh{没(有)} pour une action non réalisée dans le passé (et alors \zh{了} disparaît). \zh{他去医院。} $\rightarrow$ \zh{他没去医院。}\\
\textbf{Actif $\rightarrow$ Passif :} Complément + \zh{被} + agent + verbe + (autres éléments). \zh{病毒破坏免疫系统。} $\rightarrow$ \zh{免疫系统被病毒破坏了。}\\
\textbf{Déclaration $\rightarrow$ Question :} ajouter \zh{吗} en fin de phrase, ou \zh{呢} pour renvoyer la question (« et toi ? »).
\end{propriete}

\begin{exemplebox}[Exemples traduits et transformés]
\zh{我们要避免高危行为。} « Nous devons éviter les comportements à risque. » $\rightarrow$ \zh{只有避免高危行为，才能预防感染。} « Ce n'est qu'en évitant les comportements à risque que l'on peut prévenir l'infection. »\\
\zh{艾滋病不通过蚊子传播。} « Le sida ne se transmet pas par les moustiques. » $\rightarrow$ \zh{蚊子不能传播艾滋病。} « Les moustiques ne peuvent pas transmettre le sida. » (le passif en \zh{被} est possible — \zh{艾滋病不会被蚊子传播} — mais moins naturel ici, car \zh{被} suggère souvent une nuance négative subie.)
\end{exemplebox}

\begin{exoresolu}[Transformations guidées]
Énoncé : 1) Mettez à la forme négative : \zh{他昨天进行了检查。} 2) Transformez avec \zh{只有……才……} : \zh{多锻炼，身体健康。} 3) Transformez en question avec \zh{吗} : \zh{蚊子传播艾滋病。}
\tcblower
Solution détaillée :\\
1) Action passée non réalisée $\rightarrow$ \zh{没} et suppression de \zh{了} : \zh{他昨天没有进行检查。} \emph{Tā zuótiān méiyǒu jìnxíng jiǎnchá.} « Hier, il ne s'est pas fait dépister. »\\
2) Condition nécessaire : \zh{只有多锻炼，身体才健康。} \emph{Zhǐyǒu duō duànliàn, shēntǐ cái jiànkāng.} « Ce n'est qu'en faisant plus de sport que le corps est en bonne santé. » Le sujet \zh{身体} est conservé car il change de place logique ; on pourrait aussi dire \zh{只有多锻炼，才能保持健康。}\\
3) \zh{蚊子传播艾滋病吗？} \emph{Wénzi chuánbō Àizībìng ma?} « Les moustiques transmettent-ils le sida ? » (Réponse attendue : \zh{不传播。})
\end{exoresolu}

\courssub{Production écrite guidée : une affiche pour le 1er décembre}

\begin{methode}[Écrire une affiche de prévention en 4 étapes]
\begin{enumerate}
\item \textbf{Verbes d'action à l'impératif :} utilisez \zh{不} (ne... pas), \zh{请} (prière de), \zh{要} (il faut), \zh{多} (davantage de).
\item \textbf{Phrases courtes :} 4 à 7 caractères par ligne, pour être lues d'un coup d'œil.
\item \textbf{Rime ou parallélisme (\zh{对仗}) :} la symétrie aide la mémoire. Modèle célèbre : \zh{早检查，早治疗，早受益。} \emph{Zǎo jiǎnchá, zǎo zhìliáo, zǎo shòuyì.} « Dépistage tôt, traitement tôt, bénéfice tôt. »
\item \textbf{Un slogan final + un numéro utile :} terminez par une phrase choc et le numéro vert local (au Cameroun, le 114 pour les urgences sanitaires).
\end{enumerate}
\end{methode}

\begin{textebox}[Modèle d'affiche — 世界艾滋病日]
\zh{预防艾滋，人人有责！}\\
\emph{Yùfáng Àizī, rénrén yǒuzé !}\\
« Prévenir le sida : la responsabilité de chacun ! »\\
\zh{不共用针具。} \emph{Bù gòngyòng zhēnjù.} « Ne partagez pas les seringues. »\\
\zh{请及早检查。} \emph{Qǐng jízǎo jiǎnchá.} « Faites-vous dépister tôt. »\\
\zh{要关爱患者。} \emph{Yào guān'ài huànzhě.} « Prenez soin des malades. »\\
\zh{只有消除歧视，才能战胜艾滋病！}\\
\emph{Zhǐyǒu xiāochú qíshì, cái néng zhànshèng Àizībìng !}\\
« Ce n'est qu'en éliminant la discrimination que nous vaincrons le sida ! »\\
\zh{咨询电话：114} \emph{Zīxún diànhuà : yāo yāo sì.} « Ligne d'information : 114. »
\end{textebox}

\begin{popfigure}
\begin{tikzpicture}[scale=0.92, every node/.style={transform shape}]
\draw[fill=popcream, draw=popblue, line width=1.5pt, rounded corners=6pt] (0,0) rectangle (11,13);
\node[fill=popblue, text=white, rounded corners=3pt, font=\bfseries, minimum width=9.5cm, minimum height=1cm] at (5.5,12.1) {\zh{预防艾滋，人人有责！}};
\node[draw=poporange, fill=poporangeL, rounded corners=3pt, minimum width=8.5cm, minimum height=0.8cm] at (5.5,10.6) {\zh{不共用针具。}};
\node[draw=poporange, fill=poporangeL, rounded corners=3pt, minimum width=8.5cm, minimum height=0.8cm] at (5.5,9.5) {\zh{请及早检查。}};
\node[draw=poporange, fill=poporangeL, rounded corners=3pt, minimum width=8.5cm, minimum height=0.8cm] at (5.5,8.4) {\zh{要关爱患者。}};
\node[draw=poppurple, fill=poppurpleL, rounded corners=3pt, minimum width=9.5cm, minimum height=0.9cm, font=\bfseries] at (5.5,6.9) {\zh{只有消除歧视，才能战胜艾滋病！}};
\draw[fill=poppink, draw=poppink] (4.9,5.6) -- (5.15,4.6) -- (5.5,5.2) -- (5.85,4.6) -- (6.1,5.6) -- (5.5,5.0) -- cycle;
\node[font=\small] at (7.6,5.1) {Ruban rouge, symbole mondial};
\node[draw=popgreen, fill=popgreenL, rounded corners=3pt, minimum width=6cm, minimum height=0.8cm] at (5.5,3.4) {\zh{咨询电话：114}};
\node[font=\small\itshape, text=popmuted] at (5.5,2.4) {Lycée de Yaoundé — 1er décembre};
\node[font=\small, text=popdark, align=left] at (5.5,1.2) {Titre accrocheur $\rightarrow$ 3 actions (\zh{不} / \zh{请} / \zh{要}) $\rightarrow$ slogan \zh{只有……才……} $\rightarrow$ ruban $\rightarrow$ numéro vert};
\end{tikzpicture}
\legende{Modèle de mise en page d'une affiche de prévention : titre, trois puces d'action (\zh{不} / \zh{请} / \zh{要}), slogan final, ruban rouge, numéro vert.}
\end{popfigure}

\begin{exercice}[À vous de produire — 15 minutes]
Rédigez une affiche de prévention (\zh{标语} ou \zh{倡议书}) pour la Journée mondiale du sida (1er décembre) dans votre lycée. Contraintes obligatoires :
\begin{itemize}
\item 60 à 80 caractères chinois ;
\item 3 impératifs avec \zh{不}, \zh{请}, \zh{要} (ou \zh{多}) ;
\item 1 structure \zh{只有……才……} ;
\item au moins 4 collocations de la leçon (\zh{进行检查}, \zh{避免接触}, \zh{关爱患者}, \zh{消除歧视}…) ;
\item un numéro utile à la fin.
\end{itemize}
Puis présentez votre affiche à l'oral (1 minute) en expliquant le choix de vos mots : « \zh{我用了……因为……} » (J'ai utilisé... parce que...).
\end{exercice}

\begin{corrige}[Grille d'auto-évaluation]
Après la rédaction, évaluez-vous sur quatre critères (note de 0 à 2 chacun) :
\begin{center}
\begin{tabularx}{\linewidth}{|X|X|X|}
\hline
\rowcolor{popblueL} Critère & 2 points & 1 point \\
\hline
Vocabulaire utilisé & 4 collocations ou plus, bien employées & 2 à 3 collocations \\
\hline
Structures correctes & \zh{只有……才……} et impératifs sans faute & une petite faute \\
\hline
Lisibilité des caractères & traits nets, ordre respecté, aucun caractère inventé & quelques hésitations \\
\hline
Message clair & on comprend l'affiche en 10 secondes & message à relire \\
\hline
\end{tabularx}
\end{center}
Objectif : 7 ou 8 sur 8. En dessous de 5, reprenez la méthode en 4 étapes et le modèle d'affiche.
\end{corrige}

\begin{savaistu}[Sida et frontières : le cas de la Chine]
La Chine a levé en 2010 l'interdiction d'entrée sur son territoire pour les personnes séropositives, une mesure saluée par les organisations internationales de santé. Toutefois, le dépistage reste obligatoire pour l'obtention de certains visas de long séjour (études X1, travail Z, journalisme J1). Au Cameroun comme en Chine, les campagnes du 1er décembre insistent sur le même message : le dépistage précoce sauve des vies, et la discrimination est un obstacle à la lutte contre l'épidémie.
\end{savaistu}

\coursec{Approfondissement socioculturel : Cameroun - Chine, regards croisés}

Dans la section précédente, nous avons appris à assembler le vocabulaire de la santé dans des collocations et des structures utiles (\zh{预防疾病}, \zh{积极治疗}, \zh{保护自己和他人}…) pour produire des phrases complètes. Nous allons maintenant élargir notre regard : comment le Cameroun et la Chine, deux pays amis mais différents, organisent-ils la lutte contre le sida ? Comparer les politiques de santé publique, c'est à la fois enrichir notre vocabulaire et développer notre compétence socioculturelle, indispensable pour un vrai dialogue interculturel.

\courssub{Texte d'ouverture : deux pays, un même combat}

\begin{textebox}[喀麦隆和中国：共同抗击艾滋病]
喀麦隆和中国都是发展中国家，都重视艾滋病的预防和治疗。\\
Kāmàilóng hé Zhōngguó dōu shì fāzhǎnzhōng guójiā, dōu zhòngshì àizībìng de yùfáng hé zhìliáo.\\
« Le Cameroun et la Chine sont tous deux des pays en développement, et tous deux attachent de l'importance à la prévention et au traitement du sida. »\\[4pt]
在喀麦隆，抗病毒药（ARV）是免费的，病人可以在医院接受治疗。很多非政府组织也帮助病人，鼓励他们“积极生活”。\\
Zài Kāmàilóng, kàngbìngdú yào (ARV) shì miǎnfèi de, bìngrén kěyǐ zài yīyuàn jiēshòu zhìliáo. Hěn duō fēi zhèngfǔ zǔzhī yě bāngzhù bìngrén, gǔlì tāmen “jījí shēnghuó”.\\
« Au Cameroun, les antirétroviraux (ARV) sont gratuits et les malades peuvent être soignés à l'hôpital. De nombreuses ONG aident aussi les malades et les encouragent à "vivre positivement". »\\[4pt]
在中国，疾控中心负责监测疾病，社区卫生服务中心为居民提供健康服务。中国的“零歧视”运动反对对艾滋病患者的歧视。\\
Zài Zhōngguó, jíkòng zhōngxīn fùzé jiāncè jíbìng, shèqū wèishēng fúwù zhōngxīn wèi jūmín tígōng jiànkāng fúwù. Zhōngguó de “líng qíshì” yùndòng fǎnduì duì àizībìng huànzhě de qíshì.\\
« En Chine, les centres de contrôle des maladies surveillent les épidémies et les centres de santé communautaires offrent des services aux habitants. La campagne "Zéro discrimination" combat la stigmatisation des personnes séropositives. »\\[4pt]
两个国家都很重视母婴阻断：如果妈妈接受治疗，孩子出生时几乎不会感染病毒。\\
Liǎng ge guójiā dōu hěn zhòngshì mǔyīng zǔduàn: rúguǒ māma jiēshòu zhìliáo, háizi chūshēng shí jīhū bú huì gǎnrǎn bìngdú.\\
« Les deux pays accordent une grande importance à la PTME : si la mère suit un traitement, l'enfant ne sera presque jamais infecté à la naissance. »\\[4pt]
虽然两国的做法不完全一样，但是目标相同：让每个人都能健康、有尊严地生活。\\
Suīrán liǎng guó de zuòfǎ bù wánquán yíyàng, dànshì mùbiāo xiāngtóng: ràng měi ge rén dōu néng jiànkāng, yǒu zūnyán de shēnghuó.\\
« Bien que les méthodes des deux pays ne soient pas exactement les mêmes, l'objectif est identique : permettre à chacun de vivre en bonne santé et dans la dignité. »
\end{textebox}

\begin{center}
\begin{tabularx}{\linewidth}{|X|X|X|X|}
\hline
\rowcolor{popblueL} Caractères & Pinyin & Nature & Traduction \\
\hline
疾控中心 & jíkòng zhōngxīn & n. & centre de contrôle des maladies (CDC) \\
\hline
社区卫生服务中心 & shèqū wèishēng fúwù zhōngxīn & n. & centre de santé communautaire \\
\hline
非政府组织 & fēi zhèngfǔ zǔzhī & n. & ONG \\
\hline
歧视 & qíshì & n./v. & discrimination ; discriminer \\
\hline
母婴阻断 & mǔyīng zǔduàn & n. & prévention de la transmission mère-enfant (PTME) \\
\hline
尊严 & zūnyán & n. & dignité \\
\hline
监测 & jiāncè & v. & surveiller (épidémiologie) \\
\hline
同伴教育 & tóngbàn jiàoyù & n. & éducation par les pairs \\
\hline
\end{tabularx}
\end{center}

\courssub{Comparer les politiques de santé publique}

\begin{definition}[母婴阻断]
\zh{母婴阻断} (mǔyīng zǔduàn) désigne la \cle{prévention de la transmission mère-enfant} (PTME, en anglais PMTCT) : ensemble de mesures médicales (traitement de la mère pendant la grossesse, accouchement encadré, traitement du nouveau-né) qui empêchent le virus de passer de la mère à l'enfant. C'est un succès majeur commun aux deux pays : le taux de transmission est inférieur à 2 \% en Chine et inférieur à 5 \% au Cameroun lorsque le protocole est suivi.
\end{definition}

Observons d'abord les grandes lignes, puis lisons le tableau comparatif.

\begin{itemize}
\item \cle{Gratuité des ARV} : au Cameroun comme en Chine, les antirétroviraux sont fournis gratuitement aux personnes séropositives. En Chine, c'est le programme national « Quatre gratuits, un soin » (\zh{四免一关怀}, sì miǎn yī guānhuái) ; au Cameroun, la gratuité est effective depuis 2007.
\item \cle{Dépistage} : les deux pays encouragent le dépistage volontaire (\zh{自愿检测}, zìyuàn jiǎncè). En Chine, le dépistage prénatal est systématiquement proposé (quasi obligatoire dans le suivi de grossesse) ; au Cameroun, il est fortement recommandé lors des consultations prénatales.
\item \cle{Structures} : en Chine, le réseau repose sur les \zh{疾控中心} (CDC) et les \zh{社区卫生服务中心} ; au Cameroun, sur les hôpitaux publics, les centres de santé intégrés et les ONG communautaires.
\end{itemize}

\begin{popfigure}
\begin{center}
\begin{tabularx}{\linewidth}{|X|X|X|X|}
\hline
\rowcolor{popblueL} Thème & Cameroun & Chine & Points communs \\
\hline
Prise en charge ARV & Gratuite depuis 2007, hôpitaux et ONG & Programme 四免一关怀, gratuité nationale & ARV gratuits pour tous les malades \\
\hline
Dépistage & Volontaire (VCT), proposé en prénatal & Volontaire, quasi systématique en prénatal & Test gratuit et confidentiel \\
\hline
Éducation des jeunes & Programme d'éducation à la vie (PEV) & Plan 健康中国2030, éducation à l'école & Prévention dès le collège \\
\hline
Stigmatisation & Campagne « Vivre positivement », ONG & Campagne 零歧视 « Zéro discrimination » & Lutte contre le rejet des malades \\
\hline
Médecine traditionnelle & Tradipraticiens, recours fréquent, encadrement en cours & Médecine traditionnelle chinoise (中医) reconnue et encadrée & Soutien, jamais substitut aux ARV \\
\hline
\end{tabularx}
\end{center}
\legende{Tableau comparatif bilatéral : politiques de lutte contre le sida au Cameroun et en Chine.}
\end{popfigure}

\begin{attention}[疾控中心 ou 医院 ?]
Ne confonds pas \zh{疾控中心} (jíkòng zhōngxīn), le Centre de Contrôle des Maladies (CDC), qui est un organisme \emph{administratif et épidémiologique} (surveillance, statistiques, campagnes de prévention), et \zh{医院} (yīyuàn), l'hôpital, où l'on reçoit le \emph{soin clinique} (consultations, traitements). On dit : \zh{我去医院看病。} (Wǒ qù yīyuàn kànbìng. « Je vais à l'hôpital me faire soigner. ») mais \zh{疾控中心发布疫情报告。} (Jíkòng zhōngxīn fābù yìqíng bàogào. « Le CDC publie le rapport épidémiologique. »)
\end{attention}

\courssub{Lutter contre la stigmatisation}

La stigmatisation (\zh{歧视}, qíshì) est souvent plus dure à vivre que la maladie elle-même. Les deux pays mènent des campagnes publiques : la Chine avec le mouvement \zh{零歧视} (líng qíshì, « Zéro discrimination »), le Cameroun avec le message « Vivre positivement », porté par de nombreuses associations et ONG (\zh{非政府组织}) qui pratiquent aussi l'éducation par les pairs (\zh{同伴教育}) : des jeunes formés parlent à d'autres jeunes, ce qui est très efficace.

\begin{exemplebox}[Deux phrases à retenir]
\zh{消除对艾滋病患者的歧视，是文明社会的标志。}\\
\emph{Xiāochú duì àizībìng huànzhě de qíshì, shì wénmíng shèhuì de biāozhì.}\\
« Éliminer la discrimination envers les personnes séropositives est la marque d'une société civilisée. »\\[4pt]
\zh{中医可以辅助治疗，但不能替代抗病毒治疗。}\\
\emph{Zhōngyī kěyǐ fǔzhù zhìliáo, dàn bù néng tìdài kàngbìngdú zhìliáo.}\\
« La médecine traditionnelle chinoise peut accompagner le traitement, mais ne peut pas remplacer le traitement antirétroviral. »
\end{exemplebox}

\courssub{Médecine traditionnelle et médecine moderne}

En Chine, la médecine traditionnelle chinoise (\zh{中医}, zhōngyī) est reconnue par l'État et encadrée : elle peut aider à \zh{提高免疫力} (tígāo miǎnyìlì, « renforcer l'immunité ») et à réduire certains effets secondaires, mais le traitement principal reste les ARV (\zh{西药}, xīyào, « médicaments occidentaux »). Au Cameroun, les tradipraticiens jouent aussi un rôle important dans la société. Dans les deux pays, le danger est le même : les « remèdes miracles » (\zh{神药}, shényào) vendus par des charlatans, qui font perdre du temps et de l'argent aux malades et peuvent provoquer des interactions dangereuses avec les ARV.

\begin{savaistu}[Tu Youyou et le prix Nobel]
Le prix Nobel de médecine 2015 a été attribué à la chercheuse chinoise \zh{屠呦呦} (Tú Yōuyōu) pour la découverte de l'artémisinine (\zh{青蒿素}, qīnghāosù), issue de la pharmacopée traditionnelle chinoise. Ce médicament a sauvé des millions de vies face au paludisme — une maladie grave chez les personnes vivant avec le VIH en co-infection. Cela montre que la tradition peut nourrir la science moderne, à condition d'être vérifiée par des essais cliniques (\zh{临床试验}, línchuáng shìyàn).
\end{savaistu}

\begin{methode}[Organiser un débat structuré]
Sujet : « La médecine traditionnelle a-t-elle sa place dans le traitement du sida ? »
\begin{enumerate}
\item Former deux équipes : POUR et CONTRE. Chaque équipe prépare trois arguments.
\item Arguments POUR possibles : soutien de l'immunité (\zh{提高免疫力}), réduction des effets secondaires, confort psychologique, tradition culturelle respectée.
\item Arguments CONTRE possibles : interactions médicamenteuses (\zh{相互作用}, xiānghù zuòyòng), retard du diagnostic, charlatanisme (\zh{神药}), absence d'essais cliniques (\zh{临床试验}).
\item Chaque élève parle 30 secondes en utilisant le vocabulaire : \zh{辅助} (fǔzhù, accompagner), \zh{替代} (tìdài, remplacer), \zh{相互作用}, \zh{临床试验}.
\item Conclure ensemble avec la phrase de synthèse : \zh{中医可以辅助治疗，但不能替代抗病毒治疗。}
\end{enumerate}
\end{methode}

\courssub{Éducation à la santé et questions éthiques}

L'éducation des jeunes est la clé de la prévention. En Chine, le plan national \zh{健康中国2030} (Jiànkāng Zhōngguó 2030, « Chine en bonne santé 2030 ») prévoit l'éducation à la santé dès l'école ; au Cameroun, le Programme d'éducation à la vie (PEV) aborde la prévention du VIH avec les collégiens et lycéens. Dans les deux cas, le public cible est le même : vous, les jeunes de 12 à 20 ans.

Sur le plan éthique, deux principes se rencontrent parfois : la \cle{confidentialité médicale} et la \cle{déclaration obligatoire} des maladies transmissibles. En Chine, le médecin remplit une fiche de signalement (\zh{传染病报告卡}, chuánrǎnbìng bàogàokǎ) pour le suivi épidémiologique, mais l'identité du malade reste protégée. Les droits des malades sont reconnus : droit à l'éducation (\zh{就学权}, jiùxuéquán) et droit au travail (\zh{就业权}, jiùyèquán).

\begin{propriete}[Vocabulaire droit et santé]
\begin{itemize}
\item \zh{知情同意} (zhīqíng tóngyì) : consentement éclairé — le patient doit comprendre et accepter avant un test ou un traitement.
\item \zh{保密原则} (bǎomì yuánzé) : principe de confidentialité — le médecin ne révèle pas le statut du patient.
\item \zh{强制检测} (qiángzhì jiǎncè) : dépistage obligatoire — interdit pour l'emploi et l'école en Chine depuis 2010 ; au Cameroun, la loi de 2021 protège aussi les personnes contre les tests forcés et les discriminations.
\end{itemize}
\end{propriete}

\begin{exoresolu}[Traduire et discuter]
Énoncé : traduis en français la phrase suivante, puis dis si tu es d'accord : \zh{为了保护病人，医生应该遵守保密原则。}
\tcblower
Solution : \emph{Wèile bǎohù bìngrén, yīshēng yīnggāi zūnshǒu bǎomì yuánzé.} Traduction : « Pour protéger les malades, les médecins doivent respecter le principe de confidentialité. » Points de traduction : \zh{为了} + nom/verbe = « pour, afin de » ; \zh{应该} = « devoir » (conseil moral) ; \zh{遵守} = « respecter (une règle) », plus soutenu que \zh{听}. Discussion attendue : oui, car sans confidentialité, les personnes ont peur de se faire dépister ; le secret médical encourage le dépistage volontaire, ce qui protège toute la société.
\end{exoresolu}

\begin{exercice}[À toi de jouer]
\begin{enumerate}
\item Complète avec le bon mot (\zh{疾控中心}, \zh{医院}, \zh{歧视}, \zh{母婴阻断}) : \zh{妈妈接受了治疗，\_\_\_\_ 保护了刚出生的孩子。}
\item Traduis en chinois : « Les ARV sont gratuits au Cameroun et en Chine. »
\item Rédige deux phrases avec \zh{辅助} et \zh{替代} pour donner ton avis sur la médecine traditionnelle.
\end{enumerate}
\end{exercice}

\begin{corrige}[Exercice « À toi de jouer »]
\begin{enumerate}
\item \zh{母婴阻断} : la PTME a protégé le nouveau-né.
\item \zh{在喀麦隆和中国，抗病毒药都是免费的。} (Zài Kāmàilóng hé Zhōngguó, kàngbìngdú yào dōu shì miǎnfèi de.)
\item Exemple : \zh{传统药可以辅助治疗，但是不能替代ARV。} (Chuántǒng yào kěyǐ fǔzhù zhìliáo, dànshì bù néng tìdài ARV. « Les remèdes traditionnels peuvent accompagner le traitement, mais ne peuvent pas remplacer les ARV. »)
\end{enumerate}
\end{corrige}

\begin{aretenir}[L'essentiel de la section]
\begin{itemize}
\item Cameroun et Chine partagent les mêmes objectifs : ARV gratuits, dépistage, PTME, lutte contre la stigmatisation.
\item \zh{疾控中心} surveille les épidémies ; \zh{医院} soigne les malades.
\item La médecine traditionnelle peut \zh{辅助} (accompagner), jamais \zh{替代} (remplacer) les ARV ; méfie-toi des \zh{神药} (remèdes miracles).
\item Les droits des malades : \zh{知情同意} (consentement éclairé), \zh{保密原则} (confidentialité), \zh{就学权} et \zh{就业权} (droit à l'éducation et au travail).
\end{itemize}
\end{aretenir}

\newpage

\coursec{Activité d'intégration}

\courssub{Objectifs de l'activité}

À la fin de cette activité d'intégration, tu seras capable de :
\begin{itemize}
\item \cle{réutiliser} le vocabulaire des maladies et du VIH/SIDA dans un document de communication réel (brochure de prévention) ;
\item \cle{mobiliser} les clés (radicaux) étudiées — 疒 (maladie), 心 (cœur), 扌 (main), 口 (bouche) — pour écrire correctement les caractères du thème ;
\item \cle{employer} les structures de prévention \zh{要 + V}, \zh{不要 + V}, \zh{请 + V + O} pour formuler des conseils ;
\item \cle{traduire} des messages courts du chinois vers le français et inversement ;
\item \cle{présenter oralement} en chinois, en 2 minutes, un message de santé publique ;
\item \cle{comparer} les dispositifs de lutte contre le SIDA au Cameroun et en Chine.
\end{itemize}

\courssub{Situation}

Dans le cadre de la \emph{Journée mondiale de lutte contre le SIDA} (1\ier{} décembre), le club de chinois de ton lycée à Yaoundé a été sollicité par le \emph{Comité National de Lutte contre le SIDA} (CNLS) et par l'ambassade de Chine au Cameroun. Les deux partenaires souhaitent diffuser dans les lycées une \cle{brochure bilingue chinois-français} destinée aux élèves de Première. Ton groupe a été choisi pour concevoir le modèle de cette brochure : « Ambassadeurs Santé », à vous de jouer !

Voici le message officiel reçu par le club :

\begin{textebox}[Message du comité d'organisation]
亲爱的同学们：\\
Qīn'ài de tóngxuémen :\\
« Chers camarades,\\[4pt]
十二月一日是世界艾滋病日。\\
Shí'èr yuè yī rì shì shìjiè Àizībìng rì.\\
Le 1\ier{} décembre, c'est la Journée mondiale du SIDA.\\[4pt]
请你们做一个小册子，告诉中学生：\\
Qǐng nǐmen zuò yí ge xiǎo cèzi, gàosu zhōngxuéshēng :\\
Nous vous demandons de créer une petite brochure pour dire aux lycéens :\\[4pt]
艾滋病是什么？怎么预防？\\
Àizībìng shì shénme ? Zěnme yùfáng ?\\
Qu'est-ce que le SIDA ? Comment le prévenir ?\\[4pt]
我们要健康，也要爱病人。\\
Wǒmen yào jiànkāng, yě yào ài bìngrén.\\
Nous voulons la santé, et nous devons aussi aimer les malades.\\[4pt]
谢谢你们的帮助！\\
Xièxie nǐmen de bāngzhù !\\
Merci de votre aide ! »
\end{textebox}

\courssub{Consignes}

Travail en groupes de 4 élèves, en 45 minutes. Répartition des rôles : un \cle{rédacteur chinois}, un \cle{traducteur}, un \cle{illustrateur-maquettiste}, un \cle{présentateur}.

\begin{enumerate}
\item \textbf{Maquette.} Plie une feuille A4 en trois volets (6 pages). Répartis les contenus : page 1 = titre et slogan ; page 2 = définition ; pages 3-4 = vrai/faux ; page 5 = gestes de prévention ; page 6 = numéros utiles et message anti-discrimination.
\item \textbf{Définition simple.} Rédige en chinois (2 ou 3 phrases maximum) une définition du VIH et du SIDA en utilisant les mots \zh{病毒} (bìngdú, virus), \zh{身体} (shēntǐ, corps), \zh{病} (bìng, maladie). Le traducteur écrit la version française en regard.
\item \textbf{Vrai ou faux.} Écris 3 phrases VRAIES et 3 phrases FAUSSES sur la transmission, avec \zh{会} (huì, pouvoir [transmettre]) ou \zh{不会} (bú huì, ne pas pouvoir). \cle{Surligne en couleur} les caractères contenant la clé 疒 (clé de la maladie) et la clé 心 (clé du cœur).
\item \textbf{Prévention.} Formule 5 conseils avec les structures \zh{要 + V} et \zh{不要 / 请 + V + O} (ex. : \zh{要洗手}, \zh{请医生检查}).
\item \textbf{Numéros utiles.} Indique les numéros d'aide : Cameroun 114 et 1510 ; Chine 12320, avec une phrase d'accompagnement en chinois.
\item \textbf{Message anti-discrimination.} Termine par un slogan en chinois (4 à 8 caractères) contre le rejet des malades, traduit en français.
\item \textbf{Exposé oral.} Le présentateur présente la brochure en chinois en 2 minutes : salutation, un point par volet, remerciement.
\end{enumerate}

\courssub{Ressources à mobiliser}

\begin{aretenir}[Boîte à outils de l'Ambassadeur Santé]
\begin{itemize}
\item \textbf{Lexique :} 艾滋病 (Àizībìng, SIDA), 病毒 (bìngdú, virus), 疾病 (jíbìng, maladie), 预防 (yùfáng, prévenir), 传染 (chuánrǎn, transmettre), 检查 (jiǎnchá, examiner), 健康 (jiànkāng, santé), 病人 (bìngrén, malade), 医生 (yīshēng, médecin), 医院 (yīyuàn, hôpital).
\item \textbf{Structures :} Sujet + 会 / 不会 + Verbe (capacité, possibilité) ; 要 + Verbe (il faut) ; 不要 + Verbe (il ne faut pas) ; 请 + Verbe + Complément (politesse).
\item \textbf{Clés :} 疒 (nè, clé de la maladie : 病, 疾, 疼) ; 心 (xīn, clé du cœur : 爱, 想) ; 扌 (shǒu, clé de la main : 护, 打) ; 口 (kǒu, clé de la bouche : 吃, 告).
\item \textbf{Culture :} au Cameroun comme en Chine, la lutte contre le SIDA passe par l'information, le dépistage gratuit et le refus de la discrimination.
\end{itemize}
\end{aretenir}

\begin{attention}[Pièges à éviter]
\begin{itemize}
\item Ne dis pas \zh{很预防} : 预防 est un verbe, pas un adjectif ; on dit \zh{要预防艾滋病} (yào yùfáng Àizībìng, il faut prévenir le SIDA).
\item Attention à l'ordre des mots : \zh{不会传染} (bú huì chuánrǎn, ne se transmet pas) — la négation 不 se place \emph{avant} 会.
\item Ton de 不 : devant un 4\ieme{} ton (会), 不 se prononce au 2\ieme{} ton : bú huì.
\item N'écris pas 病 sans la clé 疒 : la clé de la maladie enveloppe le caractère en haut et à gauche.
\end{itemize}
\end{attention}

\courssub{Proposition de production et corrigé commenté}

\begin{exoresolu}[Modèle de brochure « Ambassadeurs Santé »]
Énoncé : produire le texte chinois complet (avec traduction française) des six volets de la brochure.
\tcblower
\textbf{Volet 1 — Titre et slogan :}\\
\zh{健康第一，爱心无限！}\\
\emph{Jiànkāng dìyī, àixīn wúxiàn !}\\
« La santé d'abord, l'amour sans limites ! »\\
\emph{Commentaire : slogan en deux segments de 4 caractères, rythme typique des slogans chinois ; 爱心 (àixīn) contient la clé 心.}

\textbf{Volet 2 — Définition :}\\
\zh{艾滋病是一种很严重的疾病。艾滋病毒攻击人的身体。}\\
\emph{Àizībìng shì yì zhǒng hěn yánzhòng de jíbìng. Àizī bìngdú gōngjī rén de shēntǐ.}\\
« Le SIDA est une maladie très grave. Le virus du SIDA attaque le corps de l'homme. »\\
\emph{Commentaire : emploi du classificateur 种 (zhǒng) pour « sorte de maladie » ; 很 relie le sujet à l'adjectif 严重.}

\textbf{Volet 3-4 — Vrai ou faux :}\\
\zh{握手不会传染艾滋病。（真）}\\
\emph{Wòshǒu bú huì chuánrǎn Àizībìng. (zhēn)}\\
« Se serrer la main ne transmet pas le SIDA. (Vrai) »\\
\zh{一起吃饭不会传染艾滋病。（真）}\\
\emph{Yìqǐ chīfàn bú huì chuánrǎn Àizībìng. (zhēn)}\\
« Manger ensemble ne transmet pas le SIDA. (Vrai) »\\
\zh{血液会传染艾滋病。（真）}\\
\emph{Xuèyè huì chuánrǎn Àizībìng. (zhēn)}\\
« Le sang peut transmettre le SIDA. (Vrai) »\\
\zh{蚊子会传染艾滋病。（假）}\\
\emph{Wénzi huì chuánrǎn Àizībìng. (jiǎ)}\\
« Les moustiques transmettraient le SIDA. (Faux) »\\
\zh{拥抱会传染艾滋病。（假）}\\
\emph{Yōngbào huì chuánrǎn Àizībìng. (jiǎ)}\\
« Les embrassades transmettraient le SIDA. (Faux) »\\
\zh{上学不会传染艾滋病。（真）}\\
\emph{Shàngxué bú huì chuánrǎn Àizībìng. (zhēn)}\\
« Aller à l'école ne transmet pas le SIDA. (Vrai) »\\
\emph{Commentaire : alternance 会 / 不会 ; les caractères 病, 疾, 疼 portent la clé 疒 à surligner.}

\textbf{Volet 5 — Cinq conseils :}\\
\zh{要保护自己的身体。} \emph{Yào bǎohù zìjǐ de shēntǐ.} « Protège ton corps. »\\
\zh{不要共用针头。} \emph{Bú yào gòngyòng zhēntóu.} « Ne partage pas les aiguilles. »\\
\zh{请去医院检查。} \emph{Qǐng qù yīyuàn jiǎnchá.} « Va te faire dépister à l'hôpital. »\\
\zh{要听医生的话。} \emph{Yào tīng yīshēng de huà.} « Écoute les conseils du médecin. »\\
\zh{请告诉你的朋友。} \emph{Qǐng gàosu nǐ de péngyou.} « Informe tes amis. »

\textbf{Volet 6 — Numéros et message final :}\\
\zh{喀麦隆：114，1510。中国：12320。}\\
\emph{Kāmàilóng : yāo yāo sì, yāo wǔ yāo líng. Zhōngguó : yāo èr sān èr líng.}\\
« Cameroun : 114, 1510. Chine : 12320. »\\
\zh{不歧视病人，我们都是朋友！}\\
\emph{Bù qíshì bìngrén, wǒmen dōu shì péngyou !}\\
« Ne rejetons pas les malades, nous sommes tous amis ! »\\
\emph{Commentaire : dans les numéros de téléphone, le chiffre 1 se lit yāo ; 都 (dōu, tous) renforce le message d'unité.}
\end{exoresolu}

\courssub{Bilan de l'activité}

\begin{methode}[Grille d'évaluation de la brochure (20 points)]
\begin{enumerate}
\item \textbf{Lexique (4 pts)} : au moins 8 mots du thème correctement employés.
\item \textbf{Caractères (4 pts)} : écriture correcte, clés 疒, 心, 扌, 口 repérées et surlignées.
\item \textbf{Structures (4 pts)} : 会/不会, 要/不要, 请 + V + O utilisées sans faute d'ordre des mots.
\item \textbf{Culture (4 pts)} : numéros utiles des deux pays, message anti-discrimination pertinent.
\item \textbf{Oral (4 pts)} : présentation de 2 minutes en chinois, prononciation et tons soignés.
\end{enumerate}
\end{methode}

Pour conclure, retiens l'essentiel : tu sais désormais \cle{comprendre} un texte sur les maladies et le SIDA, \cle{écrire} les caractères clés du domaine médical, \cle{formuler} des conseils de prévention et \cle{défendre} un message de solidarité en chinois. Comme le dit le slogan de votre brochure : \zh{健康第一，爱心无限！} (Jiànkāng dìyī, àixīn wúxiàn !) — « La santé d'abord, l'amour sans limites ! » Au Cameroun comme en Chine, informer, c'est déjà protéger.

\newpage

\coursec{Méthodes et exercices résolus}

\coursec{Leçon 13 — Vocabulaire et compréhension de texte : maladies en général et du SIDA en particulier}

\courssub{Texte support : « 认识疾病，预防艾滋 »}

\begin{textebox}[Texte original — Niveau HSK 2-3]
\zh{认识疾病，预防艾滋} \\[4pt]
\emph{Rènshi jíbìng, yùfáng àizībìng} \\[4pt]
\zh{疾病是人类的敌人。感冒、发烧、头疼都是常见的病，去医院看医生、吃药就能好。但是，艾滋病不一样。艾滋病是一种很危险的传染病，它在全世界都有，喀麦隆也有。得了艾滋病，身体会越来越弱，很难治好。所以，我们应该了解这种病，学会预防它。年轻人要多学习健康知识，不做危险的事，保护自己和家人。预防胜于治疗，健康最重要。} \\[4pt]
\emph{Jíbìng shì rénlèi de dírén. Gǎnmào, fāshāo, tóuténg dōu shì chángjiàn de bìng, qù yīyuàn kàn yīshēng, chī yào jiù néng hǎo. Dànshì, àizībìng bù yīyàng. Àizībìng shì yì zhǒng hěn wēixiǎn de chuánrǎnbìng, tā zài quán shìjiè dōu yǒu, Kāmàilóng yě yǒu. Děi le àizībìng, shēntǐ huì yuè lái yuè ruò, hěn nán zhì hǎo. Suǒyǐ, wǒmen yīnggāi liǎojiě zhè zhǒng bìng, xuéhuì yùfáng tā. Niánqīngrén yào duō xuéxí jiànkāng zhīshi, bù zuò wēixiǎn de shì, bǎohù zìjǐ hé jiārén. Yùfáng shèng yú zhìliáo, jiànkāng zuì zhòngyào.} \\[4pt]
« Les maladies sont l'ennemi de l'humanité. Rhume, fièvre, maux de tête sont des maladies courantes ; aller à l'hôpital voir le médecin et prendre des médicaments suffit pour guérir. Mais le sida est différent. Le sida est une maladie infectieuse très dangereuse, il existe dans le monde entier, le Cameroun aussi. Si on attrape le sida, le corps devient de plus en plus faible, c'est très difficile à guérir. Donc, nous devons comprendre cette maladie, apprendre à la prévenir. Les jeunes doivent beaucoup étudier les connaissances sur la santé, ne pas faire de choses dangereuses, se protéger eux-mêmes et leur famille. Prévenir vaut mieux que guérir, la santé est le plus important. »
\end{textebox}

\begin{methode}[Stratégie de lecture : 5 étapes pour comprendre un texte sanitaire]
\begin{enumerate}
\item \textbf{Lecture globale} : Lisez une fois sans dictionnaire. Identifiez le thème (maladies courantes vs sida) et la structure (constat $\rightarrow$ danger $\rightarrow$ conseil).
\item \textbf{Repérage lexical} : Soulignez les mots-clés du champ \emph{santé} : \zh{疾病} (maladies), \zh{传染病} (maladie infectieuse), \zh{预防} (prévenir), \zh{治疗} (traiter/soigner), \zh{健康} (santé).
\item \textbf{Analyse syntaxique} : Repérez le \cle{Sujet} + \cle{Verbe} + \cle{Complément}. Notez la place des compléments de temps/lieu \textbf{avant} le verbe (\zh{去医院看医生}).
\item \textbf{Inférence contextuelle} : Deviner \zh{传染病} (maladie qui se transmet) grâce à \zh{病} et au contexte « dangereux, monde entier ». Deviner \zh{越来越弱} (de plus en plus faible) grâce à \zh{身体} et \zh{难治好}.
\item \textbf{Reformulation/Traduction} : Traduisez mentalement phrase par phrase avant de répondre aux questions.
\end{enumerate}
\emph{Astuce Bac :} Les réponses aux questions de compréhension suivent l'ordre du texte. Citez le texte (copiez les caractères) pour justifier.
\end{methode}

\courssub{Glossaire thématique : Tableau de vocabulaire maître}

\begin{definition}[Vocabulaire du Module 2 — Unité « Maladies et SIDA »]
Le tableau ci-dessous regroupe le lexique essentiel du texte support et du programme officiel (MINESEC). Apprenez \textbf{caractères + pinyin (tons) + nature + traduction} et les \textbf{collocations} (associations habituelles) indiquées entre parenthèses.
\end{definition}

\begin{center}
\begin{tabularx}{\linewidth}{|X|X|X|X|}
\hline
\rowcolor{popblueL} \textbf{Caractères} & \textbf{Pinyin} & \textbf{Nature} & \textbf{Traduction \& Collocations} \\
\hline
\zh{疾病} & jí bìng & nom & maladie (général) — \zh{常见疾病} (maladies courantes) \\
\zh{病} & bìng & nom/verbe & maladie ; être malade — \zh{生病} (tomber malade), \zh{病人} (malade/patient) \\
\zh{感冒} & gǎn mào & verbe/nom & rhume, refroidissement — \zh{感冒了} (avoir attrapé froid) \\
\zh{发烧} & fā shāo & verbe & avoir de la fièvre — \zh{发高烧} (forte fièvre) \\
\zh{头疼} & tóu téng & verbe & avoir mal à la tête — \zh{头很疼} (très mal à la tête) \\
\zh{常见} & cháng jiàn & adj. & courant, fréquent — \zh{常见的病} \\
\zh{医院} & yī yuàn & nom & hôpital — \zh{去医院} (aller à l'hôpital), \zh{在医院里} \\
\zh{医生} & yī shēng & nom & médecin — \zh{看医生} (consulter), \zh{好医生} \\
\zh{药} & yào & nom & médicament, remède — \zh{吃药} (prendre méd.), \zh{开药} (prescrire), \zh{药店} (pharmacie) \\
\zh{艾滋病} & ài zī bìng & nom & SIDA (AIDS) — abbr. \zh{艾滋} \\
\zh{传染病} & chuán rǎn bìng & nom & maladie infectieuse/contagieuse — \zh{预防传染病} \\
\zh{危险} & wēi xiǎn & adj. & dangereux — \zh{危险的病}, \zh{危险的事} \\
\zh{全世界} & quán shì jiè & nom & le monde entier — \zh{全世界都有} \\
\zh{喀麦隆} & kǎ mài lóng & n.p. & Cameroun \\
\zh{身体} & shēn tǐ & nom & corps, santé physique — \zh{身体弱} (corps faible), \zh{锻炼身体} (faire du sport) \\
\zh{越来越} & yuè lái yuè & adv. & de plus en plus — \zh{越来越弱}, \zh{越来越好} \\
\zh{弱} & ruò & adj. & faible, affaibli — \zh{身体弱} \\
\zh{治好} & zhì hǎo & verbe & guérir (résultat) — \zh{很难治好}, \zh{治好了} \\
\zh{了解} & liǎo jiě & verbe & comprendre, connaître — \zh{了解情况}, \zh{了解这种病} \\
\zh{预防} & yù fáng & verbe & prévenir — \zh{预防疾病}, \zh{预防胜于治疗} \\
\zh{年轻人} & nián qīng rén & nom & jeunes, la jeunesse — \zh{年轻人要注意} \\
\zh{健康} & jiàn kāng & adj./nom & en bonne santé, santé — \zh{健康知识}, \zh{保持健康}, \zh{健康生活} \\
\zh{知识} & zhī shi & nom & connaissances, savoir — \zh{学习知识} \\
\zh{保护} & bǎo hù & verbe & protéger — \zh{保护自己}, \zh{保护环境} \\
\zh{自己} & zì jǐ & pronom & soi-même — \zh{保护自己} \\
\zh{家人} & jiā rén & nom & famille, proches — \zh{保护家人} \\
\zh{治疗} & zhì liáo & verbe/nom & traiter, soigner, traitement — \zh{治疗疾病} \\
\zh{重要} & zhòng yào & adj. & important — \zh{最重要}, \zh{非常重要} \\
\hline
\end{tabularx}
\end{center}

\begin{aretenir}[Classificateurs (量词) indispensables pour ce thème]
\begin{itemize}
\item \zh{种} (zhǒng) : pour les types/espèces → \zh{一种病} (une maladie), \zh{这\zh{种}传染病}.
\item \zh{个} (gè) : général (personnes, objets) → \zh{一个医生}, \zh{一个病人}.
\item \zh{些} (xiē) : pluriel indéfini « quelques/des » → \zh{一些药} (des médicaments), \zh{一些知识}.
\item \zh{次} (cì) : fois (actions) → \zh{去了一次医院} (aller une fois à l'hôpital).
\end{itemize}
\emph{Règle d'or :} Jamais de numéral nu sans classificateur (*\zh{一病} $\times$, *\zh{两医生} $\times$).
\end{aretenir}

\courssub{Clés (Radicaux) et Ordre des traits : Les briques du corps médical}

\begin{propriete}[Radicaux clés du vocabulaire médical]
La clé (radical) indique le champ sémantique. La reconnaître permet de deviner le sens et de classer les caractères dans le dictionnaire.
\begin{center}
\begin{tabularx}{\linewidth}{|X|X|X|X|}
\hline
\rowcolor{popgreenL} \textbf{Clé (Caractère)} & \textbf{Nom / Sens} & \textbf{Exemples du thème} & \textbf{Autres exemples HSK} \\
\hline
\zh{疒} (nè) & \emph{bìngzìtóu} — maladie & \zh{病} (bìng), \zh{疾} (jí), \zh{疼} (téng), \zh{痛} (tòng), \zh{症} (zhèng: symptôme) & \zh{疲} (fatigue), \zh{瘦} (maigre) \\
\zh{艹} (cǎo) & herbe, plante, médicament & \zh{药} (yào), \zh{艾} (ài (artemisia/sida)) & \zh{花} (fleur), \zh{草} (herbe), \zh{菜} (légume) \\
\zh{氵} (shuǐ) & eau, liquide & \zh{滋} (zī: croître, sida) & \zh{河} (fleuve), \zh{海} (mer), \zh{汤} (soupe) \\
\zh{月} (yuè) & \emph{ròuzìpáng} — chair, organe, corps & \zh{肚} (dù: ventre), \zh{腿} (tuǐ: jambe), \zh{脏} (zàng: organe), \zh{肺} (fèi: poumon) & \zh{肝} (foie), \zh{心} (cœur — clé propre) \\
\zh{酉} (jiǔ) & vase, alcool, fermentation & (contexte prévention: \zh{酒} jiǔ — alcool, facteur de risque) & \zh{醉} (zuì: ivre), \zh{醋} (cù: vinaigre) \\
\zh{矢} (shǐ) & flèche & \zh{医} (yī: médecine — flèche dans la boîte/soin) & \zh{知} (zhī: savoir) \\
\hline
\end{tabularx}
\end{center}
\end{propriete}

\begin{exemplebox}[Ordre des traits — 6 caractères modèles]
Écrivez en respectant : \textbf{Haut $\rightarrow$ Bas, Gauche $\rightarrow$ Droite, Horizontal $\rightarrow$ Vertical, Centre $\rightarrow$ Côtés, Fermer en dernier}.
\begin{enumerate}
\item \zh{病} (10 traits) : Clé \zh{疒} (5 traits : 丶 一 丿 丶 ㇏) + \zh{丙} (5 traits : 一 丨 一 一 丨) à l'intérieur.
\item \zh{药} (9 traits) : Clé \zh{艹} (3 traits : 一 丨 一) + \zh{乐} (6 traits : 丨 丨 一 一 丨 一) — \emph{attention :} \zh{乐} s'écrit différemment de \zh{木}.
\item \zh{医} (7 traits) : \zh{匚} (2 traits : 丨 ㇆) + \zh{矢} (5 traits : 一 丨 一 丿 丶) à l'intérieur.
\item \zh{预} (12 traits) : Clé \zh{页} (6 traits, bas) + \zh{予} (4 traits, haut) — \emph{structure haut-bas}. Écrire \zh{予} puis \zh{页}.
\item \zh{防} (7 traits) : Clé \zh{阝} (2 traits, droite : ㇆ 丨) + \zh{方} (4 traits, gauche : 丶 一 丿 丶) — \emph{structure gauche-droite}. Écrire \zh{方} puis \zh{阝}.
\item \zh{艾} (14 traits) : Clé \zh{艹} (3 traits) + \zh{艾} (partie phonétique, 11 traits : 丨 丨 一 一 丨 一 丿 丶 丿 丶 乀). Complexe : à mémoriser par répétition.
\end{enumerate}
\end{exemplebox}

\begin{popfigure}
\begin{tikzpicture}[mindmap, grow cyclic, every node/.style=concept, concept color=popblue!60,
    level 1/.append style={level distance=4.5cm, sibling angle=90},
    level 2/.append style={level distance=3cm, sibling angle=45}]
\node[concept color=popdark] {\textbf{Santé / 健康}}
    child[concept color=popgreen] { node {\textbf{Maladies / 疾病}} 
        child { node {\zh{感冒} Rhume} }
        child { node {\zh{发烧} Fièvre} }
        child { node {\zh{头疼} Mal tête} }
        child { node {\zh{艾滋病} SIDA} }
        child { node {\zh{传染病} Infectieuse} }
    }
    child[concept color=poporange] { node {\textbf{Soins / 医疗}} 
        child { node {\zh{医院} Hôpital} }
        child { node {\zh{医生} Médecin} }
        child { node {\zh{药} Médicament} }
        child { node {\zh{看病} Consulter} }
        child { node {\zh{治疗} Traiter} }
    }
    child[concept color=poppurple] { node {\textbf{Prévention / 预防}} 
        child { node {\zh{了解} Comprendre} }
        child { node {\zh{预防} Prévenir} }
        child { node {\zh{健康知识} Éduc. santé} }
        child { node {\zh{保护} Protéger} }
        child { node {\zh{不做危险的事} Éviter risques} }
    }
    child[concept color=poppink] { node {\textbf{Acteurs / Lieux}} 
        child { node {\zh{年轻人} Jeunes} }
        child { node {\zh{家人} Famille} }
        child { node {\zh{自己} Soi-même} }
        child { node {\zh{喀麦隆} Cameroun} }
        child { node {\zh{全世界} Monde} }
    };
\end{tikzpicture}
\legende{Carte mentale du vocabulaire thématique (Leçon 13)}
\end{popfigure}

\courssub{Compréhension de texte : Questions graduelles}

\begin{exoresolu}[Modèle de compréhension — Type Examen]
\textbf{Texte de référence} : Celui de la section 1 (\zh{认识疾病，预防艾滋}).

\textbf{Questions} :
\begin{enumerate}
\item 根据课文，感冒、发烧、头疼属于什么病？(D'après le texte, rhume, fièvre, maux de tête sont quel genre de maladies ?)
\item 艾滋病和感冒有什么不同？(Quelle est la différence entre le sida et le rhume ?)
\item 为什么说「预防胜于治疗」？(Pourquoi dit-on « Prévenir vaut mieux que guérir » ?)
\item 年轻人应该怎么做，才能保护自己和家人？(Que doivent faire les jeunes pour se protéger eux et leur famille ?)
\end{enumerate}
\tcblower
\textbf{Solution détaillée} :
\begin{enumerate}
\item \textbf{Repérage} : Phrase 2 : \zh{感冒、发烧、头疼都是常见的病}. \\
\textbf{Réponse} : \zh{它们是常见的病。} \emph{Tāmen shì chángjiàn de bìng.} « Ce sont des maladies courantes. »
\item \textbf{Repérage} : Phrase 3-4 : \zh{但是，艾滋病不一样。艾滋病是一种很危险的传染病…很难治好}. \\
\textbf{Réponse} : \zh{艾滋病是一种很危险的传染病，很难治好；感冒是常见的病，吃药就能好。} \emph{Àizībìng shì yì zhǒng hěn wēixiǎn de chuánrǎnbìng, hěn nán zhì hǎo; gǎnmào shì chángjiàn de bìng, chī yào jiù néng hǎo.} « Le sida est une maladie infectieuse très dangereuse, difficile à guérir ; le rhume est une maladie courante, les médicaments font guérir. »
\item \textbf{Inférence} : Phrase 5-6 : \zh{身体会越来越弱，很难治好。所以…学会预防它。最后一句：预防胜于治疗}. \\
\textbf{Réponse} : \zh{因为得了艾滋病身体会越来越弱，很难治好，所以预防比治疗更重要。} \emph{Yīnwèi děi le àizībìng shēntǐ huì yuè lái yuè ruò, hěn nán zhì hǎo, suǒyǐ yùfáng bǐ zhìliáo gèng zhòngyào.} « Parce que le sida affaiblit le corps et est difficile à soigner, la prévention est plus importante que le traitement. »
\item \textbf{Synthèse} : Dernière phrase : \zh{年轻人要多学习健康知识，不做危险的事，保护自己和家人}. \\
\textbf{Réponse} : \zh{年轻人要多学习健康知识，不做危险的事，这样才能保护自己和家人。} \emph{Niánqīngrén yào duō xuéxí jiànkāng zhīshi, bù zuò wēixiǎn de shì, zhèyàng cáinéng bǎohù zìjǐ hé jiārén.} « Les jeunes doivent étudier les connaissances sanitaires, ne pas faire de choses dangereuses, ainsi ils pourront se protéger et protéger leur famille. »
\end{enumerate}
\end{exoresolu}

\begin{exercice}[Entraînement — Compréhension écrite]
Répondez en chinois (caractères + pinyin) en vous basant \textbf{uniquement} sur le texte support.
\begin{enumerate}
\item 疾病是谁的敌人？(L'ennemi de qui ?)
\item 去医院看医生、吃药，常见的病能好吗？(Est-ce que ça guérit ?)
\item 喀麦隆有艾滋病吗？(Le Cameroun a-t-il le sida ?)
\item 得了艾滋病，身体会怎么样？(Le corps devient comment ?)
\item 最后一句话「健康最重要」用中文怎么写？(Écrivez la dernière phrase.)
\end{enumerate}
\end{exercice}

\begin{corrige}[Exercice — Compréhension écrite]
\begin{enumerate}
\item \zh{疾病是人类的敌人。} \emph{Jíbìng shì rénlèi de dírén.}
\item \zh{能，就能好。} \emph{Néng, jiù néng hǎo.} (Ou : \zh{常见的病，去医院看医生、吃药就能好。})
\item \zh{有，喀麦隆也有。} \emph{Yǒu, Kāmàilóng yě yǒu.}
\item \zh{身体会越来越弱。} \emph{Shēntǐ huì yuè lái yuè ruò.}
\item \zh{预防胜于治疗，健康最重要。} \emph{Yùfáng shèng yú zhìliáo, jiànkāng zuì zhòngyào.}
\end{enumerate}
\end{corrige}

\courssub{Collocations, Structures utiles et Production guidée}

\begin{propriete}[Collocations fixes (搭配) — À mémoriser par cœur]
En chinois, les mots ne s'associent pas librement. Ces couples Verbe + Objet ou Adj + Nom sont obligatoires.
\begin{center}
\begin{tabularx}{\linewidth}{|X|X|X|}
\hline
\rowcolor{poporangeL} \textbf{Structure / Collocation} & \textbf{Pinyin} & \textbf{Traduction / Contexte} \\
\hline
\zh{生病} & shēng bìng & Tomber malade (V + O) — \zh{我生病了} \\
\zh{看病 / 看医生} & kàn bìng / kàn yīshēng & Consulter un médecin — \zh{去医院看病} \\
\zh{吃药 / 服药} & chī yào / fú yào & Prendre des médicaments (oral / formel) \\
\zh{开药} & kāi yào & Prescrire des médicaments (médecin sujet) — \zh{医生给我开药} \\
\zh{打针} & dǎ zhēn & Faire une injection/piqûre — \zh{护士给我打针} \\
\zh{预防疾病 / 预防艾滋病} & yùfáng jíbìng / yùfáng àizībìng & Prévenir les maladies / le sida \\
\zh{传染疾病 / 传染给别人} & chuánrǎn jíbìng / chuánrǎn gěi biérén & Transmettre une maladie / contaminer qqn \\
\zh{保护自己 / 保护家人} & bǎohù zìjǐ / bǎohù jiārén & Se protéger / Protéger sa famille \\
\zh{了解情况 / 了解知识} & liǎojiě qíngkuàng / liǎojiě zhīshi & Se renseigner / Acquérir des connaissances \\
\zh{保持健康 / 身体健康} & bǎochí jiànkāng / shēntǐ jiànkāng & Rester en bonne santé / Avoir une bonne santé \\
\zh{越来越 + Adj.} & yuè lái yuè + Adj. & De plus en plus... — \zh{越来越弱}, \zh{越来越强} \\
\zh{不但…而且…} & bùdàn… érqiě… & Non seulement… mais encore… — \zh{不但危险，而且难治好} \\
\zh{因为…所以…} & yīnwèi… suǒyǐ… & Parce que… donc… (Cause $\rightarrow$ Conséquence) \\
\hline
\end{tabularx}
\end{center}
\end{propriete}

\begin{methode}[Traduction Thème (Français $\rightarrow$ Chinois) : 4 étapes]
\begin{enumerate}
\item \textbf{Découper} : Identifiez Sujet / Verbe / Compléments (Temps, Lieu, Manière, Objet).
\item \textbf{Ordonner} : Appliquez l'ordre chinois canonique : \cle{Sujet + Temps + Lieu + Manière + Verbe + Objet}. (Le temps/lieu \textbf{avant} le verbe).
\item \textbf{Sélectionner} : Choisissez le bon classificateur (\zh{种, 个, 些, 次}), le bon mot-outil (\zh{了, 的, 地, 得}), la bonne collocation (\zh{看病} pas *\zh{看医院}).
\item \textbf{Vérifier} : Tons (pinyin), ponctuation chinoise (，。), accord sujet-verbe (pas de conjugaison mais aspect \zh{了}/\zh{过}/\zh{着}).
\end{enumerate}
\end{methode}

\begin{exoresolu}[Thème guidé : Du français au chinois]
\textbf{Énoncé} : Traduis : « Hier, mon petit frère a eu de la fièvre. Maman l'a emmené à l'hôpital. Le médecin lui a prescrit des médicaments et a dit : "Il faut boire beaucoup d'eau et se reposer." Mon frère a écouté le médecin. Aujourd'hui, il va mieux. »

\tcblower
\textbf{Solution détaillée} :
\begin{enumerate}
\item \textbf{Phrase 1} : « Hier, mon petit frère a eu de la fièvre. »
Sujet: \zh{我弟弟} | Temps: \zh{昨天} | Verbe: \zh{发烧} | Aspect: \zh{了} (accompli).
$\rightarrow$ \zh{我弟弟昨天发烧了。} \emph{Wǒ dìdi zuótiān fāshāo le.}
\item \textbf{Phrase 2} : « Maman l'a emmené à l'hôpital. »
Sujet: \zh{妈妈} | Objet: \zh{他} | Verbe: \zh{带} (emmener) | Lieu: \zh{去医院} | Aspect: \zh{了}.
$\rightarrow$ \zh{妈妈带他去了医院。} \emph{Māma dài tā qù le yīyuàn.} (Structure \zh{带 + qqn + 去 + Lieu}).
\item \textbf{Phrase 3} : « Le médecin lui a prescrit des médicaments et a dit : "Il faut boire beaucoup d'eau et se reposer." »
Partie A: Sujet: \zh{医生} | Verbe: \zh{给} (donner/à) + Objet: \zh{他} + Objet2: \zh{开了一些药} (prescrit des médocs).
$\rightarrow$ \zh{医生给他开了一些药。} \emph{Yīshēng gěi tā kāi le yìxiē yào.}
Partie B: Discours direct. « Il faut » = \zh{要} / \zh{应该}. « Boire de l'eau » = \zh{喝水} (collocation). « Beaucoup » = \zh{多}. « Se reposer » = \zh{休息}.
$\rightarrow$ \zh{医生说：「你要多喝水，多休息。」} \emph{Yīshēng shuō: « Nǐ yào duō hē shuǐ, duō xiūxi. »} (Impératif conseil).
\item \textbf{Phrase 4} : « Mon frère a écouté le médecin. »
Sujet: \zh{我弟弟} | Verbe: \zh{听} (écouter/obéir) + Objet: \zh{医生的} (celui du médecin) / \zh{医生的话} (paroles du médecin). Aspect: \zh{了}.
$\rightarrow$ \zh{我弟弟听了医生的话。} \emph{Wǒ dìdi tīng le yīshēng de huà.}
\item \textbf{Phrase 5} : « Aujourd'hui, il va mieux. »
Sujet: \zh{他} | Temps: \zh{今天} | Verbe: \zh{好} (aller bien) + Complément degré: \zh{多了} (beaucoup mieux) / \zh{好多了}.
$\rightarrow$ \zh{今天他好多了。} \emph{Jīntiān tā hǎo duō le.}
\end{enumerate}
\textbf{Texte final} :
\zh{我弟弟昨天发烧了。妈妈带他去了医院。医生给他开了一些药，说：「你要多喝水，多休息。」我弟弟听了医生的话。今天他好多了。}
\end{exoresolu}

\begin{exoresolu}[Production écrite guidée : Conseil de prévention]
\textbf{Consigne} : Rédigez un court message (5-6 phrases) pour une affiche de prévention au lycée : « Comment prévenir le sida ? » Utilisez : \zh{应该}, \zh{要}, \zh{因为…所以…}, \zh{预防}, \zh{了解}, \zh{保护}, \zh{健康}.

\tcblower
\textbf{Modèle de rédaction} :
\zh{同学们，艾滋病很危险，但是我们可以预防它。我们应该了解艾滋病的知识，因为了解才能预防。年轻人要保护自己，不做危险的事。也要告诉朋友和家人，一起预防艾滋病。预防胜于治疗，健康最重要，让我们一起行动吧！} \\[4pt]
\emph{Tóngxuémen, àizībìng hěn wēixiǎn, dànshì wǒmen kěyǐ yùfáng tā. Wǒmen yīnggāi liǎojiě àizībìng de zhīshi, yīnwèi liǎojiě cáinéng yùfáng. Niánqīngrén yào bǎohù zìjǐ, bù zuò wēixiǎn de shì. Yě yào gàosu péngyou hé jiārén, yìqǐ yùfáng àizībìng. Yùfáng shèng yú zhìliáo, jiànkāng zuì zhòngyào, ràng wǒmen yìqǐ xíngdòng ba!} \\[4pt]
« Camarades, le sida est dangereux mais on peut le prévenir. Nous devons connaître le sida, car comprendre permet de prévenir. Les jeunes doivent se protéger, ne pas faire de choses dangereuses. Il faut aussi en parler aux amis et à la famille, prévenir ensemble. Prévenir vaut mieux que guérir, la santé est primordiale, agissons ensemble ! »

\textbf{Analyse grammaticale} :
\begin{itemize}
\item \zh{同学们} (Vocatif) + phrases impératives/conseils (\zh{应该}, \zh{要}).
\item \zh{因为…所以…} (ici \zh{因为…才能…} : seulement si... alors...).
\item \zh{一起} (ensemble) + Verbe.
\item \zh{让我们…吧} (Exhortation : « Faisons... »).
\end{itemize}
\end{exoresolu}

\begin{exercice}[Entraînement — Thème et Production]
\begin{enumerate}
\item \textbf{Traduisez en chinois} :
    a) Le sida est une maladie infectieuse dangereuse.
    b) Nous devons aller à l'hôpital voir le médecin.
    c) Le médecin m'a donné des médicaments hier.
    d) Les jeunes doivent protéger leur santé.
\item \textbf{Complétez les dialogues} (choisissez : \zh{预防} / \zh{治疗} / \zh{了解} / \zh{保护} / \zh{健康} / \zh{危险}) :
    A: 我们为什么要\_\_\_\_艾滋病？ \\
    B: 因为\_\_\_\_胜于\_\_\_\_，\_\_\_\_最重要。
    A: 这种事很\_\_\_\_，我们要\_\_\_\_自己。
\item \textbf{Rédigez} (4 phrases min.) : « Mon ami est malade. Je lui donne un conseil. » (Utilisez : \zh{生病}, \zh{去医院}, \zh{吃药}, \zh{多喝水}, \zh{休息}, \zh{应该}).
\end{enumerate}
\end{exercice}

\begin{corrige}[Exercice — Thème et Production]
\begin{enumerate}
\item a) \zh{艾滋病是一种危险的传染病。} \emph{Àizībìng shì yì zhǒng wēixiǎn de chuánrǎnbìng.}
    b) \zh{我们应该去医院看医生。} \emph{Wǒmen yīnggāi qù yīyuàn kàn yīshēng.}
    c) \zh{医生昨天给了我一些药。} \emph{Yīshēng zuótiān gěi le wǒ yìxiē yào.} (Ou : \zh{医生昨天给我开了一些药。})
    d) \zh{年轻人要保护健康。} \emph{Niánqīngrén yào bǎohù jiànkāng.} (Ou : \zh{保护自己的健康。})
\item A: \zh{了解} — B: \zh{预防} / \zh{治疗} / \zh{预防} / \zh{健康} / \zh{危险} / \zh{保护}.
    \zh{我们为什么要了解艾滋病？因为预防胜于治疗，健康最重要。这种事很危险，我们要保护自己。}
\item \textbf{Modèle} : \zh{我朋友生病了。他应该去医院看医生，吃药。医生说要多喝水，多休息。他听医生的话，很快就会好起来。} \emph{Wǒ péngyou shēngbìng le. Tā yīnggāi qù yīyuàn kàn yīshēng, chī yào. Yīshēng shuō yào duō hē shuǐ, duō xiūxi. Tā tīng yīshēng de huà, hěn kuài jiù huì hǎo qǐlái.}
\end{enumerate}
\end{corrige}

\courssub{Approfondissement socioculturel : Cameroun – Chine, regards croisés}

\begin{savaistu}[Santé publique et SIDA : Parallèles Cameroun-Chine]
\begin{itemize}
\item \textbf{Épidémiologie} : Le Cameroun a un taux de prévalence du VIH parmi les plus élevés d'Afrique centrale (\textasciitilde 2,7\% adultes, estimations récentes). La Chine a une prévalence faible (\textless 0,1\%) mais un nombre absolu de cas en hausse chez les jeunes et les personnes âgées. Les deux pays ciblent la \textbf{prévention chez les jeunes} (éducation sexuelle, préservatifs, dépistage).
\item \textbf{Accès aux ARV (Antirétroviraux)} : Au Cameroun, la gratuité des ARV est effective depuis 2007 (Fonds mondial, État). En Chine, le programme « Quatre gratuits et un soin » (\zh{四免一关怀}) : dépistage gratuit, ARV gratuits pour les pauvres, prévention mère-enfant gratuite, scolarité gratuite pour orphelins du sida, aide aux familles touchées.
\item \textbf{Médecine traditionnelle} : Au Cameroun, le \emph{médecin traditionnel} (herboriste, \zh{中医} \emph{zhōngyī} en chinois) est consulté en premier pour \zh{发烧}, \zh{疟疾} (paludisme). En Chine, la \zh{中医} (Médecine Traditionnelle Chinoise — MTC) est intégrée aux hôpitaux modernes : \zh{中西医结合} (intégration MTC / médecine occidentale). Pour le VIH, la MTC est utilisée en \emph{adjuvant} (renforcer l'immunité, \zh{提高免疫力}) mais \textbf{ne remplace pas} les ARV.
\item \textbf{Stigmatisation} : Dans les deux pays, la peur du regard social (\zh{歧视} qíshì) freine le dépistage. Campagnes communes : « Connaître son statut », « Zéro discrimination » (\zh{零歧视}).
\item \textbf{Vocabulaire culturel} :
    \begin{itemize}
    \item \zh{体检} (tǐjiǎn) : Check-up médical (obligatoire à l'embauche en Chine, courant au Cameroun).
    \item \zh{献血} (xiànxiě) / \zh{捐血} (juānxiě) : Don de sang (dépistage VIH systématique).
    \item \zh{安全套} (ānquántào) / \zh{避孕套} (bìyùntào) : Préservatif (terme courant / terme médical).
    \item \zh{母婴阻断} (mǔyīng zǔduàn) : Prévention de la transmission mère-enfant (PTME).
    \end{itemize}
\end{itemize}
\emph{Activité orale suggérée} : Débat en binôme : « Au Cameroun, on va d'abord voir le tradipraticien. En Chine, on va à l'hôpital public. Quels avantages/inconvénients ? » Utilisez \zh{虽然…但是…} (Bien que... cependant...).
\end{savaistu}

\begin{experience}[Activité de classe : Jeu de rôle « À l'hôpital »]
\textbf{Contexte} : Un élève joue le \zh{病人} (patient : fièvre, mal à la tête, fatigue), l'autre le \zh{医生} (médecin : questions, diagnostic, ordonnance, conseils).
\textbf{Structures imposées} :
\begin{itemize}
\item Médecin : \zh{你哪里不舒服？} / \zh{你怎么了？} / \zh{发烧了吗？} / \zh{我给你开点药。} / \zh{要多喝水，多休息。}
\item Patient : \zh{我头疼，发烧。} / \zh{我吃药了，没好。} / \zh{要吃几天药？} / \zh{谢谢医生。}
\end{itemize}
\textbf{Évaluation} : Prononciation des tons (bìng/bing, yào/yào), ordre des mots (Temps avant Verbe), classificateurs (\zh{一些药}, \zh{三天}).
\end{experience}

\courssub{Erreurs fréquentes et conseils d'examen}

\begin{attention}[Les 5 erreurs qui coûtent des points au Bac — Spécial Leçon 13]
\begin{enumerate}
\item \textbf{Ordre des mots calqué du français (Temps/Lieu)} :
\textbf{Faux} : *\zh{他去医院昨天} / *\zh{我在医院看病昨天}. \\
\textbf{Juste} : \zh{他昨天去医院。} / \zh{我昨天在医院看病。} \\
\cle{Règle} : Sujet + \textbf{Temps} + Lieu + Verbe. Jamais de complément de temps à la fin de phrase.
\item \textbf{Mauvais emploi de \zh{了} (Aspect accompli vs État)} :
\textbf{Faux} : *\zh{我生病} (pour « je suis tombé malade » passé). *\zh{我每天去了医院} (\zh{了} incompatible avec habitude \zh{每天}). \\
\textbf{Juste} : \zh{我生病了} (changement d'état/accompli). \zh{我每天去医院} (habitude, pas de \zh{了}). \\
\cle{Règle} : \zh{了} après le verbe = action terminée/état nouveau. Pas de \zh{了} avec \zh{每天}, \zh{经常}, \zh{从不}.
\item \textbf{Classificateur absent ou incorrect} :
\textbf{Faux} : *\zh{一病}, *\zh{两医生}, *\zh{三药}. \\
\textbf{Juste} : \zh{一种病} (\zh{种}), \zh{两个医生} (\zh{个}), \zh{三次打针} (\zh{次}), \zh{一些药} (\zh{些}). \\
\cle{Astuce} : \zh{病} prend \zh{种} (type) ou \zh{场} (épisode) ; \zh{医生} prend \zh{个} / \zh{位} (politesse) ; \zh{药} prend \zh{些} / \zh{盒} / \zh{瓶} / \zh{片}.
\item \textbf{Confusion homophones / Tons (Le piège \emph{yào})} :
\zh{药} (yào, 4\ieme{} ton, médicament, clé \zh{艹}) $\neq$ \zh{要} (yào, 4\ieme{} ton, vouloir/devoir, pas de clé).
\zh{病} (bìng, 4\ieme{} ton, maladie) $\neq$ \zh{冰} (bīng, 1\ier{} ton, glace).
\zh{传染} (chuánrǎn, contaminer) $\neq$ \zh{传然} (inexistant).
\textbf{Oral} : Un ton faux = mot incompris. Entraînez-vous : \zh{我要吃药} (Wǒ \textbf{yào} chī \textbf{yào}).
\item \textbf{Caractères « jumeaux » mal écrits (Clé \zh{疒})} :
\zh{病} (bìng, maladie) : clé \zh{疒} + \zh{丙}.
\zh{疼} (téng, avoir mal — verbe) : clé \zh{疒} + \zh{冬} (hiver : le mal « gèle » la tête).
\zh{痛} (tòng, douleur — nom / formel) : clé \zh{疒} + \zh{甬} (chemin : la douleur « passe »).
\zh{症} (zhèng, symptôme) : clé \zh{疒} + \zh{正} (correct/normal : signe de la maladie).
\textbf{Écrit} : Perte de points si \zh{疼} écrit au lieu de \zh{病} dans « \zh{一种病} ».
\end{enumerate}
\textbf{Check-list finale avant de rendre la copie} :
\begin{itemize}
\item[$\square$] Ordre : Sujet + Temps + Lieu + Verbe + Objet ?
\item[$\square$] Classificateur après \zh{一/两/这/那/每} ?
\item[$\square$] \zh{了} pour actions passées accomplies (pas avec habitude) ?
\item[$\square$] Tons notés sur les voyelles (ā á ǎ à) pour l'oral ?
\item[$\square$] Ponctuation chinoise (，。？) et pas d'espaces entre caractères ?
\end{itemize}
\end{attention}

\newpage

\coursec{Exercices}

\courssub{Je vérifie mes connaissances}

\begin{exercice}[QCM : le bon sens du mot]
Pour chaque question, entoure la bonne réponse (a, b ou c).
\begin{enumerate}
\item \zh{艾滋病} (àizībìng) désigne : \quad a) le paludisme \quad b) le SIDA \quad c) la tuberculose
\item \zh{病毒} (bìngdú) signifie : \quad a) un médicament \quad b) un virus \quad c) un hôpital
\item Pour dire « se faire dépister », on utilise : \quad a) \zh{检查} (jiǎnchá) \quad b) \zh{吃药} (chī yào) \quad c) \zh{发烧} (fāshāo)
\item \zh{预防} (yùfáng) veut dire : \quad a) guérir \quad b) prévenir \quad c) transmettre
\end{enumerate}
\end{exercice}

\begin{exercice}[Vrai ou faux ? Justifie]
Réponds par V (vrai) ou F (faux) et justifie ta réponse en français à l'aide d'une phrase du cours.
\begin{enumerate}
\item \zh{艾滋病可以通过握手传染。}(Àizībìng kěyǐ tōngguò wòshǒu chuánrǎn.) « Le SIDA peut se transmettre par une poignée de main. »
\item \zh{发烧是一种症状。}(Fāshāo shì yì zhǒng zhèngzhuàng.) « La fièvre est un symptôme. »
\item \zh{只要及时治疗，很多病都可以治好。}(Zhǐyào jíshí zhìliáo, hěn duō bìng dōu kěyǐ zhìhǎo.) « Pourvu qu'on se soigne à temps, beaucoup de maladies peuvent être guéries. »
\item \zh{预防疾病不需要知识。}(Yùfáng jíbìng bù xūyào zhīshi.) « Pour prévenir les maladies, on n'a pas besoin de connaissances. »
\end{enumerate}
\end{exercice}

\begin{exercice}[Vocabulaire : complète le tableau]
Recopie et complète le tableau avec les mots du glossaire de la leçon.
\begin{center}
\begin{tabularx}{\linewidth}{|X|X|X|X|}\hline
\rowcolor{popblueL} Caractères & Pinyin & Nature & Traduction \\ \hline
\zh{疾病} & \ldots\ldots & nom & \ldots\ldots \\ \hline
\ldots\ldots & chuánrǎn & \ldots\ldots & transmettre, contaminer \\ \hline
\zh{症状} & \ldots\ldots & \ldots\ldots & symptôme \\ \hline
\ldots\ldots & jiǎnchá & verbe & \ldots\ldots \\ \hline
\zh{治疗} & zhìliáo & \ldots\ldots & \ldots\ldots \\ \hline
\end{tabularx}
\end{center}
\end{exercice}

\begin{exercice}[Associe le pinyin aux caractères]
Relie chaque mot en caractères à son pinyin.
\begin{center}
\begin{tabularx}{\linewidth}{|X|X|}\hline
\rowcolor{popblueL} Caractères & Pinyin à associer \\ \hline
\zh{发烧} & yùfáng \\ \hline
\zh{咳嗽} & fāshāo \\ \hline
\zh{预防} & àizībìng \\ \hline
\zh{艾滋病} & késou \\ \hline
\zh{医院} & yīyuàn \\ \hline
\end{tabularx}
\end{center}
\end{exercice}

\courssub{Je m'entraîne}

\begin{exercice}[Traduis en chinois]
Traduis les phrases suivantes en chinois (caractères), en utilisant le vocabulaire de la leçon. Écris ensuite le pinyin sous chaque phrase.
\begin{enumerate}
\item Mon frère a de la fièvre, il doit aller à l'hôpital.
\item Il faut prévenir les maladies et se faire dépister régulièrement.
\item Le SIDA ne se transmet pas en parlant avec quelqu'un.
\item Ce médicament est très efficace contre le paludisme.
\end{enumerate}
\end{exercice}

\begin{exercice}[Transforme les phrases]
\begin{enumerate}
\item Mets la phrase à la voix passive avec \zh{被} : \zh{病毒传染了很多人。}(Bìngdú chuánrǎn le hěn duō rén.)
\item Transforme en question rhétorique avec \zh{难道……不……} : \zh{预防疾病很重要。}(Yùfáng jíbìng hěn zhòngyào.)
\item Remplace le verbe par un synonyme plus formel : \zh{这种病会传给别人。}(Zhè zhǒng bìng huì chuán gěi biérén.) (indice : \zh{传播} ou \zh{传染})
\item Complète avec \zh{只要……就……} : \zh{\ldots\ldots 经常锻炼，\ldots\ldots 不容易生病。}
\end{enumerate}
\end{exercice}

\begin{exercice}[Texte à trous]
Complète le texte suivant avec les mots de la liste : \zh{预防}、\zh{检查}、\zh{症状}、\zh{传染}、\zh{治疗}。

\begin{textebox}[La santé de Wang Ming]
王明最近身体不舒服，他有发烧和咳嗽的\ldots\ldots{}。\\[2pt]
Wáng Míng zuìjìn shēntǐ bù shūfu, tā yǒu fāshāo hé késou de \ldots\ldots\\[2pt]
« Wang Ming ne se sent pas bien ces derniers temps, il a des \ldots\ldots de fièvre et de toux. »\\[4pt]
妈妈带他去医院做\ldots\ldots{}，医生说他的病不会\ldots\ldots 给别人。\\[2pt]
Māma dài tā qù yīyuàn zuò \ldots\ldots, yīshēng shuō tā de bìng bú huì \ldots\ldots gěi biérén.\\[2pt]
« Sa mère l'emmène à l'hôpital pour un \ldots\ldots, le médecin dit que sa maladie ne \ldots\ldots pas aux autres. »\\[4pt]
只要按时吃药、好好\ldots\ldots{}，他很快就会好。平时我们要注意\ldots\ldots 疾病。\\[2pt]
Zhǐyào ànshí chī yào, hǎohao \ldots\ldots, tā hěn kuài jiù huì hǎo. Píngshí wǒmen yào zhùyì \ldots\ldots jíbìng.\\[2pt]
« Pourvu qu'il prenne ses médicaments à l'heure et se \ldots\ldots bien, il guérira vite. Dans la vie quotidienne, nous devons faire attention à \ldots\ldots les maladies. »
\end{textebox}
\end{exercice}

\begin{exercice}[Remets en ordre et apparie]
\textbf{A.} Remets les mots dans le bon ordre pour faire une phrase correcte, puis traduis-la.
\begin{enumerate}
\item \zh{疾病 / 知识 / 预防 / 需要 / 的}
\item \zh{只要 / 治疗 / 及时 / 就 / 能 / 治好 / 病 / 这种}
\end{enumerate}
\textbf{B.} Associe chaque symptôme à la maladie correspondante.
\begin{center}
\begin{tabularx}{\linewidth}{|X|X|}\hline
\rowcolor{popblueL} Symptôme & Maladie \\ \hline
\zh{发高烧、发冷} (fā gāoshāo, fālěng) & \zh{艾滋病} (àizībìng) \\ \hline
\zh{长期咳嗽、消瘦} (chángqī késou, xiāoshòu) & \zh{疟疾} (nüèjí) \\ \hline
\zh{免疫力下降} (miǎnyìlì xiàjiàng) & \zh{结核病} (jiéhébìng) \\ \hline
\end{tabularx}
\end{center}
\end{exercice}

\courssub{Je me prépare au Bac}

\begin{exercice}[Type Bac — Compréhension de texte et questions de langue]
Lis le texte suivant, puis réponds aux questions.

\begin{textebox}[Une nouvelle arme contre le VIH : la PrEP]
尽管艾滋病还不能完全治好，但是现在有很多办法可以预防它。\\[2pt]
Jǐnguǎn àizībìng hái bù néng wánquán zhìhǎo, dànshì xiànzài yǒu hěn duō bànfǎ kěyǐ yùfáng tā.\\[2pt]
« Bien que le SIDA ne puisse pas encore être complètement guéri, il existe aujourd'hui de nombreux moyens de le prévenir. »\\[4pt]
PrEP 是一种新的预防药。只要每天吃一片，就可以大大降低感染病毒的危险。\\[2pt]
PrEP shì yì zhǒng xīn de yùfáng yào. Zhǐyào měi tiān chī yí piàn, jiù kěyǐ dàdà jiàngdī gǎnrǎn bìngdú de wēixiǎn.\\[2pt]
« La PrEP est un nouveau médicament de prévention. Pourvu qu'on en prenne un comprimé par jour, on peut réduire fortement le risque d'infection par le virus. »\\[4pt]
医生告诉我们：这种药很安全，但是吃药以前，必须先做检查。早检查、早发现、早预防，健康才有保证。\\[2pt]
Yīshēng gàosu wǒmen: zhè zhǒng yào hěn ānquán, dànshì chī yào yǐqián, bìxū xiān zuò jiǎnchá. Zǎo jiǎnchá, zǎo fāxiàn, zǎo yùfáng, jiànkāng cái yǒu bǎozhèng.\\[2pt]
« Le médecin nous dit : ce médicament est sûr, mais avant de le prendre, il faut d'abord faire un test. Dépistage précoce, découverte précoce, prévention précoce : c'est ainsi que la santé est garantie. »
\end{textebox}

\textbf{Questions de compréhension}
\begin{enumerate}
\item Vrai ou faux ? Justifie par une citation du texte en chinois : « La PrEP guérit complètement le SIDA. »
\item Recopie et complète le tableau : \begin{center}\begin{tabularx}{\linewidth}{|X|X|X|}\hline
\rowcolor{popblueL} Médicament & Fonction & Précaution \\ \hline
\ldots\ldots & \ldots\ldots & \ldots\ldots \\ \hline
\end{tabularx}\end{center}
\item Que faut-il faire avant de prendre ce médicament ? Réponds en français.
\end{enumerate}

\textbf{Questions de langue}
\begin{enumerate}
\item Traduis en français : \zh{只要每天吃一片，就可以大大降低感染病毒的危险。}
\item Traduis en français : \zh{尽管艾滋病还不能完全治好，但是现在有很多办法可以预防它。}
\item Relève dans le texte une phrase avec la structure \zh{只有……才……} ou transforme la dernière phrase pour l'utiliser.
\end{enumerate}
\end{exercice}

\begin{exercice}[Type Bac — Thème et version]
\textbf{A. Version.} Traduis en français le passage suivant :

\begin{textebox}[Au centre de santé de Yaoundé]
在雅温得的一个医疗中心，医生们每天给年轻人讲解预防艾滋病的方法。\\[2pt]
Zài Yǎwēndé de yí ge yīliáo zhōngxīn, yīshēngmen měi tiān gěi niánqīngrén jiǎngjiě yùfáng àizībìng de fāngfǎ.\\[2pt]
« Dans un centre de santé de Yaoundé, les médecins expliquent chaque jour aux jeunes les méthodes de prévention du SIDA. »\\[4pt]
他们说：不要害怕检查，检查是保护自己和家人的第一步。\\[2pt]
Tāmen shuō: bú yào hàipà jiǎnchá, jiǎnchá shì bǎohù zìjǐ hé jiārén de dì yī bù.\\[2pt]
« Ils disent : n'ayez pas peur du dépistage, le dépistage est le premier pas pour se protéger soi-même et sa famille. »
\end{textebox}

\textbf{B. Thème.} Traduis en chinois (caractères) les phrases suivantes :
\begin{enumerate}
\item Beaucoup de maladies peuvent être prévenues grâce à de bonnes habitudes.
\item Bien que ce médicament soit cher, il est très efficace.
\item Mon ami ne veut pas aller à l'hôpital parce qu'il a peur du résultat du test.
\end{enumerate}
\end{exercice}

\begin{exercice}[Type Bac — Expression écrite]
\textbf{Sujet :} Ton ami(e) chinois(e) refuse de se faire dépister par peur du résultat. Écris-lui un message WeChat en chinois pour le/la convaincre.

\textbf{Consignes :}
\begin{itemize}
\item Longueur : 80 à 100 caractères chinois.
\item Tu dois employer obligatoirement les trois structures suivantes : \zh{只有……才……}, \zh{不要……要……}, \zh{早……早……}.
\item Ton message doit contenir : un mot d'encouragement, une information sur le dépistage (\zh{检查}), et une conclusion rassurante.
\item Rédige en caractères simplifiés ; le pinyin n'est pas exigé dans la copie finale.
\end{itemize}

\textbf{Aide vocabulaire :} \zh{害怕} (hàipà, avoir peur) ; \zh{勇敢} (yǒnggǎn, courageux) ; \zh{结果} (jiéguǒ, résultat) ; \zh{放心} (fàngxīn, être rassuré) ; \zh{保护} (bǎohù, protéger).
\end{exercice}

\newpage

\coursec{Corrigés des exercices}

\begin{corrige}[Exercice 1 — QCM : le bon sens du mot]
\begin{enumerate}
\item \textbf{b) le SIDA.} \zh{艾滋病} (àizībìng) : \zh{艾滋} est la transcription phonétique de « AIDS » et \zh{病} signifie « maladie ». Le paludisme se dit \zh{疟疾} (nüèjí) et la tuberculose \zh{结核病} (jiéhébìng).
\item \textbf{b) un virus.} \zh{病毒} (bìngdú) : \zh{病} « maladie » + \zh{毒} « poison, toxine ». Un médicament se dit \zh{药} (yào) et un hôpital \zh{医院} (yīyuàn).
\item \textbf{a) \zh{检查}} (jiǎnchá) : « examiner, contrôler » ; \zh{做检查} = « faire un test / se faire dépister ». \zh{吃药} (chī yào) = « prendre un médicament » ; \zh{发烧} (fāshāo) = « avoir de la fièvre ».
\item \textbf{b) prévenir.} \zh{预防} (yùfáng) = « prévenir ». Guérir se dit \zh{治好} (zhìhǎo) et transmettre \zh{传染} (chuánrǎn).
\end{enumerate}
\end{corrige}

\begin{corrige}[Exercice 2 — Vrai ou faux ? Justifie]
\begin{enumerate}
\item \textbf{F (faux).} Le SIDA ne se transmet pas par les contacts quotidiens. Le cours précise : \zh{艾滋病不会通过握手、吃饭、说话传染。} « Le SIDA ne se transmet ni par une poignée de main, ni en mangeant, ni en parlant. » Il se transmet par le sang, les rapports non protégés et de la mère à l'enfant.
\item \textbf{V (vrai).} \zh{发烧} « la fièvre » est bien un \zh{症状} « symptôme », c'est-à-dire un signe visible d'une maladie, comme la toux (\zh{咳嗽}).
\item \textbf{V (vrai).} La structure \zh{只要……就……} exprime la condition suffisante : un traitement rapide (\zh{及时治疗}) permet de guérir (\zh{治好}) beaucoup de maladies, par exemple le paludisme.
\item \textbf{F (faux).} C'est le contraire : \zh{预防疾病需要知识。} « Prévenir les maladies demande des connaissances. » Connaître les modes de transmission est la base de la prévention.
\end{enumerate}
\end{corrige}

\begin{corrige}[Exercice 3 — Vocabulaire : complète le tableau]
\begin{center}
\begin{tabularx}{\linewidth}{|X|X|X|X|}\hline
\rowcolor{popblueL} Caractères & Pinyin & Nature & Traduction \\ \hline
\zh{疾病} & jíbìng & nom & maladie (terme général) \\ \hline
\zh{传染} & chuánrǎn & verbe & transmettre, contaminer \\ \hline
\zh{症状} & zhèngzhuàng & nom & symptôme \\ \hline
\zh{检查} & jiǎnchá & verbe & examiner, dépister \\ \hline
\zh{治疗} & zhìliáo & verbe & soigner, traiter \\ \hline
\end{tabularx}
\end{center}
\textbf{Points de vigilance :} attention aux tons de \zh{疾病} (jíbìng, 2\textsuperscript{e} + 4\textsuperscript{e} ton) et à la voyelle ü de \zh{疟} dans \zh{疟疾} (nüèjí). Ne pas confondre \zh{检查} (jiǎnchá, « examen médical ») et \zh{测验} (cèyàn, « test scolaire »).
\end{corrige}

\begin{corrige}[Exercice 4 — Associe le pinyin aux caractères]
\begin{center}
\begin{tabularx}{\linewidth}{|X|X|}\hline
\rowcolor{popblueL} Caractères & Pinyin correct \\ \hline
\zh{发烧} & fāshāo \\ \hline
\zh{咳嗽} & késou \\ \hline
\zh{预防} & yùfáng \\ \hline
\zh{艾滋病} & àizībìng \\ \hline
\zh{医院} & yīyuàn \\ \hline
\end{tabularx}
\end{center}
\textbf{Points de vigilance :} dans \zh{咳嗽} (késou), le second caractère se lit au ton neutre ; \zh{发烧} (fāshāo) : 1\textsuperscript{er} ton sur les deux syllabes ; \zh{艾滋病} (àizībìng) : 4\textsuperscript{e}–1\textsuperscript{er}–4\textsuperscript{e} ton.
\end{corrige}

\begin{corrige}[Exercice 5 — Traduis en chinois]
\begin{enumerate}
\item \zh{我弟弟发烧了，他应该去医院。}\\
\emph{Wǒ dìdi fāshāo le, tā yīnggāi qù yīyuàn.}\\
Note : \zh{了} marque le changement d'état (il a maintenant de la fièvre) ; \zh{应该} (yīnggāi) = « devoir » (conseil).
\item \zh{我们要预防疾病，定期做检查。}\\
\emph{Wǒmen yào yùfáng jíbìng, dìngqī zuò jiǎnchá.}\\
Note : \zh{定期} (dìngqī) = « régulièrement » ; « se faire dépister » se rend par \zh{做检查}.
\item \zh{艾滋病不会通过和别人说话传染。}\\
\emph{Àizībìng bú huì tōngguò hé biérén shuōhuà chuánrǎn.}\\
Note : \zh{通过} (tōngguò) + nom/verbe = « par le biais de » ; la négation d'un fait futur/général est \zh{不会}.
\item \zh{这种药对疟疾很有效。}\\
\emph{Zhè zhǒng yào duì nüèjí hěn yǒuxiào.}\\
Note : structure \zh{对……有效} (duì … yǒuxiào) = « être efficace contre » ; classificateur \zh{种} (zhǒng) pour les sortes/types.
\end{enumerate}
\end{corrige}

\begin{corrige}[Exercice 6 — Transforme les phrases]
\begin{enumerate}
\item \zh{很多人被病毒传染了。}\\
\emph{Hěn duō rén bèi bìngdú chuánrǎn le.} « Beaucoup de gens ont été contaminés par le virus. »\\
Structure passive : Patient + \zh{被} + Agent + Verbe. Le patient (\zh{很多人}) passe en tête de phrase.
\item \zh{难道预防疾病不重要吗？}\\
\emph{Nándào yùfáng jíbìng bù zhòngyào ma?} « La prévention des maladies n'est-elle pas importante ? »\\
Structure : \zh{难道} + phrase négative + \zh{吗} = question rhétorique qui attend une réponse évidente (« si, bien sûr ! »).
\item \zh{这种病会传播给别人。} ou \zh{这种病会传染给别人。}\\
\emph{Zhè zhǒng bìng huì chuánbō (chuánrǎn) gěi biérén.}\\
\zh{传播} (chuánbō) est plus formel et plus large (diffuser) ; \zh{传染} (chuánrǎn) insiste sur la contagion.
\item \zh{只要经常锻炼，就不容易生病。}\\
\emph{Zhǐyào jīngcháng duànliàn, jiù bù róngyì shēngbìng.} « Pourvu qu'on fasse régulièrement du sport, on ne tombe pas facilement malade. »\\
Attention : \zh{就} se place avant le verbe de la seconde proposition, jamais en tête.
\end{enumerate}
\end{corrige}

\begin{corrige}[Exercice 7 — Texte à trous]
Texte complété (dans l'ordre des blancs) : \zh{症状}、\zh{检查}、\zh{传染}、\zh{治疗}、\zh{预防}.

\begin{textebox}[La santé de Wang Ming — corrigé]
王明最近身体不舒服，他有发烧和咳嗽的症状。\\[2pt]
Wáng Míng zuìjìn shēntǐ bù shūfu, tā yǒu fāshāo hé késou de zhèngzhuàng.\\[2pt]
« Wang Ming ne se sent pas bien ces derniers temps, il a des symptômes de fièvre et de toux. »\\[4pt]
妈妈带他去医院做检查，医生说他的病不会传染给别人。\\[2pt]
Māma dài tā qù yīyuàn zuò jiǎnchá, yīshēng shuō tā de bìng bú huì chuánrǎn gěi biérén.\\[2pt]
« Sa mère l'emmène à l'hôpital pour un examen, le médecin dit que sa maladie ne se transmet pas aux autres. »\\[4pt]
只要按时吃药、好好治疗，他很快就会好。平时我们要注意预防疾病。\\[2pt]
Zhǐyào ànshí chī yào, hǎohao zhìliáo, tā hěn kuài jiù huì hǎo. Píngshí wǒmen yào zhùyì yùfáng jíbìng.\\[2pt]
« Pourvu qu'il prenne ses médicaments à l'heure et se soigne bien, il guérira vite. Dans la vie quotidienne, nous devons veiller à prévenir les maladies. »
\end{textebox}
\textbf{Justifications :} \zh{症状} est le nom attendu après \zh{有……的} ; \zh{做检查} est la collocation « passer un examen » ; \zh{传染给} + personne = « transmettre à » ; \zh{治疗} suit l'adverbe \zh{好好} (« bien ») ; \zh{预防疾病} = « prévenir les maladies ».
\end{corrige}

\begin{corrige}[Exercice 8 — Remets en ordre et apparie]
\textbf{A. Phrases remises en ordre}
\begin{enumerate}
\item \zh{预防疾病需要知识。}\\
\emph{Yùfáng jíbìng xūyào zhīshi.} « Prévenir les maladies demande des connaissances. »\\
Ordre : Sujet (\zh{预防疾病}, verbe + nom formant un groupe sujet) + verbe \zh{需要} + complément \zh{知识}.
\item \zh{只要及时治疗，这种病就能治好。}\\
\emph{Zhǐyào jíshí zhìliáo, zhè zhǒng bìng jiù néng zhìhǎo.} « Pourvu qu'on la soigne à temps, cette maladie peut être guérie. »\\
Attention : \zh{只要} ouvre la condition ; \zh{就} précède le verbe \zh{能治好} dans la conséquence.
\end{enumerate}
\textbf{B. Appariement symptôme — maladie}
\begin{center}
\begin{tabularx}{\linewidth}{|X|X|}\hline
\rowcolor{popblueL} Symptôme & Maladie \\ \hline
\zh{发高烧、发冷} (fā gāoshāo, fālěng) & \zh{疟疾} (nüèjí) — paludisme \\ \hline
\zh{长期咳嗽、消瘦} (chángqī késou, xiāoshòu) & \zh{结核病} (jiéhébìng) — tuberculose \\ \hline
\zh{免疫力下降} (miǎnyìlì xiàjiàng) & \zh{艾滋病} (àizībìng) — SIDA \\ \hline
\end{tabularx}
\end{center}
\textbf{Justification :} le paludisme provoque fièvre élevée et frissons ; la tuberculose se manifeste par une toux prolongée et un amaigrissement ; le SIDA attaque le système immunitaire, d'où la baisse de l'immunité.
\end{corrige}

\begin{corrige}[Exercice 9 — Type Bac : compréhension de texte et questions de langue]
\textbf{Barème indicatif (10 points) :} compréhension 4 pts (Q1 : 1,5 ; Q2 : 1,5 ; Q3 : 1) ; langue 6 pts (2 pts par question).

\textbf{Questions de compréhension}
\begin{enumerate}
\item \textbf{Faux.} Le texte dit : \zh{尽管艾滋病还不能完全治好} « Bien que le SIDA ne puisse pas encore être complètement guéri ». La PrEP est un médicament de \emph{prévention} (\zh{预防药}), pas de guérison.
\item Tableau complété :
\begin{center}
\begin{tabularx}{\linewidth}{|X|X|X|}\hline
\rowcolor{popblueL} Médicament & Fonction & Précaution \\ \hline
PrEP（\zh{预防药}） & \zh{预防艾滋病，降低感染病毒的危险} & \zh{吃药以前，必须先做检查} \\ \hline
\end{tabularx}
\end{center}
\item Avant de prendre ce médicament, il faut d'abord faire un test (un dépistage) : \zh{吃药以前，必须先做检查。}
\end{enumerate}

\textbf{Questions de langue}
\begin{enumerate}
\item \zh{只要每天吃一片，就可以大大降低感染病毒的危险。} = « Pourvu qu'on en prenne un comprimé par jour, on peut réduire fortement le risque d'infection par le virus. » (\zh{只要……就……} = condition suffisante ; \zh{一片} : classificateur \zh{片} pour les comprimés.)
\item \zh{尽管艾滋病还不能完全治好，但是现在有很多办法可以预防它。} = « Bien que le SIDA ne puisse pas encore être complètement guéri, il existe aujourd'hui de nombreux moyens de le prévenir. » (\zh{尽管……但是……} = concession.)
\item Transformation avec \zh{只有……才……} : \zh{只有早检查、早发现、早预防，健康才有保证。}\\
\emph{Zhǐyǒu zǎo jiǎnchá, zǎo fāxiàn, zǎo yùfáng, jiànkāng cái yǒu bǎozhèng.} « Ce n'est qu'en se dépistant tôt, en découvrant tôt et en prévenant tôt que la santé est garantie. »\\
\textbf{Point de vigilance :} \zh{只有……才……} exprime la condition \emph{nécessaire} (« seulement si »), alors que \zh{只要……就……} exprime la condition \emph{suffisante} (« pourvu que »). Ne pas les intervertir.
\end{enumerate}
\end{corrige}

\begin{corrige}[Exercice 10 — Type Bac : thème et version]
\textbf{Barème indicatif (10 points) :} version 4 pts ; thème 6 pts (2 pts par phrase).

\textbf{A. Version — traduction de référence}

« Dans un centre de santé de Yaoundé, les médecins expliquent chaque jour aux jeunes les méthodes de prévention du SIDA. Ils disent : n'ayez pas peur du dépistage ; le dépistage est le premier pas pour se protéger soi-même et protéger sa famille. »

\textbf{Points de vigilance :} \zh{雅温得} (Yǎwēndé) est la transcription de « Yaoundé » ; \zh{给年轻人讲解} = « expliquer aux jeunes » (\zh{给} introduit le bénéficiaire) ; \zh{第一步} = « le premier pas » (\zh{第} + nombre = ordinal).

\textbf{B. Thème — traductions de référence}
\begin{enumerate}
\item \zh{很多疾病都可以通过好习惯预防。}\\
\emph{Hěn duō jíbìng dōu kěyǐ tōngguò hǎo xíguàn yùfáng.}\\
Variante passive acceptée : \zh{很多疾病都可以被好习惯预防。} (moins naturelle ; préférer la tournure active avec \zh{通过}).
\item \zh{尽管这种药很贵，但是很有效。}\\
\emph{Jǐnguǎn zhè zhǒng yào hěn guì, dànshì hěn yǒuxiào.}\\
Concession : \zh{尽管……但是……}. En chinois, contrairement au français, on peut (et on préfère) garder les deux connecteurs.
\item \zh{我的朋友不想去医院，因为他害怕检查的结果。}\\
\emph{Wǒ de péngyou bù xiǎng qù yīyuàn, yīnwèi tā hàipà jiǎnchá de jiéguǒ.}\\
Note : \zh{不想} = « ne pas vouloir » ; la cause s'exprime avec \zh{因为} (yīnwèi) ; \zh{检查的结果} = « le résultat du test » (\zh{的} de possession/appartenance).
\end{enumerate}
\end{corrige}

\begin{corrige}[Exercice 11 — Type Bac : expression écrite]
\textbf{Barème indicatif (10 points) :} respect des consignes et longueur 2 pts ; emploi correct des trois structures imposées 3 pts (1 pt chacune) ; exactitude grammaticale et vocabulaire 3 pts ; cohérence et ton encourageant 2 pts.

\textbf{Proposition de rédaction modèle (environ 90 caractères) :}

\begin{textebox}[Message WeChat modèle]
小明，你好！听说你害怕做检查，别担心，要勇敢一点！\\[2pt]
Xiǎo Míng, nǐ hǎo! Tīngshuō nǐ hàipà zuò jiǎnchá, bié dānxīn, yào yǒnggǎn yìdiǎn!\\[2pt]
« Salut Xiao Ming ! J'ai entendu que tu as peur de faire le test, ne t'inquiète pas, sois courageux ! »\\[4pt]
只有早检查，才能早知道自己的身体情况。不要害怕结果，要相信自己。\\[2pt]
Zhǐyǒu zǎo jiǎnchá, cái néng zǎo zhīdào zìjǐ de shēntǐ qíngkuàng. Bú yào hàipà jiéguǒ, yào xiāngxìn zìjǐ.\\[2pt]
« Ce n'est qu'en te dépistant tôt que tu connaîtras tôt ton état de santé. N'aie pas peur du résultat, aie confiance en toi. »\\[4pt]
早检查，早发现，早治疗，健康才有保证。检查很简单，也很安全。放心吧，我陪你一起去！\\[2pt]
Zǎo jiǎnchá, zǎo fāxiàn, zǎo zhìliáo, jiànkāng cái yǒu bǎozhèng. Jiǎnchá hěn jiǎndān, yě hěn ānquán. Fàngxīn ba, wǒ péi nǐ yìqǐ qù!\\[2pt]
« Dépistage précoce, découverte précoce, traitement précoce : c'est ainsi que la santé est garantie. Le test est simple et sûr. Sois rassuré, je t'accompagne ! »
\end{textebox}

\textbf{Vérification des structures imposées :}
\begin{itemize}
\item \zh{只有……才……} : \zh{只有早检查，才能早知道……} — condition nécessaire, correcte.
\item \zh{不要……要……} : \zh{不要害怕结果，要相信自己。} — opposition interdiction/conseil, correcte.
\item \zh{早……早……} : \zh{早检查，早发现，早治疗} — accumulation des adverbes \zh{早}, correcte.
\end{itemize}

\textbf{Points de vigilance :}
\begin{itemize}
\item Ne pas écrire \zh{不要……不要……} : la structure impose un contraste négation/conseil positif.
\item \zh{才} se place avant le verbe, jamais en tête de proposition.
\item Compter les caractères : rester entre 80 et 100 ; la ponctuation ne compte pas.
\item Terminer par une formule rassurante (\zh{放心吧！}) : c'est la conclusion attendue par la consigne.
\end{itemize}
\end{corrige}

\newpage

\coursec{Fiche bilan}

\begin{popfigure}
\begin{tikzpicture}[mindmap, grow cyclic, every node/.style={concept, align=center, font=\small},
root concept/.append style={concept color=popblue, font=\bfseries},
level 1/.append style={level distance=3.2cm, sibling angle=51.4},
level 1 concept/.append style={font=\footnotesize}]
\node[root concept, align=center]{Leçon 13\\ 认识疾病，预防艾滋\\ Maladies et SIDA}
child[concept color=poporange]{ node[align=center]{Lexique maladies\\ 病 bìng\\ 感冒 gǎnmào\\ 发烧 fāshāo} }
child[concept color=popgreen]{ node[align=center]{Lexique SIDA\\ 艾滋病 àizībìng\\ 病毒 bìngdú\\ 传染 chuánrǎn} }
child[concept color=poppurple]{ node[align=center]{Prévention\\ 预防 yùfáng\\ 保护 bǎohù\\ 安全 ānquán} }
child[concept color=popgold]{ node[align=center]{Radicaux\\ 疒 (maladie)\\ 心 (cœur)\\ 疋/正 (traits)} }
child[concept color=poppink]{ node[align=center]{Structures\\ 会 + V\\ 应该 + V\\ 要 + V} }
child[concept color=popblue!60]{ node[align=center]{Compréhension\\ lire, repérer\\ répondre en\\ phrases complètes} }
child[concept color=popgreen!60]{ node[align=center]{Socio-culture\\ Cameroun -- Chine\\ santé publique\\ respect, pas de\\ discrimination} };
\end{tikzpicture}
\legende{Carte mentale de la leçon 13 : maladies en général et SIDA en particulier.}
\end{popfigure}

\begin{bilan}
\end{bilan}

\courssub{Lexique clé à connaître par cœur}
\begin{itemize}
\item \zh{生病} shēngbìng : tomber malade ; \zh{病人} bìngrén : malade (n.) ; \zh{医院} yīyuàn : hôpital ; \zh{医生} yīshēng : médecin.
\item \zh{感冒} gǎnmào : rhume ; \zh{发烧} fāshāo : avoir de la fièvre ; \zh{头疼} tóuténg : avoir mal à la tête ; \zh{咳嗽} késou : tousser.
\item \zh{艾滋病} àizībìng : SIDA ; \zh{病毒} bìngdú : virus ; \zh{传染} chuánrǎn : transmettre, contagieux ; \zh{预防} yùfáng : prévenir.
\item \zh{保护} bǎohù : protéger ; \zh{健康} jiànkāng : santé ; \zh{危险} wēixiǎn : dangereux ; \zh{安全} ānquán : sûr, sécurité.
\end{itemize}

\courssub{Règles de grammaire à maîtriser}
\begin{center}
\begin{tabularx}{\linewidth}{|X|X|X|}
\hline
\rowcolor{popblueL} Structure & Exemple (caractères + pinyin) & Traduction \\
\hline
Sujet + 会 + verbe (possibilité, capacité) & \zh{艾滋病会传染。} Àizībìng huì chuánrǎn. & Le SIDA peut se transmettre. \\
\hline
Sujet + 应该 + verbe (devoir, conseil) & \zh{我们应该预防疾病。} Wǒmen yīnggāi yùfáng jíbìng. & Nous devons prévenir les maladies. \\
\hline
Sujet + 要 + verbe (il faut, volonté) & \zh{生病了要去看医生。} Shēngbìng le yào qù kàn yīshēng. & Quand on est malade, il faut aller voir le médecin. \\
\hline
Verbe + 了 (action accomplie) & \zh{他感冒了。} Tā gǎnmào le. & Il a attrapé un rhume. \\
\hline
不 / 没有 + verbe (négation) & \zh{他没有发烧。} Tā méiyou fāshāo. & Il n'a pas de fièvre. \\
\hline
\end{tabularx}
\end{center}

\courssub{Écriture : clés et traits}
\begin{itemize}
\item Clé de la maladie \zh{疒} (nè) : présente dans \zh{病}, \zh{疼}, \zh{痛} ; elle se trace en 5 traits (point, horizontal, puis l'enveloppe).
\item Ordre général : haut $\rightarrow$ bas, gauche $\rightarrow$ droite, extérieur avant intérieur, horizontal avant vertical.
\item \zh{病} : clé 疒 + \zh{丙} (bǐng) à l'intérieur ; \zh{感} : \zh{咸} au-dessus de la clé du cœur \zh{心}.
\end{itemize}

\courssub{Méthodes express}
\begin{enumerate}
\item Compréhension : lire une première fois sans s'arrêter, repérer le thème, puis souligner les mots-clés (maladie, symptôme, conseil) avant de répondre.
\item Répondre en phrases complètes : reprendre les mots de la question (问什么，答什么).
\item Vocabulaire : apprendre par familles (maladie / symptôme / prévention) et par collocations : \zh{预防感冒}, \zh{保护自己}, \zh{去看医生}.
\item Tons : vérifier à voix haute les paires sensibles : \zh{bìng} (4\up{e} ton, maladie) et \zh{bìngrén}.
\end{enumerate}


\begin{aretenir}[Checklist avant le Bac]
\begin{itemize}
\item[$\square$] Je sais lire et écrire les 15 mots clés du thème (maladies, symptômes, SIDA, prévention).
\item[$\square$] Je connais le pinyin et les tons de \zh{艾滋病}, \zh{传染}, \zh{预防}, \zh{发烧} sans hésiter.
\item[$\square$] Je reconnais la clé 疒 et je sais tracer \zh{病} et \zh{疼} dans le bon ordre.
\item[$\square$] J'emploie correctement 会, 应该 et 要 devant un verbe.
\item[$\square$] Je place 了 juste après le verbe pour une action accomplie (\zh{他感冒了。}).
\item[$\square$] Je fais la négation avec 不 (présent/futur) et 没有 (passé) sans les confondre.
\item[$\square$] Je sais répondre à une question de compréhension par une phrase complète en chinois.
\item[$\square$] Je connais trois conseils de prévention en chinois : \zh{注意卫生}, \zh{保护自己}, \zh{去看医生}.
\item[$\square$] Je sais parler du SIDA avec respect, sans discrimination, en chinois comme en français.
\item[$\square$] Je peux comparer brièvement la santé publique au Cameroun et en Chine.
\end{itemize}
\end{aretenir}

\findecours

\end{document}
